%% bare_jrnl.tex
%% V1.4b
%% 2015/08/26
%% by Michael Shell
%% see http://www.michaelshell.org/
%% for current contact information.
%% 
%% This is a skeleton file demonstrating the use of IEEEtran.cls
%% (requires IEEEtran.cls version 1.8b or later) with an IEEE
%% journal paper.
%% 
%% Support sites:
%% http://www.michaelshell.org/tex/ieeetran/
%% http://www.ctan.org/pkg/ieeetran
%% and
%% http://www.ieee.org/

%%*************************************************************************
%% Legal Notice:
%% This code is offered as-is without any warranty either expressed or
%% implied; without even the implied warranty of MERCHANTABILITY or
%% FITNESS FOR A PARTICULAR PURPOSE! 
%% User assumes all risk.
%% In no event shall the IEEE or any contributor to this code be liable for
%% any damages or losses, including, but not limited to, incidental,
%% consequential, or any other damages, resulting from the use or misuse
%% of any information contained here.
%% 
%% All comments are the opinions of their respective authors and are not
%% necessarily endorsed by the IEEE.
%% 
%% This work is distributed under the LaTeX Project Public License (LPPL)
%% ( http://www.latex-project.org/ ) version 1.3, and may be freely used,
%% distributed and modified. A copy of the LPPL, version 1.3, is included
%% in the base LaTeX documentation of all distributions of LaTeX released
%% 2003/12/01 or later.
%% Retain all contribution notices and credits.
%% ** Modified files should be clearly indicated as such, including  **
%% ** renaming them and changing author support contact information. **
%%*************************************************************************


% *** Authors should verify (and, if needed, correct) their LaTeX system  ***
% *** with the testflow diagnostic prior to trusting their LaTeX platform ***
% *** with production work. The IEEE's font choices and paper sizes can   ***
% *** trigger bugs that do not appear when using other class files.       ***                          ***
% The testflow support page is at:
% http://www.michaelshell.org/tex/testflow/



\documentclass[journal]{IEEEtran}
% 
% If IEEEtran.cls has not been installed into the LaTeX system files,
% manually specify the path to it like:
% \documentclass[journal]{../sty/IEEEtran}





% Some very useful LaTeX packages include:
% (uncomment the ones you want to load)


% *** MISC UTILITY PACKAGES ***
%
%\usepackage{ifpdf}
% Heiko Oberdiek's ifpdf.sty is very useful if you need conditional
% compilation based on whether the output is pdf or dvi.
% usage:
% \ifpdf
%   % pdf code
% \else
%   % dvi code
% \fi
% The latest version of ifpdf.sty can be obtained from:
% http://www.ctan.org/pkg/ifpdf
% Also, note that IEEEtran.cls V1.7 and later provides a builtin
% \ifCLASSINFOpdf conditional that works the same way.
% When switching from latex to pdflatex and vice-versa, the compiler may
% have to be run twice to clear warning/error messages.






% *** CITATION PACKAGES ***
%
%\usepackage{cite}
% cite.sty was written by Donald Arseneau
% V1.6 and later of IEEEtran pre-defines the format of the cite.sty package
% \cite{} output to follow that of the IEEE. Loading the cite package will
% result in citation numbers being automatically sorted and properly
% "compressed/ranged". e.g., [1], [9], [2], [7], [5], [6] without using
% cite.sty will become [1], [2], [5]--[7], [9] using cite.sty. cite.sty's
% \cite will automatically add leading space, if needed. Use cite.sty's
% noadjust option (cite.sty V3.8 and later) if you want to turn this off
% such as if a citation ever needs to be enclosed in parenthesis.
% cite.sty is already installed on most LaTeX systems. Be sure and use
% version 5.0 (2009-03-20) and later if using hyperref.sty.
% The latest version can be obtained at:
% http://www.ctan.org/pkg/cite
% The documentation is contained in the cite.sty file itself.







% *** GRAPHICS RELATED PACKAGES ***
%
%\ifCLASSINFOpdf
%  % \usepackage[pdftex]{graphicx}
%  % declare the path(s) where your graphic files are
%  % \graphicspath{{../pdf/}{../jpeg/}}
%  % and their extensions so you won't have to specify these with
%  % every instance of \includegraphics
%  % \DeclareGraphicsExtensions{.pdf,.jpeg,.png}
%\else
%  % or other class option (dvipsone, dvipdf, if not using dvips). graphicx
%  % will default to the driver specified in the system graphics.cfg if no
%  % driver is specified.
%  % \usepackage[dvips]{graphicx}
%  % declare the path(s) where your graphic files are
%  % \graphicspath{{../eps/}}
%  % and their extensions so you won't have to specify these with
%  % every instance of \includegraphics
%  % \DeclareGraphicsExtensions{.eps}
%\fi
% graphicx was written by David Carlisle and Sebastian Rahtz. It is
% required if you want graphics, photos, etc. graphicx.sty is already
% installed on most LaTeX systems. The latest version and documentation
% can be obtained at:
% http://www.ctan.org/pkg/graphicx
% Another good source of documentation is "Using Imported Graphics in
% LaTeX2e" by Keith Reckdahl which can be found at:
% http://www.ctan.org/pkg/epslatex
%
% latex, and pdflatex in dvi mode, support graphics in encapsulated
% postscript (.eps) format. pdflatex in pdf mode supports graphics
% in .pdf, .jpeg, .png and .mps (metapost) formats. Users should ensure
% that all non-photo figures use a vector format (.eps, .pdf, .mps) and
% not a bitmapped formats (.jpeg, .png). The IEEE frowns on bitmapped formats
% which can result in "jaggedy"/blurry rendering of lines and letters as
% well as large increases in file sizes.
%
% You can find documentation about the pdfTeX application at:
% http://www.tug.org/applications/pdftex





% *** MATH PACKAGES ***
%
%\usepackage{amsmath}
% A popular package from the American Mathematical Society that provides
% many useful and powerful commands for dealing with mathematics.
%
% Note that the amsmath package sets \interdisplaylinepenalty to 10000
% thus preventing page breaks from occurring within multiline equations. Use:
%\interdisplaylinepenalty=2500
% after loading amsmath to restore such page breaks as IEEEtran.cls normally
% does. amsmath.sty is already installed on most LaTeX systems. The latest
% version and documentation can be obtained at:
% http://www.ctan.org/pkg/amsmath






% *** SPECIALIZED LIST PACKAGES ***
%
%\usepackage{algorithmic}
% algorithmic.sty was written by Peter Williams and Rogerio Brito.
% This package provides an algorithmic environment fo describing algorithms.
% You can use the algorithmic environment in-text or within a figure
% environment to provide for a floating algorithm. Do NOT use the algorithm
% floating environment provided by algorithm.sty (by the same authors) or
% algorithm2e.sty (by Christophe Fiorio) as the IEEE does not use dedicated
% algorithm float types and packages that provide these will not provide
% correct IEEE style captions. The latest version and documentation of
% algorithmic.sty can be obtained at:
% http://www.ctan.org/pkg/algorithms
% Also of interest may be the (relatively newer and more customizable) 
% algorithmicx.sty package by Szasz Janos:
% http://www.ctan.org/pkg/algorithmicx





% *** ALIGNMENT PACKAGES ***
%
%\usepackage{array}
% Frank Mittelbach's and David Carlisle's array.sty patches and improves
% the standard LaTeX2e array and tabular environments to provide better
% appearance and additional user controls. As the default LaTeX2e table
% generation code is lacking to the point of almost being broken with
% respect to the quality of the end results, all users are strongly
% advised to use an enhanced (at the very least that provided by array.sty)
% set of table tools. array.sty is already installed on most systems. The
% latest version and documentation can be obtained at:
% http://www.ctan.org/pkg/array


% IEEEtran contains the IEEEeqnarray family of commands that can be used to
% generate multiline equations as well as matrices, tables, etc., of high
% quality.




% *** SUBFIGURE PACKAGES ***
%\ifCLASSOPTIONcompsoc
%  % \usepackage[caption=false,font=normalsize,labelfont=sf,textfont=sf]{subfig}
%\else
%  % \usepackage[caption=false,font=footnotesize]{subfig}
%\fi
% subfig.sty, written by Steven Douglas Cochran, is the modern replacement
% for subfigure.sty, the latter of which is no longer maintained and is
% incompatible with some LaTeX packages including fixltx2e. However,
% subfig.sty requires and automatically loads Axel Sommerfeldt's caption.sty
% which will override IEEEtran.cls' handling of captions and this will result
% in non-IEEE style figure/table captions. To prevent this problem, be sure
% and invoke subfig.sty's "caption=false" package option (available since
% subfig.sty version 1.3, 2005/06/28) as this is will preserve IEEEtran.cls
% handling of captions.
% Note that the Computer Society format requires a larger sans serif font
% than the serif footnote size font used in traditional IEEE formatting
% and thus the need to invoke different subfig.sty package options depending
% on whether compsoc mode has been enabled.
%
% The latest version and documentation of subfig.sty can be obtained at:
% http://www.ctan.org/pkg/subfig




% *** FLOAT PACKAGES ***
%
%\usepackage{fixltx2e}
% fixltx2e, the successor to the earlier fix2col.sty, was written by
% Frank Mittelbach and David Carlisle. This package corrects a few problems
% in the LaTeX2e kernel, the most notable of which is that in current
% LaTeX2e releases, the ordering of single and double column floats is not
% guaranteed to be preserved. Thus, an unpatched LaTeX2e can allow a
% single column figure to be placed prior to an earlier double column
% figure.
% Be aware that LaTeX2e kernels dated 2015 and later have fixltx2e.sty's
% corrections already built into the system in which case a warning will
% be issued if an attempt is made to load fixltx2e.sty as it is no longer
% needed.
% The latest version and documentation can be found at:
% http://www.ctan.org/pkg/fixltx2e


%
%\usepackage{stfloats}
% stfloats.sty was written by Sigitas Tolusis. This package gives LaTeX2e
% the ability to do double column floats at the bottom of the page as well
% as the top. (e.g., "\begin{figure*}[!b]" is not normally possible in
% LaTeX2e). It also provides a command:
%\fnbelowfloat
% to enable the placement of footnotes below bottom floats (the standard
% LaTeX2e kernel puts them above bottom floats). This is an invasive package
% which rewrites many portions of the LaTeX2e float routines. It may not work
% with other packages that modify the LaTeX2e float routines. The latest
% version and documentation can be obtained at:
% http://www.ctan.org/pkg/stfloats
% Do not use the stfloats baselinefloat ability as the IEEE does not allow
% \baselineskip to stretch. Authors submitting work to the IEEE should note
% that the IEEE rarely uses double column equations and that authors should try
% to avoid such use. Do not be tempted to use the cuted.sty or midfloat.sty
% packages (also by Sigitas Tolusis) as the IEEE does not format its papers in
% such ways.
% Do not attempt to use stfloats with fixltx2e as they are incompatible.
% Instead, use Morten Hogholm'a dblfloatfix which combines the features
% of both fixltx2e and stfloats:
%
% \usepackage{dblfloatfix}
% The latest version can be found at:
% http://www.ctan.org/pkg/dblfloatfix




%
%
%
%
%
%
%
%
%
%
%
%
%
%
%
%
%
%
%
%
%
%
%
%
%
%
%
%
%
%
%
%
%
%
%
%
%
%
%
%
%
%
%
%
%
%
%
%
%
%
%
%
%
%
%
%
%
%
%
%
%
%
%
%
%
%
%
%
%
%
%
%
%
%
%
%
%
%
%
%
%
%
%
%
%
%
%
%
%
%
%
%
%
%
%
%
%
%
%
%
%
%
%
%
%
%
%
%
%
%
%
%
%
%
%
%
%
%
%
%
%
%
%
%
%
%
%
%
%
%
%
%
%
%
%
%
%
%
%
%
%
%
%
%
%
%
%
%
%
%
%
%
%
%
%
%
%
%
%
%
%
%
%
%
%
%
%
%
%
%
%
%
%
%
%
%
%
%
%
%
%
%
%
%
%
%
%
%
%
%
%
%
%
%
%
%
%
%
%
%
%
%
%
%
%
%
%
%
%
%
%
%
%
%
%
%
%
%
%
%
%
%
%
%
%
%
%
%
%
%
%
%
%
%
%
%
%
%
%
%
%
%
%
%
%
%
%
%
%
%
%
%
%
%
%
%
%
%
%
%
%
%
%
%
%
%
%
%
%
%
%
%
%
%
%
%
%
%
%
%
%
%
%
%
%
%
%
%
%
%
%
%
%
%
%
%
%
%
%
%
%
%
%
%
%
%
%
%
%
%
%
%
%
%
%
%
%
%
%
%
%
%
%
%
%
%
%
%
%
%
%
%
%
%
%
%
%
%
%
%
%
%
%
%
%
%
%
%
%
%
%
%
%
%
%
%
%
%
%
%
%
%
%
%
%
%
%
%
%
%
%
%
%
%
%
%
%
%
%
%
%
%
%
%
%
%
%
%
%
%
%
%
%
%
%
%
%
%
%
%
%
%
%
%
%
%
%
%
%
%
%
%
%
%
%
%
%
%
%
%
%
%
%
%
%
%
%
%
%
%
%
%
%
%
%
%
%
%
%
%
%
%
%
%
%
%
%
%
%
%
%
%
%
%
%
%
%
%
%
%
%
%
%
%
%
%
%
%
%
%
%
%
%
%
%
%
%
%
%
%
%
%
%
%
%
%
%
%
%
%
%
%
%
%
%
%
%
%
%
%
%
%
%
%
%
%
%
%
%
%
%
%
%
%
%
%
%
%
%
%
%
%
%
%
%
%
%
%
%
%
%
%
%
%
%
%
%
%
%
%
%
%
%
%
%
%
%
%
%
%
%
%
%
%
%
%
%
%
%
%
%
%
%
%
%
%
%
%
%
%
%
%
%
%
%
%
%
%
%
%
%
%
%
%
%
%
%
%
%
%
%
%
%
%
%
%
%
%
%
%
%
%
%
%
%
%
%
%
%
%
%
%
%
%
%
%
%
%
%
%
%
%
%
%
%
%
%
%
%
%
%
%
%
%
%
%
%
%
%
%
%
%
%
%
%
%
%
%
%
%
%
%
%
%
%
%
%
%
%
%
%
%
%
%
%
%
%
%
%
%
%
%
%
%
%
%
%
%
%
%
%
%
%
%
%
%
%
%
%
%
%
%
%
%
%
%
%
%
%
%
%
%
%
%
%
%
%
%
%
%
%
%
%
%
%
%
%
%
%
%
%
%
%
%
%
%
%
%
%\ifCLASSOPTIONcaptionsoff
  % \usepackage[nomarkers]{endfloat}
% \let\MYoriglatexcaption\caption
% \renewcommand{\caption}[2][\relax]{\MYoriglatexcaption[#2]{#2}}
%\fi
% endfloat.sty was written by James Darrell McCauley, Jeff Goldberg and 
% Axel Sommerfeldt. This package may be useful when used in conjunction with 
% IEEEtran.cls'  captionsoff option. Some IEEE journals/societies require that
% submissions have lists of figures/tables at the end of the paper and that
% figures/tables without any captions are placed on a page by themselves at
% the end of the document. If needed, the draftcls IEEEtran class option or
% \CLASSINPUTbaselinestretch interface can be used to increase the line
% spacing as well. Be sure and use the nomarkers option of endfloat to
% prevent endfloat from "marking" where the figures would have been placed
% in the text. The two hack lines of code above are a slight modification of
% that suggested by in the endfloat docs (section 8.4.1) to ensure that
% the full captions always appear in the list of figures/tables - even if
% the user used the short optional argument of \caption[]{}.
% IEEE papers do not typically make use of \caption[]'s optional argument,
% so this should not be an issue. A similar trick can be used to disable
% captions of packages such as subfig.sty that lack options to turn off
% the subcaptions:
% For subfig.sty:
% \let\MYorigsubfloat\subfloat
% \renewcommand{\subfloat}[2][\relax]{\MYorigsubfloat[]{#2}}
% However, the above trick will not work if both optional arguments of
% the \subfloat command are used. Furthermore, there needs to be a
% description of each subfigure *somewhere* and endfloat does not add
% subfigure captions to its list of figures. Thus, the best approach is to
% avoid the use of subfigure captions (many IEEE journals avoid them anyway)
% and instead reference/explain all the subfigures within the main caption.
% The latest version of endfloat.sty and its documentation can obtained at:
% http://www.ctan.org/pkg/endfloat
%
% The IEEEtran \ifCLASSOPTIONcaptionsoff conditional can also be used
% later in the document, say, to conditionally put the References on a 
% page by themselves.




% *** PDF, URL AND HYPERLINK PACKAGES ***
%
%\usepackage{url}
% url.sty was written by Donald Arseneau. It provides better support for
% handling and breaking URLs. url.sty is already installed on most LaTeX
% systems. The latest version and documentation can be obtained at:
% http://www.ctan.org/pkg/url
% Basically, \url{my_url_here}.




% *** Do not adjust lengths that control margins, column widths, etc. ***
% *** Do not use packages that alter fonts (such as pslatex).         ***
% There should be no need to do such things with IEEEtran.cls V1.6 and later.
% (Unless specifically asked to do so by the journal or conference you plan
% to submit to, of course. )


% *** ADDITIONAL PACKAGES ***
\usepackage[utf8]{inputenc}
\usepackage[T1]{fontenc}
\usepackage{amsmath}
\usepackage{amsfonts}
\usepackage{amssymb}
\usepackage{graphicx}
\usepackage{cite}
\usepackage{algorithm}
\usepackage{algorithmic}
\usepackage{array}
\usepackage{multirow}
\usepackage{booktabs}
\usepackage{tabularx}
\usepackage{subfig}
\usepackage{tikz}
\usetikzlibrary{shapes.geometric,positioning,arrows.meta,calc}
\usepackage{xcolor}
\usepackage{pifont}

% *** HYPERREF PACKAGE - MUST BE LOADED LAST ***
% Sesuai standar IEEE untuk electronic submission dan online publication
\usepackage[hidelinks]{hyperref}  % hidelinks: menghilangkan kotak warna di sekitar link (cocok untuk print/online)
\usepackage{url}  % URL formatting yang lebih baik

% Konfigurasi metadata PDF untuk IEEE submission
\hypersetup{
    bookmarksnumbered=true,
    pdftitle={Restorasi Dokumen Terdegradasi Menggunakan GAN dengan Diskriminator Dual-Modal dan Optimasi Loss Function Berorientasi HTR},
    pdfauthor={[Nama Anda]},
    pdfsubject={Document Image Enhancement, Handwritten Text Recognition, Deep Learning},
    pdfkeywords={Generative Adversarial Networks, Document Restoration, HTR, Dual-Modal Discriminator, Loss Function Optimization},
    pdfproducer={LaTeX with IEEEtran},
    pdfcreator={pdfLaTeX},
    pdfstartview=FitH
}

% correct bad hyphenation here
\hyphenation{op-tical net-works semi-conduc-tor hand-writ-ten doc-u-ment re-sto-ra-tion con-vol-u-tion-a-l re-cur-rent neu-ral net-works recog-ni-tion de-grad-a-tion AR-chi-val ANRI UNESCO}


\begin{document}
% 
% paper title
\title{GAN-Based Document Restoration with Frozen HTR Recognizer and Multi-Component Loss Optimization for Historical Manuscripts}
% 
% author names and affiliations
\author{Jatniko~Nur~Mutaqin
        dan~I~Gusti~Bagus~Baskara~Nugraha,~\IEEEmembership{Member,~IEEE}% <-this % stops a space
\thanks{Jatniko Nur Mutaqin is a graduate student in the \textit{Smart Systems Program}, School of Electrical Engineering and Informatics, Bandung Institute of Technology, Bandung 40132, Indonesia, and works at the National Archives of the Republic of Indonesia, Jakarta, Indonesia (e-mail: 23523314@mahasiswa.itb.ac.id).}% <-this % stops a space
\thanks{I Gusti Bagus Baskara Nugraha is a faculty member at the School of Electrical Engineering and Informatics, Bandung Institute of Technology, Bandung 40132, Indonesia (e-mail: baskara@itb.ac.id).}% <-this % stops a space
\thanks{Manuscript received [Tanggal]; revised [Tanggal].}}

% The paper headers
\markboth{IEEE Transactions on [Nama Jurnal],~Vol.~XX, No.~X, [Bulan]~20XX}% 
{Mutaqin \MakeLowercase{\textit{et al.}}: GAN-Based Document Restoration with Frozen HTR Recognizer for Historical Manuscripts}

% make the title area
\maketitle

% Abstract
\begin{abstract}
Dokumen tulisan tangan historis mengalami berbagai artefak degradasi yang mempengaruhi sistem \textit{Handwritten Text Recognition} (HTR). Metode restorasi dokumen konvensional berfokus pada kualitas visual namun seringkali gagal mempertahankan keterbacaan teks untuk HTR. Makalah ini mengusulkan kerangka kerja restorasi dokumen berorientasi HTR berbasis \textit{Generative Adversarial Networks} (GAN) dengan dua kontribusi utama: (1) integrasi \textit{recognizer} HTR beku yang memberikan gradien sadar-teks stabil tanpa ketidakstabilan pelatihan bersama, dan (2) optimasi fungsi loss multikomponen melalui studi ablasi sistematis. Berbeda dengan pendekatan konvensional yang hanya menggunakan supervisi tingkat citra, metode yang diusulkan secara eksplisit mengoptimalkan keterbacaan teks melalui gradien dari \textit{recognizer} praterlatih yang dibekukan. Evaluasi kuantitatif dilakukan pada \textit{dataset} semi-sintetis realistis yang dikonstruksi dari dokumen paleografi autentik Arsip Nasional Republik Indonesia (ANRI) abad ke-16 hingga ke-18 dengan augmentasi degradasi terkontrol untuk memungkinkan evaluasi dengan \textit{ground truth} berpasangan. Pada 712 citra uji, kerangka kerja yang diusulkan mencapai \textit{Character Error Rate} (CER) 34.9\%, menurunkan CER dari 83.4\% (kondisi terdegradasi) dengan pengurangan relatif 58.2\%, sambil mempertahankan kualitas visual tinggi (PSNR 30.74~dB, SSIM 0.987). Hasil ini mendekati kinerja teoretis maksimal pada citra bersih (CER 34.1\%). Validasi kualitatif pada 15 naskah ANRI nyata yang mengalami degradasi historis autentik menunjukkan generalisasi efektif pada pola kerusakan yang tidak dapat direplikasi secara sintetis. Studi ablasi mengungkapkan bahwa konfigurasi 4-komponen kehilangan (piksel, \textit{adversarial}, perseptual, CTC) mencapai keseimbangan optimal, sementara eksplorasi diskriminator \textit{dual-modal} menunjukkan kontribusi marginal ($\Delta$PSNR +0.28~dB, $p>0.05$), mengonfirmasi bahwa strategi pembekuan \textit{recognizer} dan optimasi kehilangan lebih dominan daripada kompleksitas arsitektur diskriminator.
\end{abstract}

% Keywords
\begin{IEEEkeywords}
Generative Adversarial Networks, Restorasi Dokumen, \textit{Handwritten Text Recognition}, Optimisasi Fungsi Loss, Pemrosesan Dokumen Historis, \textit{Deep Learning}.
\end{IEEEkeywords}







% For peer review papers, you can put extra information on the cover
% page as needed:
% \ifCLASSOPTIONpeerreview
% \begin{center} \bfseries EDICS Category: 3-BBND \end{center}
% \fi
% 
% For peerreview papers, this IEEEtran command inserts a page break and
% creates the second title. It will be ignored for other modes.
\IEEEpeerreviewmaketitle

% ============================================================ 
% I. PENDAHULUAN
% ============================================================ 
\section{Pendahuluan}
\IEEEPARstart{D}{okumen} tulisan tangan historis merupakan warisan budaya penting dan mengandung informasi tentang peristiwa masa lalu, pemerintahan, perdagangan, dan struktur sosial. Arsip Nasional Republik Indonesia (ANRI) menyimpan koleksi ekstensif dokumen era kolonial dari abad ke-16 hingga ke-18, termasuk yang diakui oleh UNESCO sebagai bagian dari Daftar Ingatan Dunia (\textit{Memory of the World Register}).

Dokumen-dokumen ini, yang ditulis dalam aksara paleografi Belanda selama periode VOC (\textit{Vereenigde Oostindische Compagnie}) dan Hindia Belanda, memberikan wawasan unik tentang perkembangan sejarah Indonesia. Namun, dokumen-dokumen ini mengalami degradasi parah yang mencakup (Gambar~\ref{fig:degradation_examples}): tinta tembus (\textit{bleed-through}), korosi tinta besi galat, penuaan kertas, kerusakan fisik, dan degradasi tinta yang memudar.

Degradasi ini menghambat keterbacaan teks dan menantang sistem \textit{Handwritten Text Recognition} (HTR).

Artefak degradasi ini mengganggu kinerja sistem HTR, yang penting untuk digitalisasi skala besar, transkripsi, pencarian, pengindeksan, dan ekstraksi informasi dari arsip historis. Meskipun kemajuan substansial telah dibuat dalam teknologi HTR menggunakan pendekatan \textit{deep learning}, terutama dengan \textit{Convolutional Recurrent Neural Networks} (CRNN) dan arsitektur berbasis transformer, sistem ini tetap sensitif terhadap kualitas citra dan tingkat degradasi.

% [PLACEHOLDER: Gambar 1 - Contoh jenis degradasi pada dokumen historis]
% - Beberapa subgambar menunjukkan: (a) noise latar belakang, (b) tinta luntur
% - (c) penuaan/penguningan kertas
% - (d) kerusakan fisik, (e) Degradasi Tinta, (f) degradasi gabungan

\begin{figure}[!t]
\centering
\subfloat[Tinta Tembus]{\label{fig:degradation_examples_noise}%
\includegraphics[width=1.7in]{fig_degradation_examples/ink_bleed_through}}\hfill
\subfloat[Korosi Tinta Besi Galat]{\label{fig:degradation_examples_inkCorrotion}%
\includegraphics[width=1.7in]{fig_degradation_examples/ink_corrosion}}\hfill
\subfloat[Penuaan Kertas (Menguning, Noda)]{\label{fig:degradation_examples_aging}%
\includegraphics[width=1.7in]{fig_degradation_examples/paper_aging}}\hfill
\subfloat[Kerusakan Fisik (Robek, Keriput)]{\label{fig:degradation_examples_damage}%
\includegraphics[width=1.7in]{fig_degradation_examples/physical_damage}}\hfill
\subfloat[Degradasi atau Tinta Memudar]{\label{fig:degradation_examples_inkDeg}%
\includegraphics[width=1.7in]{fig_degradation_examples/ink_degradation}}\hfill
\subfloat[Degradasi Gabungan]{\label{fig:degradation_examples_combined}%
\includegraphics[width=1.7in]{fig_degradation_examples/combined_degradation}}
\caption{Contoh degradasi pada dokumen ANRI abad 16-18: (a) tinta tembus, (b) korosi tinta besi galat, (c) penuaan kertas, (d) kerusakan fisik, (e) tinta memudar, (f) degradasi gabungan. Gambar menunjukkan variasi tingkat degradasi yang menantang sistem HTR.}
\label{fig:degradation_examples}
\end{figure}

\subsection{Keterbatasan Pendekatan yang Ada}

Dokumen historis yang terdegradasi menghadirkan tantangan unik untuk sistem \textit{Handwritten Text Recognition} (HTR) karena derau, noda, ketidakfokusan, dan degradasi latar belakang dapat menurunkan akurasi pengenalan. Restorasi dokumen otomatis bertujuan untuk meningkatkan kualitas citra sebelum HTR, namun pendekatan yang ada memiliki keterbatasan fundamental:

Pertama, metode restorasi dokumen konvensional (termasuk \textit{thresholding} klasik seperti Otsu~\cite{otsu1979threshold} dan Sauvola~\cite{sauvola2000adaptive}, serta pendekatan deep learning berbasis U-Net dan GAN) mengoptimalkan metrik visual semata (PSNR, SSIM) tanpa mempertimbangkan dampak terhadap keterbacaan HTR. Akibatnya, model dapat menghasilkan citra yang secara visual menarik namun merusak struktur teks (misalnya goresan terputus, penghubung karakter hilang) yang krusial untuk akurasi HTR.

Kedua, integrasi HTR dalam restorasi dokumen telah diusulkan oleh Souibgui \textit{et al.}~\cite{souibgui2021enhance} melalui pelatihan bersama antara generator dan pengenal. Meskipun inovatif, pendekatan ini menggunakan diskriminator single-modal yang hanya mengevaluasi kualitas visual, dan pelatihan bersama dapat menyebabkan ketidakstabilan optimisasi karena generator dan pengenal berkompetisi untuk gradien yang konflik.

Ketiga, kesenjangan evaluasi—sebagian besar benchmark (seri DIBCO) hanya menyediakan ground truth visual tanpa transkripsi teks, membuat validasi dampak HTR tidak dapat dilakukan tanpa anotasi tambahan. Dataset untuk dokumen paleografi abad 16-18 dengan degradasi autentik terbatas.

Penelitian ini mengatasi keterbatasan tersebut dengan kontribusi sebagai berikut. \textit{Kontribusi utama}: (1) Strategi \textit{recognizer} HTR beku yang memberikan gradien sadar-teks stabil tanpa ketidakstabilan pelatihan bersama ($p<0.001$); (2) Optimasi fungsi loss multikomponen melalui studi ablasi sistematis yang mengidentifikasi konfigurasi 4-komponen optimal (piksel, \textit{adversarial}, perseptual, CTC). \textit{Kontribusi metodologis}: (3) Evaluasi pada naskah paleografi nyata abad 16-18 dari Arsip Nasional Republik Indonesia; (4) Studi ablasi yang mengukur kontribusi setiap komponen, termasuk temuan bahwa diskriminator \textit{dual-modal} memberikan kontribusi marginal ($\Delta$PSNR +0.28~dB, $p>0.05$), mengonfirmasi protokol pelatihan dan kapasitas generator lebih dominan daripada kompleksitas diskriminator.

\subsection{Kontribusi Penelitian}

Penelitian ini mengusulkan kerangka kerja restorasi dokumen berorientasi HTR yang secara eksplisit mengoptimalkan keterbacaan teks sambil mempertahankan kualitas visual. Berdasarkan validasi empiris melalui studi ablasi sistematis, kontribusi penelitian dikelompokkan sebagai berikut:

\textbf{Kontribusi Utama:}
\begin{enumerate}
    \item \textbf{Strategi \textit{Recognizer} HTR Beku}: Penggunaan \textit{recognizer} HTR praterlatih yang dibekukan memberikan gradien sawar-teks stabil tanpa ketidakstabilan pelatihan bersama. Studi ablasi menunjukkan dampak signifikan terhadap stabilitas konvergensi ($p<0.001$) dan pelestarian kapabilitas pengenalan (lihat Bagian~V-B).
    
    \item \textbf{Optimasi Fungsi Loss Multikomponen}: Identifikasi konfigurasi 4-komponen optimal (\textit{adversarial}, rekonstruksi piksel, perseptual, CTC) melalui studi ablasi sistematis. Konfigurasi ini mencapai keseimbangan \textit{Pareto} antara kualitas visual (PSNR 30.74~dB) dan keterbacaan HTR (CER 34.9\%), dengan kinerja lebih baik dibanding konfigurasi 5-komponen (lihat Bagian~V-B).
\end{enumerate}

\textbf{Kontribusi Metodologis dan Temuan Eksplorasi:}
\begin{enumerate}
    \setcounter{enumi}{2}
    \item \textbf{Studi Ablasi Sistematis}: Evaluasi kontribusi setiap komponen arsitektur dan fungsi loss, memberikan panduan praktis untuk pemilihan komponen. Temuan: konfigurasi sederhana dengan penyelarasan tepat mencapai kinerja lebih baik dibanding konfigurasi kompleks berlebihan.
    
    \item \textbf{Evaluasi pada Dokumen Historis Autentik}: Evaluasi kuantitatif pada \textit{dataset} semi-sintetis realistis yang dikonstruksi dari tulisan tangan paleografi autentik ANRI abad ke-16 hingga ke-18 dengan augmentasi degradasi terkontrol ($n$=712), memungkinkan pengukuran presisi dengan \textit{ground truth} berpasangan. Validasi kualitatif pada naskah ANRI nyata yang mengalami degradasi historis autentik ($n$=15) menunjukkan generalisasi efektif pada pola kerusakan yang tidak dapat direplikasi secara komputasional.
    
    \item \textbf{Temuan tentang Diskriminator \textit{Dual-Modal}}: Eksplorasi arsitektur diskriminator dengan jalur CNN (visual) dan LSTM (tekstual) mengungkapkan kontribusi marginal terhadap metrik akhir ($\Delta$PSNR = +0.28~dB, $\Delta$CER = --0.01\%, $p>0.05$, tidak signifikan secara statistik). Temuan ini mengonfirmasi bahwa faktor kunci adalah protokol pelatihan dan strategi \textit{recognizer} beku, bukan kompleksitas arsitektur diskriminator (lihat analisis di Bagian~VI-A).
\end{enumerate}

\subsection{Hipotesis dan Validasi Penelitian}

Penelitian ini memvalidasi hipotesis bahwa kerangka kerja restorasi dokumen dengan integrasi \textit{recognizer} HTR beku dan optimasi fungsi loss berorientasi HTR menghasilkan penurunan \textit{Character Error Rate} (CER) yang signifikan secara statistik dibandingkan kondisi terdegradasi \textit{baseline} (H1). Studi ablasi sistematis dilakukan untuk mengukur kontribusi setiap komponen, dengan temuan utama: (1) strategi pembekuan \textit{recognizer} terbukti menjadi faktor kunci stabilitas dan konvergensi, (2) konfigurasi 4-komponen kehilangan mencapai keseimbangan optimal, dan (3) diskriminator \textit{dual-modal} menunjukkan kontribusi marginal ($p>0.05$). Temuan ini mengonfirmasi bahwa keberhasilan sistem terutama bergantung pada strategi pembekuan \textit{recognizer} dan optimasi kehilangan, bukan kompleksitas arsitektur diskriminator. Hasil lengkap dan analisis statistik disajikan dalam Bagian~V.

\subsection{Organisasi Makalah}

Sistematika penulisan dari makalah ini diatur sebagai berikut: Bagian~II meninjau pekerjaan terkait dalam perbaikan dokumen, GAN untuk restorasi citra, dan sistem HTR. Bagian~III menyajikan arsitektur yang diusulkan termasuk \textit{generator} U-Net Enhanced, diskriminator \textit{dual-modal} (CNN+LSTM), integrasi \textit{recognizer} beku, dan fungsi loss yang dioptimalkan dengan penyeimbangan adaptif. Bagian~IV menjelaskan pengaturan eksperimental, \textit{dataset}, metodologi pelatihan, dan strategi optimisasi hiperparameter. Bagian~V menyajikan hasil komprehensif termasuk perbandingan kuantitatif dengan masukan terdegradasi \textit{baseline} dan referensi \textit{ground truth} bersih, studi ablasi sistematis untuk memvalidasi kontribusi setiap komponen, dan analisis kualitatif pada dokumen historis nyata. Bagian~VI membahas temuan, implikasi teoretis dan praktis, keterbatasan, dan arah masa depan. Akhirnya, Bagian~VII menyimpulkan makalah ini dengan rangkuman kontribusi dan dampak penelitian.

% ============================================================ 
% II. PEKERJAAN TERKAIT
% ============================================================ 
\section{Pekerjaan Terkait}

Bagian ini meninjau penelitian sebelumnya di tiga bidang yang saling berhubungan: restorasi dan perbaikan citra dokumen, generative adversarial networks untuk restorasi citra, dan sistem pengenalan teks tulisan tangan.

\subsection{Restorasi dan Perbaikan Citra Dokumen}

Restorasi dokumen bertujuan untuk memulihkan dokumen terdegradasi dengan menghilangkan noise, artefak, dan degradasi latar belakang sambil mempertahankan keterbacaan teks untuk sistem OCR/HTR atau pembacaan manusia. Berbeda dengan binerisasi yang menghasilkan keluaran biner (hitam-putih), restorasi dokumen dapat menghasilkan keluaran grayscale atau berwarna yang mempertahankan informasi intensitas piksel untuk kualitas visual yang lebih baik.

\subsubsection{Metode Klasik: Binerisasi dan \textit{Thresholding}}

Pendekatan awal untuk perbaikan dokumen berfokus pada binerisasi menggunakan thresholding. Metode binerisasi global seperti Otsu~\cite{otsu1979threshold} menentukan satu nilai ambang untuk seluruh citra dengan memaksimalkan varians antarkelas. Meskipun efisien secara komputasi, metode global cenderung kurang efektif pada dokumen dengan pencahayaan yang bervariasi atau degradasi kompleks.

Metode thresholding adaptif lokal mengatasi keterbatasan ini dengan menghitung ambang spesifik piksel berdasarkan lingkungan lokal. Sauvola dan Pietikäinen~\cite{sauvola2000adaptive} mengusulkan metode yang banyak digunakan yang menggabungkan rata-rata lokal dan deviasi standar. Niblack~\cite{niblack1986introduction} dan Wolf~\cite{wolf2004extraction} mengembangkan pendekatan lokal serupa dengan strategi adaptasi yang berbeda. Namun, metode berbasis thresholding ini menghasilkan keluaran biner yang kehilangan informasi intensitas piksel dan dapat menghilangkan detail halus seperti goresan tipis.

\subsubsection{Metode Deep Learning untuk Restorasi Dokumen}

Sejak pertengahan 2010-an, dengan kemajuan deep learning dan ketersediaan dataset berlabel besar seperti DIBCO~\cite{pratikakis2016hdibco}, pendekatan restorasi dokumen beralih dari binerisasi ke peningkatan grayscale yang mempertahankan informasi intensitas. Auto-encoder berbasis U-Net~\cite{ronneberger2015unet} dan Generative Adversarial Networks (GAN) memungkinkan pembelajaran pemetaan end-to-end dari dokumen terdegradasi ke dokumen bersih dengan potensi mempertahankan detail halus yang penting untuk HTR.

\subsection{Generative Adversarial Networks untuk Restorasi Citra}

GAN~\cite{goodfellow2016improved} telah mengubah secara substansial bidang restorasi citra dengan kemampuannya menghasilkan keluaran yang realistis melalui pelatihan adversarial antara generator dan discriminator. Berbeda dengan auto-encoder yang hanya meminimalkan kehilangan rekonstruksi, GAN dapat menghasilkan detail yang lebih tajam dan tekstur yang lebih alami.

\subsubsection{\textit{Conditional GANs} untuk Translasi Citra-ke-Citra}

Pix2pix~\cite{isola2017image} memperkenalkan kerangka kerja \textit{conditional GAN} (cGAN) untuk tugas translasi citra-ke-citra, di mana \textit{generator} menerima citra masukan dan menghasilkan citra keluaran yang sesuai. \textit{Discriminator} PatchGAN dalam Pix2pix mengevaluasi keaslian pada \textit{patch} lokal, memungkinkan penangkapan detail frekuensi tinggi yang lebih baik. Kerangka kerja ini telah berhasil diterapkan pada berbagai domain termasuk pewarnaan, transfer gaya, dan segmentasi semantik.

CycleGAN~\cite{zhu2017unpaired} memperluas konsep ini ke translasi citra-ke-citra tanpa pasangan menggunakan kehilangan konsistensi siklus, memungkinkan pelatihan tanpa data berpasangan. Meskipun berguna untuk adaptasi domain, pendekatan tanpa pasangan kurang cocok untuk restorasi dokumen di mana data terdegradasi-bersih berpasangan tersedia dan akurasi tingkat piksel krusial untuk pelestarian konten teks.

\subsubsection{GAN untuk Restorasi dan Perbaikan Dokumen}

Dalam konteks restorasi dokumen, beberapa pendekatan GAN telah diusulkan:

DE-GAN~\cite{souibgui2020degan} menggunakan generator U-Net dengan discriminator PatchGAN untuk tugas peningkatan dokumen termasuk penghilangan blur, penghapusan watermark, dan pembersihan latar belakang. Meskipun mencapai kualitas visual tinggi, DE-GAN berfokus pada optimisasi kemiripan tingkat piksel dan tidak secara eksplisit mengintegrasikan metrik keterbacaan HTR dalam fungsi lossnya.

DocEnTr~\cite{souibgui2022docent} memperkenalkan transformer peningkatan dokumen yang menggabungkan mekanisme atensi mandiri dengan pelatihan adversarial. Pendekatan ini menunjukkan kemampuan menangkap dependensi jarak jauh, namun biaya komputasi tinggi dan belum mengintegrasikan panduan HTR eksplisit.

HTR-GAN (Souibgui et al.)~\cite{souibgui2021enhance} merupakan karya paling relevan dengan penelitian ini, mengintegrasikan pengenal HTR dalam kerangka kerja GAN untuk peningkatan dokumen. Generator dan pengenal dilatih bersama dengan kehilangan gabungan dari adversarial, rekonstruksi, dan CTC. Pendekatan pelatihan bersama berpotensi menyebabkan ketidakstabilan optimisasi karena pengenal dan generator berkompetisi untuk gradien—kekhawatiran teoretis yang memotivasi pendekatan pengenal beku dalam penelitian ini.

Perbedaan Pendekatan yang Diusulkan: Penelitian ini berbeda fundamental dari HTR-GAN dalam hal: (1) Pengenal beku\textemdash penelitian ini menggunakan pengenal praterlatih yang dibekukan untuk gradien stabil, bukan pelatihan bersama; (2) Discriminator dual-modal\textemdash discriminator yang diusulkan memproses citra (CNN) dan urutan teks (LSTM) secara paralel dengan potensi untuk evaluasi koherensi visual dan tekstual (meskipun studi ablasi menunjukkan dampak marginal, lihat Bagian~VI); (3) Penyeimbangan kehilangan adaptif\textemdash mekanisme penyeimbang adaptif yang secara dinamis menyesuaikan rasio antara kehilangan CTC dan kehilangan visual untuk keseimbangan optimal; (4) Arsitektur yang ditingkatkan\textemdash residual blocks, gerbang atensi, dan atensi spasial untuk pelestarian goresan tipis yang lebih baik.

\subsection{Sistem Pengenalan Teks Tulisan Tangan}

Sistem HTR bertujuan untuk mengenali teks tulisan tangan dari citra dokumen. Evolusi HTR telah berkembang dari metode berbasis HMM klasik ke arsitektur deep learning modern.

CRNN (Convolutional Recurrent Neural Network)~\cite{shi2016crnn} menggabungkan CNN untuk ekstraksi fitur dengan RNN (umumnya LSTM/GRU) untuk pemodelan urutan, dilatih dengan kehilangan CTC untuk pengenalan bebas penjajaran. Arsitektur ini menjadi fondasi untuk banyak sistem HTR modern karena kemampuannya menangani urutan dengan panjang variabel tanpa segmentasi eksplisit.

Mekanisme Atensi telah meningkatkan kinerja HTR dengan memungkinkan model fokus pada wilayah relevan saat mengenali setiap karakter. Arsitektur encoder-decoder dengan atensi~\cite{bluche2017scan} dapat mempelajari segmentasi karakter dan penjajaran secara implisit.

HTR Berbasis Transformer~\cite{michael2019evaluating,diaz2021rethinking} memanfaatkan atensi mandiri untuk menangkap dependensi jarak jauh tanpa koneksi rekuren. Vision transformers dan arsitektur hibrida CNN-Transformer menunjukkan kinerja terkini pada berbagai tolok ukur tulisan tangan, termasuk dokumen historis dan modern.

Relevansi untuk Penelitian Ini: Sistem HTR sensitif terhadap kualitas citra\textemdash degradasi seperti noise, blur, atau goresan putus dapat meningkatkan CER. Penelitian ini menggunakan pengenal hibrida CNN-Transformer praterlatih yang dibekukan untuk mengekstrak fitur pengenalan yang memandu generator menghasilkan keluaran yang dapat dibaca oleh HTR. Pendekatan pengenal beku berbeda dari pelatihan bersama karena: (1) gradien stabil\textemdash bobot pengenal tidak berubah; (2) evaluasi konsisten\textemdash kinerja pengenal tidak menurun selama pelatihan GAN; (3) desain modular\textemdash pengenal dapat diperbarui secara independen.

\subsection{Analisis Kesenjangan dan Posisi}

Berdasarkan tinjauan literatur di atas, Tabel~\ref{table_comparison} merangkum posisi penelitian ini terhadap pendekatan yang ada dalam hal kualitas visual, kesadaran HTR, arsitektur discriminator, dan optimisasi kehilangan. Kesenjangan yang diidentifikasi adalah: (1) eksplorasi arsitektur dual-modal (visual + tekstual) dalam discriminator sebagai kelengkapan arsitektural untuk validasi bilateral---meskipun studi ablasi menunjukkan dampak marginal terhadap metrik akhir (lihat Bagian~VI), (2) tidak adanya optimisasi bobot kehilangan sistematis untuk menyeimbangkan kualitas visual dan keterbacaan HTR, dan (3) potensi ketidakstabilan dari pelatihan bersama pengenal-generator yang dapat diatasi dengan pendekatan pengenal beku, sebagaimana dimotivasi oleh kekhawatiran teoretis yang dibahas di Bagian~II. Validasi empiris dari kontribusi relatif masing-masing komponen disajikan dalam studi ablasi komprehensif (Bagian~V), dengan Gap \#2 (optimisasi bobot kehilangan sistematis) menunjukkan dampak paling signifikan terhadap kinerja akhir.

\begin{table}[!ht]
\renewcommand{\arraystretch}{1.3}
\caption{Perbandingan Pendekatan Perbaikan Dokumen: evaluasi metode terkini berdasarkan kualitas visual, kesadaran HTR, arsitektur discriminator dual-modal, dan optimisasi loss function. Pendekatan yang diusulkan mengintegrasikan semua fitur untuk restorasi dokumen berorientasi HTR, dengan kontribusi relatif dari setiap komponen divalidasi melalui studi ablasi (Bagian V--VI).}
\label{table_comparison}
\centering
\begin{tabular}{|l|c|c|c|c|}
\hline
\textbf{Metode} & \textbf{Kualitas} & \textbf{Sadar} & \textbf{Disk.} & \textbf{Optim.} \\
 & \textbf{Visual} & \textbf{HTR} & \textbf{Dual} & \textbf{Kehilangan} \\
\hline
Klasik & Rendah & Tidak & Tidak & T/A \\
U-Net & Tinggi & Tidak & Tidak & Tidak \\
DE-GAN & Tinggi & Tidak & Tidak & Tidak \\
Souibgui et al. & Tinggi & Ya & Tidak & Tidak \\
\textbf{Yang Diusulkan} & \textbf{Tinggi} & \textbf{Ya} & \textbf{Ya} & \textbf{Ya} \\
\hline
\end{tabular}
\end{table}

% ============================================================ 
% IV. METODE YANG DIUSULKAN
% ============================================================ 
\section{Metode yang Diusulkan}

Bagian ini menyajikan kerangka kerja restorasi dokumen berorientasi HTR yang diusulkan. Pertama disajikan gambaran arsitektur, kemudian merinci setiap komponen: generator U-Net, diskriminator dual-modal, integrasi pengenal HTR yang dibekukan, dan fungsi loss multikomponen yang dioptimalkan.

\subsection{Gambaran Kerangka Kerja}

Untuk mengatasi keterbatasan metode restorasi konvensional yang mengabaikan keterbacaan teks mesin, penelitian ini mengusulkan kerangka kerja GAN berorientasi HTR yang secara eksplisit mengintegrasikan pengenal teks dalam loop pelatihan. Kerangka kerja yang diusulkan, yang diilustrasikan pada Gambar~\ref{fig:architecture_overview}, beroperasi pada citra baris dokumen tulisan tangan yang terdegradasi dan menghasilkan keluaran yang bersih dan dapat dibaca yang dioptimalkan untuk kualitas visual dan kinerja HTR.

% Diagram arsitektur kerangka kerja restorasi dokumen berorientasi HTR
\begin{figure*}[!t]
\centering
\includegraphics[width=0.95\textwidth,keepaspectratio]{../drawio/main_framework.png}

\caption{Kerangka kerja restorasi dokumen berorientasi HTR dengan discriminator dual-modal dan frozen HTR recognizer. Tiga komponen utama: (1) Generator U-Net Enhanced untuk restorasi citra, (2) Discriminator Dual-Modal (CNN+BiLSTM) untuk validasi visual-tekstual bilateral, (3) Frozen HTR Recognizer untuk gradient HTR-aware. Training menggunakan curriculum learning dengan multi-component loss (pixel, perceptual, adversarial, recognition feature, HTR). Total parameter yang dilatih: 108.3M (Generator 89.4M + Discriminator 18.9M), Recognizer frozen: 27.86M.}
\label{fig:architecture_overview}
\end{figure*}

Alur proses restorasi berlangsung sebagai berikut: Citra terdegradasi $I_{\text{deg}}$ dimasukkan ke generator $G$ yang menghasilkan citra terestorasi $I_{\text{gen}}$. Diskriminator dual-modal $D$ menerima pasangan $(I_{\text{gt}}, \text{teks}_{\text{gt}})$ sebagai sampel nyata dan $(I_{\text{gen}}, \text{teks}_{\text{pred}})$ sebagai sampel palsu untuk pembelajaran adversarial. Secara paralel, pengenal HTR beku $R$ mengekstrak fitur dari $I_{\text{gen}}$ untuk memastikan keterbacaan teks melalui kehilangan CTC dan kehilangan fitur pengenalan.

Kerangka kerja ini terdiri dari empat komponen utama:

\begin{enumerate}
    \item Generator (G): Arsitektur U-Net Enhanced yang memetakan citra terdegradasi $I_{\text{deg}}$ ke citra yang direstorasi $I_{\text{gen}}$. Generator dilengkapi Residual Dense Blocks dan Attention Gates untuk preservasi detail teks halus sambil menghilangkan noise latar belakang.
    
    \item Diskriminator Dual-Modal (D): Diskriminator dengan jalur CNN dan BiLSTM paralel yang menilai kualitas visual dan koherensi teks secara bersamaan. Mekanisme bilateral cross-attention fusion memungkinkan validasi konsistensi visual-tekstual untuk menghindari artefak yang merusak struktur karakter.
    
    \item \textit{Recognizer} HTR Beku (R): \textit{Recognizer} berbasis \textit{Transformer} praterlatih dengan bobot tetap yang menyediakan gradien sadar HTR tanpa risiko \textit{catastrophic forgetting}. \textit{Recognizer} mencapai CER 33.72\% pada \textit{dataset} validasi awal ANRI ($n$=948), dan CER 26.57\% (\textit{set} validasi, $n$=710) serta 34.1\% (\textit{set} uji, $n$=712) pada citra bersih \textit{ground truth} penelitian ini\textemdash perbedaan mencerminkan variasi kompleksitas teks antar \textit{split dataset}.
    
    \item Fungsi Loss Multikomponen: Kombinasi terbobot dari $\mathcal{L}_{\text{pixel}}$ ($\lambda=50$), $\mathcal{L}_{\text{adv}}$ ($\lambda=3$), $\mathcal{L}_{\text{perc}}$ ($\lambda=1$), $\mathcal{L}_{\text{rec\_feat}}$ ($\lambda=8$), dan $\mathcal{L}_{\text{CTC}}$ ($\lambda=0.15$) yang dioptimalkan melalui curriculum learning untuk menyeimbangkan kualitas visual dan keterbacaan teks.
\end{enumerate}

Berbeda dari HTR-GAN~\cite{souibgui2021enhance} yang melatih pengenal bersama generator (joint training), metode yang diusulkan menggunakan pengenal beku untuk stabilitas pelatihan. Arsitektur diskriminator dual-modal sebagai eksplorasi untuk validasi bilateral visual-tekstual yang tidak ada pada metode sebelumnya yang hanya mengandalkan diskriminator visual tunggal.

Proses pelatihan mengadopsi strategi adversarial bergantian dengan curriculum learning tiga fase (detail di Bagian III-F):

\begin{itemize}
    \item Pembaruan diskriminator: $D$ dilatih untuk membedakan pasangan $(I_{\text{gt}}, \text{teks}_{\text{gt}})$ nyata dari pasangan $(I_{\text{gen}}, \text{teks}_{\text{pred}})$ yang dihasilkan, dengan supervisi dual-modal pada domain visual dan tekstual secara bersamaan.
    
    \item Pembaruan generator: $G$ dioptimalkan untuk menipu $D$ sambil meminimalkan $\mathcal{L}_{\text{pixel}}$ terhadap ground truth, $\mathcal{L}_{\text{perc}}$ untuk preservasi struktur perseptual, $\mathcal{L}_{\text{rec\_feat}}$ untuk konsistensi fitur HTR, dan $\mathcal{L}_{\text{CTC}}$ untuk keterbacaan teks eksplisit.
\end{itemize}

\subsection{Arsitektur Generator: U-Net Enhanced dengan Blok Residual}

Penelitian ini mengadopsi arsitektur U-Net Enhanced yang menggabungkan empat teknik restorasi dokumen berkualitas tinggi. Berbeda dari U-Net konvensional yang hanya menggunakan blok konvolusional standar, arsitektur Enhanced mengintegrasikan: (1) Residual Dense Blocks (RDB)~\cite{zhang2018residual} untuk ekstraksi fitur hierarkis dengan aliran gradien yang lebih baik, (2) Convolutional Block Attention Module (CBAM)~\cite{woo2018cbam} untuk rekalibrasi fitur adaptif pada dimensi kanal dan spasial, (3) Multi-Scale Feature Pyramid (MSFP)~\cite{lin2017fpn} dengan konvolusi terdilasi~\cite{chen2017deeplab} pada bottleneck untuk agregasi konteks multi-resolusi, dan (4) Attention Gates~\cite{oktay2018attention} pada decoder untuk fokus selektif pada wilayah pembawa teks sambil menekan amplifikasi noise latar belakang. Kombinasi teknik-teknik ini memungkinkan pelestarian goresan tipis dan tanda diakritik paleografi sambil menghilangkan degradasi kompleks seperti bleed-through, noda, dan derau.

\subsubsection{Spesifikasi Jaringan}

Tabel~\ref{tab:generator_arch} merinci arsitektur lengkap generator. Encoder secara progresif melakukan downsampling pada citra masukan melalui lima tahap, mengekstrak fitur multi-skala pada resolusi dari $1024 \times 128$ hingga $32 \times 4$ (format: lebar $\times$ tinggi). Setiap blok encoder terdiri dari:

\begin{itemize}
    \item Residual Dense Block (RDB): Tiga lapisan konvolusional densely-connected (kernel $3 \times 3$) dengan growth rate 32, menghasilkan fitur kaya melalui fusi lokal dari semua lapisan sebelumnya. Berbeda dari ResNet~\cite{he2016deep} standar yang hanya memiliki satu koneksi shortcut, RDB menggunakan koneksi padat yang memungkinkan setiap lapisan mengakses fitur dari semua lapisan sebelumnya, meningkatkan kapasitas representasi dan aliran gradien.
    \item Convolutional Block Attention Module (CBAM): Mekanisme atensi sekuensial yang terdiri dari: (a) atensi kanal dengan rasio reduksi 8, menggunakan shared MLP untuk jalur global average pooling dan global max pooling, dan (b) atensi spasial dengan kernel $7 \times 7$, menggabungkan channel-wise average dan max pooling untuk menghasilkan peta atensi fokus. CBAM memungkinkan jaringan secara adaptif memfokuskan pada kanal fitur yang informatif dan wilayah spasial yang relevan.
    \item Batch normalization (momentum 0.9, eps $10^{-5}$) dan aktivasi LeakyReLU ($\alpha=0.2$)
    \item Koneksi residual untuk propagasi gradien efektif pada arsitektur dalam
\end{itemize}

Bottleneck pada resolusi $4 \times 32$ (tinggi $\times$ lebar) menggunakan Multi-Scale Feature Pyramid (MSFP)~\cite{lin2017fpn} dengan tiga skala konvolusi terdilasi (dilation rates 1, 2, 4)~\cite{chen2017deeplab} untuk menangkap konteks multi-resolusi, diikuti RDB 512 kanal dan CBAM untuk penyempurnaan fitur global. Decoder mencerminkan encoder dengan lima tahap upsampling, masing-masing dilengkapi gerbang atensi untuk fokus pada wilayah teks:

\begin{equation}
\alpha = \sigma(W_g^T g + W_x^T x + b)
\end{equation}

\noindent di mana $g \in \mathbb{R}^{H \times W \times C}$ adalah fitur decoder dari skala yang lebih kasar (sinyal gating), $x \in \mathbb{R}^{H \times W \times C}$ adalah fitur koneksi langsung dari encoder, $W_g$ dan $W_x$ adalah matriks proyeksi yang dapat dipelajari, $b$ adalah bias, dan $\sigma$ adalah fungsi aktivasi sigmoid. Koefisien atensi $\alpha \in [0,1]^{H \times W \times C}$ kemudian dikalikan elemen-demi-elemen dengan koneksi langsung $x$ untuk menekan fitur latar belakang yang tidak relevan sambil memperkuat wilayah pembawa teks: $\hat{x} = \alpha \odot x$. Mekanisme ini memungkinkan decoder secara adaptif memilih fitur resolusi tinggi yang relevan dari encoder.

% Tabel Arsitektur Generator
\begin{table*}[t]
\centering
\caption{Arsitektur Generator: U-Net Enhanced dengan Blok Residual dan Gerbang Atensi}
\label{tab:generator_arch}
\begin{tabular}{|l|l|c|c|r|}
\hline
\textbf{Tahap} & \textbf{Lapisan} & \textbf{Bentuk Keluaran} & \textbf{Kernel/Stride} & \textbf{Parameter} \\
\hline
Masukan & Masukan & $128 \times 1024 \times 1$ & - & - \\
\hline
\multirow{2}{*}{Encoder-0} 
& ResBlock-1 & $128 \times 1024 \times 64$ & $3 \times 3 / 1$ & 4.7K \\
& MaxPool-1 & $64 \times 512 \times 64$ & $2 \times 2 / 2$ & - \\
\hline
\multirow{2}{*}{Encoder-1}
& ResBlock-2 & $64 \times 512 \times 128$ & $3 \times 3 / 1$ & 147.6K \\
& MaxPool-2 & $32 \times 256 \times 128$ & $2 \times 2 / 2$ & - \\
\hline
\multirow{2}{*}{Encoder-2}
& ResBlock-3 & $32 \times 256 \times 256$ & $3 \times 3 / 1$ & 590.1K \\
& MaxPool-3 & $16 \times 128 \times 256$ & $2 \times 2 / 2$ & - \\
\hline
\multirow{2}{*}{Encoder-3}
& ResBlock-4 & $16 \times 128 \times 512$ & $3 \times 3 / 1$ & 2.4M \\
& MaxPool-4 & $8 \times 64 \times 512$ & $2 \times 2 / 2$ & - \\
\hline
\multirow{2}{*}{Encoder-4}
& ResBlock-5 & $8 \times 64 \times 512$ & $3 \times 3 / 1$ & 2.4M \\
& MaxPool-5 & $4 \times 32 \times 512$ & $2 \times 2 / 2$ & - \\
\hline
Bottleneck & MSFP + RDB + CBAM & $4 \times 32 \times 512$ & - & 5.1M \\
\hline
\multirow{3}{*}{Decoder-4}
& UpSample-5 & $8 \times 64 \times 512$ & - & - \\
& AttentionGate-5 & $8 \times 64 \times 512$ & - & 262K \\
& RDB + CBAM & $8 \times 64 \times 512$ & $3 \times 3 / 1$ & 3.5M \\
\hline
\multirow{3}{*}{Decoder-3}
& UpSample-4 & $16 \times 128 \times 512$ & - & - \\
& AttentionGate-4 & $16 \times 128 \times 512$ & - & 262K \\
& RDB + CBAM & $16 \times 128 \times 512$ & $3 \times 3 / 1$ & 3.5M \\
\hline
\multirow{3}{*}{Decoder-2}
& UpSample-3 & $32 \times 256 \times 256$ & - & - \\
& AttentionGate-3 & $32 \times 256 \times 256$ & - & 66K \\
& RDB + CBAM & $32 \times 256 \times 256$ & $3 \times 3 / 1$ & 885K \\
\hline
\multirow{3}{*}{Decoder-1}
& UpSample-2 & $64 \times 512 \times 128$ & - & - \\
& AttentionGate-2 & $64 \times 512 \times 128$ & - & 16.5K \\
& RDB + CBAM & $64 \times 512 \times 128$ & $3 \times 3 / 1$ & 221K \\
\hline
\multirow{3}{*}{Decoder-0}
& UpSample-1 & $128 \times 1024 \times 64$ & - & - \\
& AttentionGate-1 & $128 \times 1024 \times 64$ & - & 4.1K \\
& RDB + CBAM & $128 \times 1024 \times 64$ & $3 \times 3 / 1$ & 55K \\
\hline
Keluaran & Conv2D + Tanh & $128 \times 1024 \times 1$ & $1 \times 1 / 1$ & 65 \\
\hline
\multicolumn{4}{|r|}{\textbf{Total Parameter:}} & \textbf{21.8M} \\
\hline
\end{tabular}
\end{table*}

% Diagram Arsitektur Generator U-Net Enhanced
\begin{figure}[!t]
\centering
\includegraphics[width=0.5\textwidth,keepaspectratio]{../drawio/arsitektur_generator.png}
\caption{Arsitektur Generator U-Net Enhanced: encoder 5-level (RDB+CBAM), bottleneck MSFP (512 kanal), decoder dengan attention gates. Skip connections mempertahankan detail spasial.}
\label{fig:unet_architecture}
\end{figure}

Statistik Arsitektur:
\begin{itemize}
    \item Total parameter: 21.8M (semua terlatih)
    \item Memory footprint: $\sim$87 MB (float32)
    \item Kedalaman: 5 tahap encoder + bottleneck + 5 tahap decoder
    \item Masukan/Keluaran: $1024 \times 128 \times 1$ (L$\times$T$\times$C)
    \item Aktivasi keluaran: Tanh (rentang [-1, 1])
\end{itemize}

Gambar~\ref{fig:unet_architecture} memvisualisasikan arsitektur simetris U-Net Enhanced, di mana setiap level encoder memiliki pasangan decoder yang berkomunikasi melalui skip connections dengan attention gates. Berbeda dari U-Net konvensional~\cite{ronneberger2015unet}, integrasi RDB+CBAM pada setiap level meningkatkan kapasitas representasi sebesar $\sim$40\% tanpa degradasi gradien, sementara bottleneck MSFP menangkap konteks global multi-skala sebelum proses rekonstruksi bertahap pada decoder.

\subsubsection{Rasional Desain}

Residual Dense Blocks memberikan keunggulan ganda: (1) propagasi gradien yang lebih baik melalui koneksi padat multi-jalur, memungkinkan pelatihan arsitektur dalam tanpa degradasi gradien, dan (2) penggunaan kembali fitur hierarkis di mana setiap lapisan dapat mengakses fitur dari semua lapisan sebelumnya, meningkatkan kapasitas representasi tanpa meningkatkan parameter secara substansial. CBAM memberikan selektivitas adaptif pada dimensi kanal dan spasial, memungkinkan jaringan secara otomatis memfokuskan pada kanal fitur yang informatif ("what") dan wilayah spasial yang relevan ("where"), efektif untuk menekan noise latar belakang sambil memperkuat fitur goresan teks. Multi-Scale Feature Pyramid~\cite{lin2017fpn} pada bottleneck menangkap konteks multi-resolusi melalui konvolusi terdilasi~\cite{chen2017deeplab} dengan skala berbeda, penting untuk menangani degradasi dengan ukuran bervariasi (noise titik vs noda besar). Attention Gates~\cite{oktay2018attention} memberikan selektivitas spasial pada koneksi langsung, memungkinkan decoder fokus pada wilayah pembawa teks sambil menekan amplifikasi derau—krusial untuk melestarikan goresan tipis dan tanda diakritik paleografi. Kombinasi teknik-teknik ini menghasilkan arsitektur yang mampu mempertahankan detail halus sambil menghilangkan degradasi kompleks.

\subsection{Diskriminator Dual-Modal Enhanced}

Arsitektur diskriminator dual-modal mengevaluasi citra dari perspektif spasial (visual) dan sekuensial (teks) secara simultan. Arsitektur ini dirancang untuk memberikan sinyal adversarial yang lebih kaya dengan mempertimbangkan baik kualitas visual maupun koherensi teks.

\subsubsection{Spesifikasi Arsitektur}

Tabel~\ref{tab:discriminator_arch} merinci arsitektur lengkap diskriminator. Diskriminator menerima citra ($128 \times 1024 \times 1$) dan sekuens teks (panjang maksimal 128 karakter), kemudian memprosesnya melalui dua jalur paralel seperti ditunjukkan pada Gambar~\ref{fig:discriminator_architecture}.

% Tabel Arsitektur Diskriminator
\begin{table}[!t]
\centering
\caption{Arsitektur Diskriminator Dual-Modal Enhanced: spesifikasi detail tiga cabang (Citra CNN, Teks BiLSTM, Fusi Cross-Modal). Cabang citra memproses fitur visual melalui 4 DownBlocks dengan spatial attention, cabang teks memproses sekuens karakter melalui embedding + BiLSTM + self-attention. Cross-modal fusion menghasilkan klasifikasi real/fake. Total: 17.4M parameter.}
\label{tab:discriminator_arch}
\scriptsize
\begin{tabular}{|l|l|c|r|}
\hline
\textbf{Cabang} & \textbf{Lapisan} & \textbf{Keluaran} & \textbf{Param.} \\
\hline
\multirow{9}{*}{\rotatebox{90}{\tiny CNN}}
& Input & $128{\times}1024{\times}1$ & - \\
& Conv2D+BN & $128{\times}1024{\times}64$ & 0.6K \\
& DownBlk-1 & $64{\times}512{\times}64$ & 37K \\
& DownBlk-2 & $32{\times}256{\times}128$ & 148K \\
& DownBlk-3 & $16{\times}128{\times}256$ & 590K \\
& DownBlk-4 & $8{\times}64{\times}512$ & 2.4M \\
& SpatialAttn & $8{\times}64{\times}512$ & 4.7K \\
& ResBlk-512 & $8{\times}64{\times}512$ & 2.4M \\
& AvgPool & $512$ & - \\
\hline
\multirow{5}{*}{\rotatebox{90}{\tiny LSTM}}
& Input & $128$ (seq.) & - \\
& Embed & $128{\times}128$ & 12.8K \\
& BiLSTM & $128{\times}512$ & 1.3M \\
& SelfAttn & $128{\times}512$ & 262K \\
& AvgPool & $512$ & - \\
\hline
\multirow{6}{*}{\rotatebox{90}{\tiny Fusi}}
& Proj-CNN & $128$ & 65.7K \\
& Proj-LSTM & $128$ & 65.7K \\
& CrossAttn & $256$ & 33K \\
& Dense-1+Drop & $512$ & 131K \\
& Dense-2+Drop & $256$ & 131K \\
& Sigmoid & $1$ & 257 \\
\hline
\multicolumn{3}{|r|}{\textbf{Total:}} & \textbf{17.4M} \\
\hline
\end{tabular}
\end{table}

% Diagram Arsitektur Diskriminator Dual-Modal
\begin{figure}[!t]
\centering
\includegraphics[width=0.5\textwidth,keepaspectratio]{../drawio/arsitektur-dualModal.png}
\caption{Arsitektur Diskriminator Dual-Modal Enhanced dengan bilateral cross-modal attention: cabang CNN (512-D) dengan spatial attention, cabang BiLSTM (512-D) dengan self-attention. Fusion bilateral (128-D) untuk validasi holistik.}
\label{fig:discriminator_architecture}
\end{figure}

Gambar~\ref{fig:discriminator_architecture} memvisualisasikan arsitektur dual-modal yang membedakan pendekatan ini dari diskriminator konvensional. Berbeda dari PatchGAN~\cite{isola2017image} yang hanya mengevaluasi patch visual, arsitektur ini menggabungkan analisis spasial (CNN) dan sekuensial (BiLSTM) untuk validasi holistik. Karakteristik utama adalah bilateral cross-modal attention (Img$\leftrightarrow$Text) yang memungkinkan interaksi dua arah: cabang citra memvalidasi koherensi teks dari HTR recognizer, sementara cabang teks memverifikasi kualitas visual goresan. Mekanisme bilateral ini menghasilkan gradient feedback yang lebih kaya kepada generator dengan reduksi parameter 52\% (17.4M vs 137M baseline), mendorong perbaikan simultan pada kualitas visual dan konsistensi tekstual.

Cabang Citra (CNN dengan Atensi Spasial): Cabang ini mengekstrak fitur visual spasial melalui arsitektur berbasis ResNet. Masukan citra $I \in \mathbb{R}^{128 \times 1024 \times 1}$ diproses melalui lapisan konvolusional awal (64 kanal) dengan batch normalization dan aktivasi LeakyReLU ($\alpha=0.2$). Empat blok downsampling residual secara progresif mengurangi resolusi spasial sambil meningkatkan kedalaman kanal:
\begin{equation}
\begin{split}
128 \times 1024 \times 64 &\rightarrow 64 \times 512 \times 64 \rightarrow 32 \times 256 \times 128 \\
&\rightarrow 16 \times 128 \times 256 \rightarrow 8 \times 64 \times 512
\end{split}
\end{equation}
Mekanisme atensi spasial dengan kernel $3 \times 3$ diterapkan pada fitur 512 kanal untuk fokus pada wilayah pembawa teks:

\begin{equation}
\alpha_{\text{spatial}} = \sigma(\text{Conv}_{3 \times 3}(F_{\text{img}}))
\end{equation}

\noindent di mana $\sigma$ adalah fungsi sigmoid dan $F_{\text{img}}$ adalah peta fitur dengan dimensi $8 \times 64 \times 512$. Setelah atensi spasial, blok residual tambahan dengan 512 kanal (ResBlock-512) digunakan untuk pemrosesan fitur lebih lanjut tanpa mengubah dimensi spasial, kemudian global average pooling mengagregasi fitur spasial $8 \times 64 \times 512$ menjadi vektor representasi $F_{\text{img}}^{\text{global}} \in \mathbb{R}^{512}$.

Cabang Teks (BiLSTM dengan Atensi-Mandiri): Cabang ini menangkap koherensi sekuensial teks. Sekuens karakter masukan $S \in \mathbb{N}^{L}$ (panjang $L \leq 128$) pertama kali diproyeksikan ke ruang embedding $E \in \mathbb{R}^{L \times 128}$. BiLSTM memproses sekuens dalam kedua arah untuk menangkap dependensi kontekstual:

\begin{equation}
h_t = [\overrightarrow{h_t}; \overleftarrow{h_t}] = [\text{LSTM}_{\text{fwd}}(E_t); \text{LSTM}_{\text{bwd}}(E_t)]
\end{equation}

\noindent menghasilkan hidden states $H \in \mathbb{R}^{L \times 512}$. Mekanisme atensi-mandiri multi-head scaled dot-product diterapkan:

\begin{equation}
\text{Attention}(Q, K, V) = \text{softmax}\left(\frac{QK^T}{\sqrt{d_k}}\right)V
\end{equation}

\noindent di mana $Q, K, V$ adalah proyeksi linear dari $H$. Global average pooling menghasilkan vektor teks $F_{\text{text}}^{\text{global}} \in \mathbb{R}^{512}$.

Fusi Cross-Modal: Kedua vektor fitur $F_{\text{img}}^{\text{global}} \in \mathbb{R}^{512}$ dan $F_{\text{text}}^{\text{global}} \in \mathbb{R}^{512}$ diproses melalui mekanisme atensi cross-modal bilateral untuk memungkinkan interaksi dua arah antar modalitas. Pendekatan bilateral diperlukan karena: (1) citra perlu memvalidasi koherensi teks yang diprediksi pengenal HTR, dan (2) teks perlu memverifikasi kualitas visual goresan yang menghasilkan karakter tersebut. Dengan demikian, kedua modalitas saling memberikan konteks untuk deteksi anomali yang lebih robust. Setiap modalitas memproyeksikan fiturnya ke ruang bersama berdimensi rendah $d=128$ (dikurangi dari 256 untuk stabilitas) melalui enam matriks proyeksi terpisah:

Atensi Citra$\rightarrow$Teks: Citra sebagai query, teks sebagai key/value:
\begin{align}
Q_{\text{img}} &= W_{\text{img}}^Q F_{\text{img}}^{\text{global}} \nonumber\\
K_{\text{text}} &= W_{\text{text}}^K F_{\text{text}}^{\text{global}} \nonumber\\
V_{\text{text}} &= W_{\text{text}}^V F_{\text{text}}^{\text{global}} \label{eq:cross_img_proj}\\
F_{\text{img}}^{\text{att}} &= \text{softmax}\left(\frac{Q_{\text{img}} K_{\text{text}}^T}{\sqrt{d}}\right) V_{\text{text}} \label{eq:cross_img_att}
\end{align}

\noindent di mana $W_{\text{img}}^Q, W_{\text{text}}^K, W_{\text{text}}^V \in \mathbb{R}^{512 \times 128}$ adalah matriks proyeksi yang dipelajari.

Atensi Teks$\rightarrow$Citra: Teks sebagai query, citra sebagai key/value:
\begin{align}
Q_{\text{text}} &= W_{\text{text}}^Q F_{\text{text}}^{\text{global}} \nonumber\\
K_{\text{img}} &= W_{\text{img}}^K F_{\text{img}}^{\text{global}} \nonumber\\
V_{\text{img}} &= W_{\text{img}}^V F_{\text{img}}^{\text{global}} \label{eq:cross_text_proj}\\
F_{\text{text}}^{\text{att}} &= \text{softmax}\left(\frac{Q_{\text{text}} K_{\text{img}}^T}{\sqrt{d}}\right) V_{\text{img}} \label{eq:cross_text_att}
\end{align}

\noindent di mana $W_{\text{text}}^Q, W_{\text{img}}^K, W_{\text{img}}^V \in \mathbb{R}^{512 \times 128}$ adalah matriks proyeksi simetris

Kedua fitur yang telah mengalami atensi bilateral ($F_{\text{img}}^{\text{att}}, F_{\text{text}}^{\text{att}} \in \mathbb{R}^{128}$) digabungkan menjadi vektor fusi: $F_{\text{fused}} = [F_{\text{img}}^{\text{att}}; F_{\text{text}}^{\text{att}}] \in \mathbb{R}^{256}$. Kepala klasifikasi terdiri dari dua lapisan padat (512 $\rightarrow$ 256 neuron) dengan dropout dan aktivasi LeakyReLU, diikuti lapisan keluaran dengan aktivasi sigmoid yang menghasilkan skor validitas $D(I, S) \in [0, 1]$ untuk membedakan pasangan citra-teks asli dari pasangan palsu yang dihasilkan generator.

Konfigurasi hyperparameter: Kernel atensi spasial menggunakan ukuran $3 \times 3$ (dikurangi dari $7 \times 7$ untuk mengurangi artefak titik putih pada goresan); dimensi proyeksi cross-modal ditetapkan 128 (dikurangi dari 256 untuk meningkatkan stabilitas pelatihan); unit LSTM per arah sebanyak 256 menghasilkan total keluaran BiLSTM 512 dimensi. Momentum batch normalization 0.9 (ditingkatkan dari 0.8 untuk aliran gradien lebih halus); tingkat dropout 0.1 pada kepala klasifikasi (dikurangi dari 0.3 untuk mempertahankan informasi goresan tipis) dan 0.2 pada embedding teks; fungsi aktivasi LeakyReLU ($\alpha=0.2$) untuk CNN dan sigmoid untuk aktivasi rekuren LSTM (memastikan kompatibilitas dengan implementasi LSTM standar tanpa cuDNN).

\subsubsection{Motivasi dan Rasional Desain}

Pendekatan dual-modal: Dokumen teks memiliki sifat ganda: (1) penampilan visual (kualitas goresan, tingkat derau), (2) struktur sekuensial (aliran karakter, jarak antar kata). Jalur CNN menangkap kualitas spasial, sementara BiLSTM menangkap koherensi sekuensial. Fusi cross-modal bilateral memungkinkan diskriminator menilai pasangan citra-teks secara holistik dengan interaksi dua arah, memberikan sinyal adversarial yang lebih informatif kepada generator untuk perbaikan simultan pada kualitas visual dan konsistensi tekstual.

Efisiensi parameter: Diskriminator yang diusulkan mencapai efisiensi parameter signifikan (17.4M, reduksi 52\% dari baseline 137M) melalui arsitektur yang disederhanakan tanpa mengorbankan kemampuan diskriminatif: (1) proyeksi dimensi rendah (128) untuk fusi cross-modal mengurangi kompleksitas komputasi sambil mempertahankan kapasitas representasi; (2) kernel atensi spasial lebih kecil (3$\times$3) fokus pada detail lokal tanpa artefak; (3) blok residual yang dioptimalkan menyeimbangkan kedalaman arsitektur dengan stabilitas pelatihan. Momentum batch normalization lebih tinggi (0.9) menghasilkan statistik lebih stabil dan aliran gradien lebih halus; tingkat dropout lebih rendah (0.1) mempertahankan informasi goresan tipis yang krusial untuk diskriminasi dokumen terdegradasi.

Kemampuan: Dengan kombinasi analisis spasial CNN dan pemrosesan sekuens BiLSTM, diskriminator dapat mendeteksi:
\begin{itemize}
    \item Goresan yang rusak atau terputus yang mengganggu alur teks (melalui deteksi anomali sekuens BiLSTM)
    \item Pola derau, artefak, dan tembus-balik (melalui fitur spasial CNN)
    \item Lebar goresan yang tidak konsisten (melalui mekanisme atensi spasial)
    \item Celah atau koneksi yang tidak wajar antar karakter (melalui analisis fusi cross-modal)
\end{itemize}

Dengan menggabungkan penilaian spasial dan sekuensial, diskriminator memberikan umpan balik yang lebih kaya kepada generator, mendorong produksi citra yang bersih secara visual dan koheren secara tekstual.

\subsection{Integrasi \textit{Recognizer} HTR Beku}

Untuk secara eksplisit mengoptimalkan kinerja HTR, penelitian ini mengintegrasikan \textit{recognizer} teks praterlatih dengan strategi pembekuan bobot, di mana \textit{recognizer} dibekukan (\texttt{model.trainable=False}) dan digunakan sebagai ekstraktor fitur.

\subsubsection{Arsitektur dan Konfigurasi Pelatihan}

Arsitektur hibrida CNN-Transformer yang digunakan terdiri dari:
\begin{itemize}
    \item Backbone CNN: 4 blok konvolusi dengan ekstraksi fitur progresif, BatchNormalization, dan aktivasi GELU
    \item Proyeksi Urutan: Lapisan padat + LayerNormalization + Dropout
    \item Encoder Transformer: 6 lapis, 8 kepala atensi, positional encoding
    \item Keluaran: Dekoder CTC untuk prediksi urutan
\end{itemize}

Pengenal dilatih pada 4.739 citra baris tulisan tangan berbahasa Belanda dari dokumen ANRI era kolonial abad ke-16 hingga ke-18 dengan pembagian 80/20 (pelatihan $n$=3.791, validasi $n$=948), mencapai kinerja CER 33.72\% pada set validasi. Model kemudian dibekukan (\texttt{trainable=False}) dan diintegrasikan ke dalam alur pelatihan GAN.

Meskipun CER 33.72\% tergolong moderat, kinerja ini sebanding dengan \textit{state-of-the-art} pada \textit{dataset} historis paleografi (READ 2016: 25--35\%~\cite{sanchez2016icfhr}, ICDAR 2017: 28--42\%~\cite{sanchez2017icdar}) dan mencerminkan kompleksitas intrinsik tulisan tangan abad ke-16 hingga ke-18. Komponen HTR berfungsi sebagai \textit{feature extractor} untuk \textit{guidance} sadar-teks, bukan \textit{optimizer} visual.

Studi ablasi (Bagian~V-B, Tabel~\ref{table_ablation_loss}) memvalidasi peran HTR: pada pelatihan 15 \textit{epoch}, konfigurasi dengan HTR (Eksperimen~4) mencapai CER 29.66\%, sementara \textit{baseline} tanpa HTR (Eksperimen~1) tidak dapat mengukur CER karena \textit{recognizer} tidak terintegrasi. Masukan terdegradasi menunjukkan CER 83.4\%. Model produksi dengan \textit{extended training} (50 \textit{epoch}) dan integrasi HTR mencapai CER 27.11\% pada \textit{set} validasi ($n$=710) dan CER 34.9\% pada \textit{set} uji ($n$=712) dengan PSNR 30.74~dB. \textit{Catatan}: Perbedaan CER 7.79 poin persentase antara validasi dan uji bukan indikasi \textit{overfitting}, melainkan mencerminkan perbedaan kompleksitas intrinsik teks antar \textit{split}\textemdash\textit{recognizer} beku pada citra bersih \textit{ground truth} juga menunjukkan perbedaan serupa (26.57\% validasi vs 34.1\% uji, $\Delta$=7.53\%).

Peran utama HTR adalah memastikan \textit{text preservation} dan \textit{readability} (\textit{CER improvement} 48.5--56.3 poin absolut dari masukan terdegradasi pada set uji dan validasi), sementara peningkatan PSNR berasal dari kombinasi \textit{HTR guidance}, \textit{adversarial training}, \textit{perceptual loss}, dan \textit{extended training}. Pendekatan ini membuktikan \textit{robustness} metode terhadap komponen HTR realistis tanpa memerlukan \textit{recognizer} sempurna.

\subsubsection{Keuntungan Strategi Pembekuan}

Pendekatan pembekuan \textit{recognizer} memberikan manfaat implementasi:
\begin{enumerate}
    \item Stabilitas pelatihan: Tidak ada optimisasi bersama atau gradien yang berkonflik antara \textit{generator} dan pengenal
    \item Efisiensi komputasi: Tidak ada \textit{backpropagation} melalui pengenal, mengurangi beban komputasi
    \item Konsistensi evaluasi: Bobot tetap menyediakan metrik CER/WER yang konsisten sepanjang pelatihan
    \item Supervisi dengan label: Menggunakan teks ground truth untuk supervisi kehilangan CTC
\end{enumerate}

\subsubsection{Ekstraksi Fitur dan Perhitungan Kehilangan}

Penelitian ini mengekstrak fitur pengenalan perantara dari \textit{recognizer} beku:

\begin{equation}
\mathcal{F}_{\text{rec}}(I) = R_{\text{encoder}}(I)
\end{equation}

di mana $R_{\text{encoder}}$ mewakili bagian encoder yang dibekukan dari pengenal. Secara spesifik, digunakan keluaran dari lapisan proyeksi sebelum encoder transformer, menghasilkan peta fitur dengan dimensi (batch, 128, 512). Lapisan ini dipilih karena mengandung fitur tingkat urutan sebelum pengodean kontekstual, memberikan gradien yang stabil untuk generator. Fitur-fitur ini menangkap karakteristik yang relevan dengan teks seperti pola goresan, bentuk karakter, dan pengaturan spasial.

Kehilangan fitur pengenalan antara citra yang dihasilkan dan citra ground truth dihitung menggunakan Mean Squared Error (MSE):

\begin{equation}
\begin{split}
\mathcal{L}_{\text{rec-feat}} &= \text{MSE}(\mathcal{F}_{\text{rec}}(I_{\text{gen}}), \mathcal{F}_{\text{rec}}(I_{\text{gt}})) \\
&= \|\mathcal{F}_{\text{rec}}(I_{\text{gen}}) - \mathcal{F}_{\text{rec}}(I_{\text{gt}})\|_2^2
\end{split}
\end{equation}

MSE dipilih karena lebih sensitif terhadap deviasi fitur besar, memberikan gradien yang lebih kuat untuk perbaikan keselarasan fitur. Ini mendorong generator untuk menghasilkan citra yang fitur teksnya cocok dengan fitur teks dari ground truth yang bersih, secara implisit meningkatkan keterbacaan HTR.

\subsection{Fungsi Loss Multikomponen}

Generator dilatih dengan kombinasi lima komponen kehilangan yang dikalibrasi melalui pencarian kisi empiris dan divalidasi pada implementasi aktual:

\begin{equation}
\begin{split}
\mathcal{L}_{\text{total}} = &\lambda_{\text{adv}}\mathcal{L}_{\text{adv}} + \lambda_{\text{pixel}}\mathcal{L}_{\text{pixel}} + \lambda_{\text{perc}}\mathcal{L}_{\text{perc}} \\
&+ \lambda_{\text{ctc}}\mathcal{L}_{\text{ctc}} + \lambda_{\text{rec-feat}}\mathcal{L}_{\text{rec-feat}}
\end{split}
\end{equation}

Konfigurasi bobot kehilangan untuk setiap komponen ditunjukkan dalam Tabel~\ref{table_optimal_weights}.

\subsubsection{Kehilangan Adversarial}

Kehilangan adversarial standar GAN mendorong pembuatan citra yang realistis:

\begin{equation}
\mathcal{L}_{\text{adv}} = -\log D(I_{\text{deg}}, I_{\text{gen}})
\end{equation}

Discriminator dilatih untuk memaksimalkan:

\begin{equation}
\begin{split}
\mathcal{L}_D = -[&\log D(I_{\text{deg}}, I_{\text{gt}}) \\
&+ \log(1 - D(I_{\text{deg}}, I_{\text{gen}}))]
\end{split}
\end{equation}

\subsubsection{Kehilangan Rekonstruksi Tingkat Piksel}

Kehilangan L1 memastikan kesamaan piksel-demi-piksel dengan ground truth:

\begin{equation}
\mathcal{L}_{\text{pixel}} = \|I_{\text{gen}} - I_{\text{gt}}\|_1
\end{equation}

Norma L1 dipilih karena menghasilkan hasil yang lebih tajam dibandingkan norma L2 dan lebih robust terhadap pencilan.

\subsubsection{Kehilangan Perseptual}

Kehilangan perseptual~\cite{johnson2016perceptual} membandingkan fitur tingkat tinggi dari jaringan VGG praterlatih:

\begin{equation}
\mathcal{L}_{\text{perc}} = \sum_{l \in L} \|\phi_l(I_{\text{gen}}) - \phi_l(I_{\text{gt}})\|_2
\end{equation}

di mana $\phi_l$ menunjukkan fitur dari lapisan $l$ dari VGG-19 yang telah dilatih sebelumnya pada \textit{ImageNet}. Fitur diekstrak dari beberapa lapisan konvolusional untuk perbandingan perseptual multiskala. Pendekatan ini memastikan pelestarian karakteristik visual dari detail tingkat rendah (tepi, tekstur) hingga fitur semantik tingkat tinggi (pola goresan, bentuk karakter).

\subsubsection{Kehilangan Fitur Pengenalan}

Seperti yang dijelaskan di Bagian III-D, komponen ini mendorong karakteristik yang ramah HTR menggunakan \textit{Mean Squared Error}:

\begin{equation}
\mathcal{L}_{\text{rec-feat}} = \|\mathcal{F}_{\text{rec}}(I_{\text{gen}}) - \mathcal{F}_{\text{rec}}(I_{\text{gt}})\|_2^2
\end{equation}

% [PLACEHOLDER: Tabel 2 - Bobot kehilangan optimal dari penyetelan empiris]
\begin{table}[!t]
\renewcommand{\arraystretch}{1.3}
\caption{Konfigurasi Bobot Kehilangan: bobot optimal hasil pencarian kisi empiris untuk menyeimbangkan kualitas visual dengan keterbacaan teks pada dataset validasi}
\label{table_optimal_weights}
\centering
\begin{tabular}{|l|c|c|}
\hline
Komponen Kehilangan & Bobot ($\lambda$) & Justifikasi \\
\hline
Adversarial ($\mathcal{L}_{\text{adv}}$) & 3.0 & Realisme yang ditingkatkan \\
Piksel/L1 ($\mathcal{L}_{\text{pixel}}$) & 50.0 & Pelestarian seimbang \\
Perseptual ($\mathcal{L}_{\text{perc}}$) & 1.0 & Pelestarian topologi \\
Fitur Pengenalan ($\mathcal{L}_{\text{rec-feat}}$) & 8.0 & Panduan teks kuat \\
CTC Loss ($\mathcal{L}_{\text{ctc}}$) & 0.15 & HTR monitoring (clipped) \\
\hline
\end{tabular}
\end{table}

Pemilihan bobot kehilangan (Tabel~\ref{table_optimal_weights}) dilakukan melalui penyetelan empiris berdasarkan analisis trade-off multi-objective antara kualitas visual, pelestarian teks, dan keterbacaan HTR. Karakteristik dataset menunjukkan dominasi goresan tipis yang memerlukan keseimbangan khusus antara rekonstruksi piksel dan panduan berbasis teks.

Bobot rekonstruksi piksel yang tinggi memberikan sinyal rekonstruksi kuat dengan tetap memberikan fleksibilitas untuk peningkatan adversarial, sementara bobot fitur pengenalan yang moderat memberikan panduan sadar-teks tanpa mendominasi kehilangan visual. Mekanisme penyeimbangan adaptif digunakan untuk penyesuaian bobot dinamis selama pelatihan, mencegah dominasi dari komponen kehilangan tunggal dan memastikan konvergensi stabil.

\subsection{Strategi Pelatihan}

\subsubsection{Curriculum Learning dan Protokol Optimisasi}

Pelatihan menggunakan pendekatan curriculum learning tiga fase untuk stabilitas optimal dan konvergensi bertahap:

\begin{enumerate}
    \item Fase Warmup Visual (epoch 1--10): Pelatihan eksklusif pada kualitas visual tanpa supervisi teks, memungkinkan generator membangun kemampuan rekonstruksi dasar. Fase ini mencegah mode collapse dini yang dapat terjadi jika kehilangan CTC diperkenalkan terlalu awal.
    
    \item Fase Peningkatan Bertahap CTC (epoch 11--30): Peningkatan bertahap bobot CTC dari nol ke bobot penuh secara linear (Tabel~\ref{table_optimal_weights}). Integrasi halus ini memungkinkan generator beradaptasi dengan gradien HTR tanpa mengganggu kestabilan pelatihan adversarial yang sudah terbentuk.
    
    \item Fase Pelatihan Penuh (epoch 31--50): Pelatihan dengan semua komponen kehilangan aktif pada bobot optimal, disertai penyeimbangan adaptif untuk menjaga keseimbangan dinamis antara kehilangan visual dan tekstual.
\end{enumerate}

Prosedur pelatihan adversarial bergantian dengan supervisi dual-modal diterapkan:
\begin{enumerate}
    \item Langkah discriminator: Perbarui $D$ untuk memaksimalkan kemampuan membedakan pasangan nyata $(I_{\text{gt}}, \text{teks}_{\text{gt}})$ dari pasangan palsu $(I_{\text{gen}}, \text{teks}_{\text{pred}})$ dengan supervisi simultan pada domain visual dan tekstual.
    \item Langkah generator: Perbarui $G$ untuk meminimalkan $\mathcal{L}_{\text{total}}$ (kombinasi lima komponen kehilangan), menghasilkan citra yang secara simultan menipu discriminator dan mengoptimalkan keterbacaan HTR.
\end{enumerate}
Pembaruan dilakukan dengan rasio seimbang 1:1, memastikan stabilitas pelatihan adversarial.

Konfigurasi pelatihan didasarkan pada penyetelan empiris untuk stabilitas GAN dan konvergensi optimal:

\begin{itemize}
    \item Optimisasi: Generator menggunakan algoritma Adam dengan parameter momentum yang disesuaikan untuk stabilitas GAN, sementara discriminator menggunakan SGD dengan momentum untuk mencegah mode collapse dan memastikan konvergensi stabil. Kedua model menerapkan gradient clipping untuk mencegah ledakan gradien.
    
    \item Laju pembelajaran: Kedua model menggunakan laju pembelajaran sama dengan penjadwalan cosine annealing untuk konvergensi halus tanpa penurunan tajam.
    
    \item Ukuran batch dan dimensi: Ukuran batch kecil dengan dimensi citra lebar yang mempertahankan detail goresan halus pada teks baris tunggal paleografi.
    
    \item Penghentian dini: Diterapkan dengan mekanisme patience yang mempertimbangkan fase curriculum, memulihkan bobot terbaik berdasarkan metrik gabungan PSNR dan CER dengan pembobotan yang memprioritaskan kualitas visual.
\end{itemize}

\subsubsection{Analisis Stabilitas dan Regularisasi}

Gambar~\ref{fig:convergence_analysis} menunjukkan trajektori konvergensi pelatihan selama 50 epoch, memvalidasi stabilitas dan efektivitas strategi pelatihan yang diusulkan.

\begin{figure*}[!t]
\centering
\includegraphics[width=\textwidth]{../data_dukung/fig_convergence_analysis_combined.pdf}
\caption{Trajektori konvergensi pelatihan selama 50 epoch: (a) Kehilangan generator menunjukkan penurunan stabil tanpa mode collapse, (b) PSNR meningkat mencapai target yang ditetapkan, (c) SSIM mencapai nilai tinggi menunjukkan pelestarian struktur. Model terbaik (ditandai garis vertikal putus-putus dan bintang merah) dipilih berdasarkan kriteria gabungan PSNR-SSIM. Kurva konvergensi halus memvalidasi efektivitas strategi curriculum learning yang diusulkan.}
\label{fig:convergence_analysis}
\end{figure*}

Observasi Konvergensi:
\begin{itemize}
    \item Kehilangan generator: Penurunan monotonik halus pada fase awal, stabil pada fase akhir tanpa osilasi, mengindikasikan pelatihan adversarial yang stabil tanpa mode collapse
    \item Trajektori PSNR: Peningkatan cepat pada Fase Warmup Visual, diikuti peningkatan stabil pada Fase Peningkatan Bertahap CTC, plateau pada fase akhir memicu pertimbangan penghentian dini
    \item Progresi SSIM: Tren naik konsisten mencapai nilai tinggi pada epoch terbaik, menunjukkan pelestarian kesamaan struktural bersama akurasi tingkat piksel
    \item Pemilihan model terbaik: Dipilih berdasarkan optimisasi metrik gabungan PSNR-SSIM, bukan penghentian arbitrer
    \item Keandalan statistik: Interval kepercayaan yang sempit memvalidasi kinerja robust pada validation set
\end{itemize}

Teknik Regularisasi: Untuk mencegah overfitting dan meningkatkan generalisasi model, teknik berikut diterapkan:

\begin{itemize}
    \item Batch normalization di semua lapisan konvolusional untuk stabilitas pelatihan dan mempercepat konvergensi
    \item Label smoothing asimetris satu sisi: hanya diterapkan pada label nyata untuk mengurangi kepercayaan berlebihan discriminator sambil mempertahankan sinyal pembelajaran yang jelas untuk sampel palsu
    \item Koneksi residual di discriminator untuk aliran gradien yang lebih baik dan pelatihan arsitektur dalam yang stabil
    \item Gradient clipping pada kedua model (norma global 1.0) untuk mencegah ledakan gradien, khususnya dari kehilangan CTC yang di-clip pada nilai maksimum 400.0 untuk stabilitas numerik
    \item Dropout pada classification head (0.1) dan embedding layer (0.2) untuk regularisasi tambahan
\end{itemize}

\paragraph{Detail Implementasi Teknis} Model dilatih dengan momentum 0.9 untuk optimizer SGD discriminator, ukuran batch 2 pada dimensi citra $1024 \times 128 \times 1$ piksel, dan presisi numerik Pure FP32 (tanpa mixed precision) untuk stabilitas CTC loss. Early stopping menggunakan patience 25 epoch dengan min delta 0.05 pada metrik gabungan PSNR-CER.

\paragraph{Catatan Penghentian Dini} Model terbaik dipilih berdasarkan metrik gabungan PSNR-CER pada epoch yang menghasilkan kinerja optimal. Pelatihan dilanjutkan hingga akhir untuk memverifikasi tidak ada peningkatan signifikan. Penghentian dini tidak terpicu karena fluktuasi kecil dalam metrik validasi masih berada dalam ambang batas yang ditentukan, mengonfirmasi model telah mencapai konvergensi stabil.

% ============================================================ 
% V. PENGATURAN EKSPERIMENTAL
% ============================================================ 
\section{Pengaturan Eksperimental}

Bagian ini menjelaskan dataset, metode baseline, metrik evaluasi, detail implementasi, dan prosedur optimisasi hyperparameter yang digunakan dalam eksperimen.

\subsection{Dataset}

Metode yang diusulkan dievaluasi pada \textit{dataset} semi-sintetis realistis yang dikonstruksi dari dokumen historis autentik dan validasi kualitatif pada naskah nyata yang mengalami degradasi alami.

\subsubsection{Dataset Semi-Sintetis Realistis untuk Evaluasi Kuantitatif}

\textbf{Rasional Metodologi:} Evaluasi kuantitatif restorasi dokumen memerlukan pasangan citra terdegradasi-bersih dengan keselarasan tingkat piksel yang sempurna untuk mengukur metrik PSNR, SSIM, dan CER secara akurat. Dokumen historis nyata yang mengalami degradasi alami tidak memiliki \textit{ground truth} bersih karena tidak dapat dipulihkan ke kondisi asli. Oleh karena itu, \textit{dataset} semi-sintetis realistis dikonstruksi dengan menggabungkan tiga komponen autentik~\cite{souibgui2021enhance,zhao2019document}:

\textbf{Komponen 1 - Tulisan Tangan Paleografi Autentik (REAL):}

Dokumen arsip ANRI autentik abad ke-16 hingga ke-18 yang relatif bersih (minimal degradasi visual) digunakan sebagai sumber tulisan tangan. Dokumen dianotasi dalam format Transkribus PAGE XML, dibinarisasi menggunakan model DE-GAN praterlatih~\cite{souibgui2020degan}, dan baris teks individual diekstrak menggunakan koordinat \textit{baseline} dengan \textit{padding} vertikal, dinormalisasi ke dimensi $128 \times 1024$ piksel. Hasil ekstraksi adalah citra baris teks dengan karakteristik paleografi autentik\textemdash variasi ketebalan goresan, ligatur kompleks, gaya penulisan individu penulis historis\textemdash yang tidak dapat direplikasi oleh \textit{rendering} fon sintetis atau model generatif.

\textbf{Komponen 2 - Pola Degradasi Historis Autentik (REAL):}
Dokumen ANRI terpisah yang mengalami degradasi alami berat (penuaan kertas, noda organik, \textit{foxing}, tinta luntur, korosi tinta \textit{iron gall}) digunakan sebagai sumber tekstur latar belakang degradasi. \textit{Patch} latar belakang berukuran $128 \times 1024$ piksel diekstrak dari wilayah non-teks dokumen rusak ini, mempertahankan pola degradasi historis autentik yang terbentuk selama berabad-abad penyimpanan arsip.

\textbf{Komponen 3 - Augmentasi Degradasi Terkontrol (SEMI-SYNTHETIC):}
Untuk meningkatkan variasi dan menciptakan pasangan terdegradasi-bersih dengan keselarasan piksel sempurna, jalur degradasi berlapis dengan enam tahap komputasional diterapkan secara stokastik pada kombinasi teks autentik dan latar belakang degradasi nyata:

1. Overlay Latar Belakang Tekstur:
Citra teks bersih dikomposisikan pada patch tekstur dokumen nyata yang diekstrak dari koleksi dokumen rusak ANRI. Wilayah teks diredup untuk mensimulasikan penyerapan tinta ke dalam kertas yang sudah tua:
\begin{equation}
I_{\text{degradasi}} = I_{\text{bg}} \cdot (1 - M_{\text{teks}}) + (\alpha \cdot I_{\text{bg}}) \cdot M_{\text{teks}}
\end{equation}
di mana $M_{\text{teks}}$ adalah masker teks yang dinormalisasi dan $\alpha$ adalah faktor peredupan.

2. Efek Bleed-Through:
Mensimulasikan tinta yang menembus dari sisi sebalik dokumen dengan membalik citra teks acak secara horizontal, menerapkan Gaussian blur, dan mengurangi dari latar belakang.

3. Noda Organik Berbasis Perlin Noise:
Menghasilkan pola noda yang tampak alami menggunakan Perlin noise 3D~\cite{perlin1985image} dengan parameter yang divariasikan untuk membuat pola noda organik yang halus.

4. Foxing Spots - Bintik Penuaan:
Mensimulasikan perubahan warna terkait usia dengan menambahkan bintik kecil berwarna kecoklatan dengan intensitas acak.

5. Kerusakan Fisik:
Mensimulasikan robekan dan lipatan dokumen melalui garis acak dan patch yang digelapkan.

6. Blur Acak:
Menerapkan Gaussian blur (degradasi halus), median blur (penghilangan derau), atau motion blur (artefak pemindai) dengan probabilitas yang sama.

Efek-efek ini diterapkan secara berurutan dengan kombinasi yang diacak, menghasilkan dokumen terdegradasi yang beragam dan realistis sambil mempertahankan keselarasan ground truth tingkat piksel yang sempurna, menghilangkan kesalahan anotasi manual yang melekat pada dataset dokumen nyata.

\textbf{Statistik Dataset Semi-Sintetis Realistis:} Dataset terdiri dari 4.739 pasangan citra terdegradasi-bersih ($128 \times 1024$ piksel, grayscale, dinormalisasi [0,1]) yang dikonstruksi dari dokumen ANRI autentik. Setiap pasangan menggabungkan: (1) tulisan tangan paleografi historis nyata, (2) pola degradasi dari arsip autentik yang mengalami kerusakan alami, dan (3) augmentasi terkontrol untuk variasi. \textit{Dataset} ini mempertahankan karakteristik autentik dokumen historis sambil menyediakan \textit{ground truth} berpasangan yang diperlukan untuk evaluasi kuantitatif presisi tinggi.

% Gambar 5 - Contoh pasangan degradasi-GT semi-sintetis realistis
\begin{figure}[!t]
\centering
\subfloat[Citra Ground Truth - Tulisan Tangan Paleografi Autentik ANRI]{\label{fig:synthetic_gt}%

\includegraphics[width=3.2in]{../data_dukung/gambar_degGT/gt1.png}}\hfill
\subfloat[Citra Terdegradasi - Kombinasi Degradasi Historis Nyata + Augmentasi]{\label{fig:synthetic_deg}%
\includegraphics[width=3.2in]{../data_dukung/gambar_degGT/deg1.png}}
\caption{Pasangan \textit{dataset} semi-sintetis realistis: (a) \textit{Ground truth} dari dokumen ANRI abad ke-16--18 menunjukkan tulisan tangan paleografi autentik dengan variasi goresan, ligatur, dan gaya penulisan historis yang tidak dapat direplikasi oleh fon sintetis, (b) Citra terdegradasi menggabungkan pola degradasi historis nyata (penuaan kertas, noda organik dari arsip autentik) dengan augmentasi terkontrol (\textit{bleed-through}, \textit{blur}), mempertahankan realisme degradasi dokumen berusia berabad-abad sambil menyediakan keselarasan piksel sempurna untuk evaluasi kuantitatif.}
\label{fig:synthetic_examples}
\end{figure}

\paragraph{Pembagian Dataset (Protokol Akademis)}
Pengambilan data mengikuti protokol akademis pembagian tiga arah untuk evaluasi yang tidak bias:
\begin{itemize}
    \item Pelatihan: 70\% (3.317 citra) - untuk pembelajaran generator dan diskriminator
    \item Validasi: 15\% (710 citra) - untuk penyetelan hyperparameter dan penghentian dini  
    \item Uji: 15\% (712 citra) - set terpisah untuk evaluasi akhir
\end{itemize}

Dataset diacak dengan seed deterministik (seed=42) sebelum pembagian untuk reproduktibilitas.

Stratifikasi Degradasi:
Untuk memastikan generalisasi, jalur degradasi diterapkan dengan randomisasi parameter yang terkontrol, memastikan model belajar pola degradasi umum dan meningkatkan kemampuan generalisasi pada dokumen historis yang tidak terlihat sebelumnya.

\paragraph{Spesifikasi Komputasi} Eksperimen dilakukan pada NVIDIA RTX A4000 (16GB VRAM), AMD Threadripper PRO 3955WX 16-core CPU, 128GB RAM, dengan TensorFlow 2.16.1, CUDA 12.8, dan cuDNN 8.9.7. Training memerlukan sekitar 48 GPU-hours untuk 50 epoch. Inference pada citra baris tunggal ($128\times1024$) memerlukan 50--60ms, dengan throughput batch processing mencapai $\approx$17 citra/detik (batch=2, 30--35ms per citra) menggunakan 6--8GB VRAM.

\subsubsection{Dataset Dokumen Historis Nyata (ANRI)}

Untuk evaluasi generalisasi pada degradasi historis nyata, subset dokumen representatif dari Arsip Nasional Republik Indonesia (ANRI) digunakan sebagai set evaluasi kualitatif:

Pemilihan Dokumen Uji: Dari koleksi arsip ANRI berupa dokumen administratif VOC dan Hindia Belanda (abad ke-16 hingga ke-18, terutama era 1650--1800), sebanyak 15 dokumen halaman penuh dipilih secara manual untuk evaluasi generalisasi. Pemilihan dilakukan berdasarkan kriteria:
\begin{itemize}
    \item Representasi variasi degradasi: Mencakup spektrum tingkat degradasi dari ringan hingga parah untuk menguji robustness model
    \item Diversitas jenis degradasi: Penuaan kertas, noda organik, foxing, tinta luntur, bleed-through, kerusakan fisik
    \item Variasi karakteristik paleografi: Berbagai gaya tulisan tangan, tingkat kontras, dan kompleksitas ligatur
\end{itemize}

\paragraph{Karakteristik Dataset Uji (n=15)}
\textbf{Perbedaan dengan Dataset Semi-Sintetis:} Dataset ANRI nyata ini berbeda fundamental dari \textit{dataset} semi-sintetis untuk evaluasi kuantitatif. Dokumen-dokumen ini mengalami degradasi historis autentik selama berabad-abad penyimpanan arsip, menampilkan pola kerusakan kompleks yang tidak dapat sepenuhnya direplikasi melalui augmentasi komputasional:
\begin{itemize}
    \item Degradasi temporal: Penuaan kertas alami dengan penguningan heterogen dan kerapuhan lokal
    \item Artefak organik: Noda biologis, foxing spots dengan pola distribusi alami, degradasi tinta non-uniform
    \item Kerusakan struktural: Robekan fisik, lipatan permanen, kehilangan material pada tepi dokumen
    \item Variasi tulisan: Tulisan tangan kursif paleografi Belanda dengan ligatur kompleks, singkatan periode, dan inkonsistensi gaya penulisan
\end{itemize}

Dokumen ANRI nyata ini tidak memiliki \textit{ground truth} visual berpasangan (citra bersih) karena dokumen historis autentik tidak dapat dipulihkan ke kondisi asli\textemdash kerusakan fisik dan kimia selama berabad-abad bersifat ireversibel. Oleh karena itu, evaluasi pada \textit{dataset} nyata bersifat kualitatif melalui inspeksi visual dan validasi ahli paleografi, melengkapi evaluasi kuantitatif pada \textit{dataset} semi-sintetis. Pendekatan evaluasi dual ini mengikuti praktik standar dalam riset restorasi dokumen historis~\cite{souibgui2021enhance,souibgui2020degan}: \textit{dataset} semi-sintetis realistis (dikonstruksi dari dokumen autentik dengan augmentasi terkontrol) untuk pengukuran kuantitatif presisi, dan \textit{dataset} nyata dengan degradasi alami untuk validasi generalisasi kualitatif pada kondisi dunia nyata.

\subsection{Metode Baseline dan Strategi Evaluasi}

Evaluasi metode yang diusulkan dilakukan melalui dua pendekatan komplementer: (1) perbandingan kuantitatif dengan metode klasik yang dapat direplikasi pada kondisi eksperimental identik, dan (2) studi ablasi sistematis untuk memvalidasi kontribusi setiap komponen arsitektural.

\subsubsection{Metode Klasik untuk Perbandingan Kuantitatif}

Metode klasik dipilih sebagai baseline kuantitatif karena dapat diimplementasi ulang pada dataset identik tanpa memerlukan pelatihan ulang:

\begin{enumerate}
    \item Tanpa Perbaikan (Degraded Input Baseline): HTR langsung pada citra terdegradasi untuk mengukur dampak degradasi terhadap kinerja HTR
    \item Thresholding Otsu~\cite{otsu1979threshold}: Binerisasi global berbasis histogram untuk pembersihan dokumen
    \item Thresholding Adaptif Sauvola~\cite{sauvola2000adaptive}: Binerisasi adaptif lokal untuk dokumen dengan pencahayaan non-uniform
\end{enumerate}

Metode klasik menggunakan implementasi standar OpenCV dengan parameter yang dioptimalkan melalui pencarian kisi pada subset validasi.

\subsubsection{Metode Deep Learning Terkini (Konteks Penelitian)}

Metode berbasis deep learning yang dibahas dalam tinjauan literatur (Bagian II) berfungsi sebagai konteks penelitian untuk memposisikan kontribusi yang diusulkan:

\begin{itemize}
    \item U-Net~\cite{ronneberger2015unet}: Arsitektur encoder-decoder dengan skip connections untuk restorasi citra
    \item DE-GAN~\cite{souibgui2020degan}: GAN bersyarat untuk peningkatan dokumen dengan fokus kualitas visual
    \item HTR-GAN~\cite{souibgui2021enhance}: GAN dengan pengenal yang dilatih bersama untuk restorasi berorientasi HTR
    \item DocEnTr: Arsitektur berbasis Transformer dengan mekanisme atensi mandiri untuk peningkatan dokumen
\end{itemize}

Perbandingan kuantitatif langsung dengan metode deep learning terkini tidak dilakukan karena: (1) metode tersebut dilatih pada dataset berbeda dengan karakteristik degradasi yang tidak identik, sehingga perbandingan langsung tidak fair~\cite{tensmeyer2017document}, (2) implementasi official yang dapat direproduksi tidak tersedia untuk sebagian besar metode, dan (3) perbedaan arsitektur, protokol pelatihan, dan fungsi loss membuat isolasi kontribusi spesifik sangat sulit dilakukan. Pendekatan ini mengikuti praktik evaluasi yang umum dalam riset restorasi dokumen di mana perbandingan dilakukan hanya pada kondisi eksperimental yang dapat direplikasi secara identik.

\subsubsection{Studi Ablasi sebagai Validasi Komponen}

Untuk memvalidasi kontribusi setiap komponen arsitektural secara sistematis, studi ablasi inkremental dilakukan dengan menambahkan komponen kehilangan secara bertahap (Bagian V-B):

\begin{enumerate}
    \item eksperimen 01 (Hanya Piksel): Setara dengan U-Net standar untuk restorasi citra
    \item eksperimen 02 (+Adversarial): Setara dengan GAN standar tanpa komponen HTR
    \item eksperimen 03 (+Perseptual): Peningkatan dengan kehilangan perseptual berbasis VGG
    \item eksperimen 04 (+CTC): Konfigurasi optimal dengan integrasi HTR
    \item eksperimen 05 (+Fitur Pengenal): Konfigurasi lengkap untuk validasi redundansi komponen
\end{enumerate}

Pendekatan ablasi ini memberikan pemahaman komprehensif tentang kontribusi relatif setiap komponen dan memvalidasi keputusan desain arsitektural secara empiris. Hasil ablasi (Tabel~\ref{table_ablation_loss}) menunjukkan bahwa konfigurasi 4 komponen (eksperimen 04) mencapai kinerja optimal dengan trade-off terbaik antara kualitas visual dan keterbacaan HTR.

\subsection{Metrik Evaluasi}

Metrik kualitas visual dan pengenalan teks digunakan untuk evaluasi.

\subsubsection{Metrik Kualitas Visual}

\begin{itemize}
    \item Peak Signal-to-Noise Ratio (PSNR): Mengukur kesamaan tingkat piksel
    \begin{equation}
    \text{PSNR} = 10 \log_{10}\left(\frac{\text{MAX}^2}{\text{MSE}}\right)
    \end{equation}
    di mana MAX=1.0 untuk citra yang dinormalisasi ke rentang [0,1], dan MSE adalah Mean Squared Error.

    \item Structural Similarity Index (SSIM): Menilai kualitas yang dirasakan dengan mempertimbangkan luminans, kontras, dan struktur~\cite{wang2004ssim}
    \begin{equation}
    \text{SSIM}(x,y) = \frac{(2\mu_x\mu_y + c_1)(2\sigma_{xy} + c_2)}{(\mu_x^2 + \mu_y^2 + c_1)(\sigma_x^2 + \sigma_y^2 + c_2)}
    \end{equation}
\end{itemize}

\subsubsection{Metrik Deteksi Artefak (Monitoring Sekunder)}

Untuk memantau kualitas generasi selama pelatihan dan memastikan generator tidak menghasilkan artefak visual yang mengganggu keterbacaan teks, tiga metrik tambahan digunakan:

\begin{itemize}
    \item \textit{Noise Variance}: Varians dari operator Laplacian untuk mengukur derau frekuensi tinggi
    
    \item \textit{Isolated White Ratio}: Rasio piksel putih terisolasi terhadap total piksel, mendeteksi artefak derau garam
    
    \item \textit{Local Variance}: Varians dari deviasi lokal terhadap citra yang dihaluskan, mengukur kehalusan hasil restorasi
\end{itemize}

Metrik ini digunakan untuk pemantauan selama pelatihan dan validasi kualitas, tetapi tidak dilaporkan sebagai metrik evaluasi utama karena bersifat diagnostik.

\subsubsection{Metrik Kinerja HTR}

Model hibrida CNN-\textit{Transformer} yang disesuaikan pada dokumen historis digunakan untuk mengevaluasi kinerja HTR. Model ini dibekukan (\textit{frozen}) selama evaluasi untuk memastikan konsistensi pengukuran:

\begin{itemize}
    \item \textit{Character Error Rate} (CER): Persentase kesalahan tingkat karakter, dihitung menggunakan jarak penyuntingan Levenshtein:
    \begin{equation}
    \text{CER} = \frac{\text{JarakPenyuntingan}(\text{prediksi}, \text{referensi})}{\text{panjang}(\text{referensi})}
    \end{equation}
    
    \item \textit{Word Error Rate} (WER): Persentase kesalahan tingkat kata, dihitung dengan jarak penyuntingan pada tingkat kata:
    \begin{equation}
    \text{WER} = \frac{\text{JarakPenyuntingan}(\text{kata prediksi}, \text{kata referensi})}{\text{panjang}(\text{kata referensi})}
    \end{equation}
\end{itemize}

\paragraph{Klarifikasi \textit{Baseline} CER \textit{Recognizer}} Untuk menghindari kebingungan, paper ini menggunakan tiga nilai \textit{baseline} CER \textit{recognizer} beku yang berbeda:\vspace{-0.5em}
\begin{itemize}[itemsep=0pt]
    \item \textbf{CER\textsubscript{pretrain} = 33.72\%}: Kinerja \textit{recognizer} hibrida CNN-Transformer (encoder CNN dengan 6-layer transformer dan CTC decoder) saat dievaluasi pada \textit{dataset} validasi ANRI terpisah ($n$=948) untuk validasi generalisasi awal. Nilai ini digunakan untuk memverifikasi bahwa \textit{recognizer} dapat menangani tulisan tangan paleografi Belanda.
    \item \textbf{CER\textsubscript{val-GT} = 26.57\%}: Kinerja \textit{recognizer} beku pada citra bersih \textit{ground truth} dari \textit{set} validasi penelitian ini ($n$=710). Nilai ini merupakan batas atas teoretis untuk \textit{set} validasi.
    \item \textbf{CER\textsubscript{test-GT} = 34.1\%}: Kinerja \textit{recognizer} beku pada citra bersih \textit{ground truth} dari \textit{set} uji penelitian ini ($n$=712). Nilai ini merupakan batas atas teoretis untuk \textit{set} uji dan menjadi referensi utama untuk menilai kualitas restorasi.
\end{itemize}
\vspace{0.5em}
Perbedaan CER\textsubscript{val-GT} (26.57\%) vs CER\textsubscript{test-GT} (34.1\%), yaitu $\Delta$=7.53 poin persentase, mencerminkan bahwa \textit{set} uji mengandung sampel dengan kompleksitas paleografi lebih tinggi (ligatur lebih kompleks, variasi gaya tulisan lebih besar). Pola perbedaan ini konsisten dengan hasil restorasi: CER\textsubscript{val-restored} = 27.11\% vs CER\textsubscript{test-restored} = 34.9\% ($\Delta$=7.79\%), mengonfirmasi bahwa gap bukan disebabkan \textit{overfitting} melainkan perbedaan karakteristik \textit{dataset}.

\begin{table}[!ht]
\centering
\caption{Ringkasan Metrik CER pada Berbagai \textit{Dataset} dan Kondisi}
\label{table:cer_summary}
\small
\begin{tabular}{|l|c|c|c|}
\hline
\textbf{Kondisi} & \textbf{\textit{Dataset}} & \textbf{$n$} & \textbf{CER (\%)} \\
\hline
\multicolumn{4}{|c|}{\textit{\textbf{Baseline Recognizer} pada Citra Bersih}} \\
\hline
CER\textsubscript{pretrain} & ANRI validasi awal & 948 & 33.72 \\
CER\textsubscript{val-GT} & Validasi GT penelitian & 710 & 26.57 \\
CER\textsubscript{test-GT} & Uji GT penelitian & 712 & 34.1 \\
\hline
\multicolumn{4}{|c|}{\textit{\textbf{Input Terdegradasi} (Tanpa Restorasi)}} \\
\hline
CER\textsubscript{degraded} & Validasi \& Uji & 1422 & 83.4 \\
\hline
\multicolumn{4}{|c|}{\textit{\textbf{Hasil Restorasi} Metode Usulan}} \\
\hline
CER\textsubscript{val-restored} & Validasi & 710 & 27.11 \\
CER\textsubscript{test-restored} & Uji & 712 & \textbf{34.9} \\
\hline
\multicolumn{4}{|c|}{\textit{\textbf{Improvement} Relatif}} \\
\hline
Validasi & 83.4\% $\rightarrow$ 27.11\% & - & 67.5\% \\
Uji & 83.4\% $\rightarrow$ 34.9\% & - & \textbf{58.2\%} \\
\hline
\end{tabular}
\end{table}

Tabel~\ref{table:cer_summary} merangkum semua nilai CER yang dilaporkan dalam paper ini untuk transparansi dan reproduktibilitas. Penggunaan CER\textsubscript{test-GT} = 34.1\% sebagai referensi utama di seluruh paper (abstract, diskusi, kesimpulan) konsisten dengan protokol evaluasi standar yang menggunakan \textit{test set} sebagai metrik akhir.

\paragraph{Pelaporan Statistik dan Protokol Evaluasi}

Semua metrik dilaporkan sebagai rata-rata $\pm$ simpangan baku dengan interval kepercayaan 95\%, dihitung dari set validasi lengkap (n=710) atau set uji (n=712). Signifikansi statistik untuk perbandingan model diuji menggunakan \textit{t-test} sampel independen dengan $\alpha = 0.05$. Semua eksperimen menggunakan \textit{seed} tetap (\textit{seed}=42) untuk reproduktibilitas penuh.

Protokol Evaluasi Final: Semua hasil yang dilaporkan dalam makalah ini (tabel, abstract, hypothesis testing, discussion) menggunakan evaluasi konsisten pada \textit{checkpoint} terbaik (epoch 44, dipilih berdasarkan metrik validasi dengan \textit{early stopping patience}=25) yang dievaluasi pada set uji \textit{held-out} (n=712, tidak pernah dilihat selama pelatihan). Pendekatan ini mengikuti protokol akademik standar untuk mencegah \textit{overfitting} dan memastikan hasil yang tidak bias.

% ============================================================ 
% VI. HASIL DAN PEMBAHASAN
% ============================================================ 
\section{Hasil dan Pembahasan}

Bagian ini menyajikan hasil eksperimental komprehensif termasuk perbandingan kuantitatif, studi ablasi, analisis kualitatif, dan pembahasan temuan.

\subsection{Hasil Kuantitatif pada Set Uji Semi-Sintetis Realistis}

Tabel~\ref{table_synthetic_results} menyajikan hasil kuantitatif perbandingan metode yang diusulkan dengan metode klasik pada \textit{set} uji semi-sintetis realistis yang dikonstruksi dari dokumen ANRI autentik. Evaluasi difokuskan pada metode \textit{baseline} klasik (\textit{thresholding} Otsu dan Sauvola adaptif) karena metode berbasis \textit{deep learning} lainnya (U-Net, GAN standar, DE-GAN, HTR-GAN) dilatih pada dataset berbeda sehingga perbandingan langsung tidak fair. Pendekatan ini mengikuti praktik terbaik dalam riset \textit{computer vision} di mana perbandingan hanya dilakukan dengan metode yang dapat direplikasi pada kondisi eksperimental identik.

% Main quantitative results on synthetic test set
\begin{table*}[!t]
\renewcommand{\arraystretch}{1.3}
\caption{Hasil Kuantitatif pada Set Uji Semi-Sintetis Realistis (n=712, dikonstruksi dari dokumen ANRI autentik)}
\label{table_synthetic_results}
\centering
\begin{tabular}{|l|c|c|c|c|c|}
\hline
\textbf{Metode} & \textbf{PSNR (dB) $\uparrow$} & \textbf{SSIM $\uparrow$} & \textbf{\textit{F-measure} $\uparrow$} & \textbf{CER (\%)} & \textbf{WER (\%)} \\
\hline
Tanpa Perbaikan & 7.91$\pm$1.01 & 0.340$\pm$0.077 & -\textsuperscript{a} & 83.4$\pm$18.5 $\downarrow$ & 98.7$\pm$7.3 $\downarrow$ \\
Otsu \textit{Thresholding} & 4.75$\pm$2.00 & 0.261$\pm$0.099 & 0.748$\pm$0.085 & 85.7$\pm$56.7 $\downarrow$ & 112.9$\pm$69.5 $\downarrow$ \\
Sauvola Adaptif & 9.92$\pm$1.99 & 0.437$\pm$0.149 & 0.933$\pm$0.030 & 109.5$\pm$179.5 $\downarrow$ & 171.5$\pm$172.2 $\downarrow$ \\
\hline
\textbf{Yang Diusulkan} & \textbf{30.74$\pm$5.09} & \textbf{0.987$\pm$0.014} & \textbf{0.921$\pm$0.067} & \textbf{34.9$\pm$21.8} $\downarrow$ & \textbf{82.4$\pm$26.1} $\downarrow$ \\
\hline
\textit{Batas Atas (GT Bersih)} & - & - & - & \textit{34.1$\pm$22.7} & \textit{82.1$\pm$27.3} \\
\hline
\end{tabular}
\textit{Keterangan:} \textit{F-measure} tidak dapat diterapkan pada citra terdegradasi tanpa binarisasi. Metode klasik menggunakan implementasi OpenCV standar dengan parameter optimal. Metode usulan menggunakan \textit{checkpoint} epoch 44. Semua metode dievaluasi pada \textit{set} uji identik ($n$=712) menggunakan \textit{recognizer} beku yang sama. Batas Atas GT Bersih (CER\textsubscript{test-GT} = 34.1\%) menunjukkan performa teoretis maksimal \textit{recognizer} pada citra tanpa degradasi\textemdash lihat klarifikasi \textit{baseline} CER di Bagian~IV-B untuk penjelasan perbedaan angka \textit{baseline} di berbagai \textit{dataset}.
\end{table*}

Pengamatan kunci berdasarkan hasil kuantitatif:
\begin{enumerate}
    \item Perbandingan dengan metode klasik: Metode yang diusulkan mencapai PSNR 30.74$\pm$5.09 dB dan SSIM 0.987$\pm$0.014, mengungguli metode klasik terbaik (Sauvola: PSNR 9.92 dB, SSIM 0.437) dengan margin substansial. Hasil ini mengindikasikan keunggulan pendekatan berbasis \textit{deep learning} untuk restorasi dokumen terdegradasi parah.
    
    \item Peningkatan substansial pada metrik HTR: CER turun dari 83.4\% (tanpa perbaikan) menjadi 34.9\% (metode usulan), menunjukkan pengurangan absolut 48.5 poin persentase atau pengurangan relatif 58.2\%. Metode klasik justru memperburuk performa HTR (Otsu: 85.7\%, Sauvola: 109.5\%) karena binarisasi yang tidak tepat menghasilkan artefak dan hilangnya detail goresan penting. Hal ini memvalidasi pentingnya pendekatan restorasi berorientasi HTR dibanding binarisasi konvensional.
    
    \item Mendekati batas atas teoretis: CER metode usulan (34.9\%) mendekati batas atas GT bersih (34.1\%), mengindikasikan bahwa restorasi berhasil mengembalikan keterbacaan mendekati kondisi bersih asli dengan selisih hanya 0.8 poin persentase. Kedekatan ini menunjukkan bahwa hambatan utama kinerja HTR bukan lagi degradasi citra, melainkan keterbatasan kapabilitas \textit{recognizer} HTR itu sendiri.
    
    \item Peningkatan kuantitatif substansial pada metrik visual: Dibanding kondisi terdegradasi, metode usulan meningkatkan PSNR sebesar +22.83 dB (dari 7.91 ke 30.74 dB) dan SSIM sebesar +0.647 poin (dari 0.340 ke 0.987), memvalidasi efektivitas kombinasi \textit{recognizer} HTR beku dengan optimasi fungsi loss multikomponen berorientasi HTR untuk restorasi simultan kualitas visual dan keterbacaan teks.
    
    \item Keterbatasan metode klasik pada degradasi kompleks: Metode \textit{thresholding} konvensional gagal pada degradasi kompleks karena asumsi distribusi intensitas bimodal tidak berlaku pada dokumen historis dengan degradasi spasial non-uniform (\textit{foxing}, \textit{bleed-through}, penuaan heterogen). Hal ini tercermin dari CER yang lebih buruk dibanding kondisi tanpa perbaikan (Otsu: 85.7\% vs 83.4\%; Sauvola: 109.5\% vs 83.4\%).
    
    \item Varians tinggi pada metode klasik: Simpangan baku CER yang besar pada Sauvola (179.5) dan Otsu (56.7) mengindikasikan ketidakstabilan performa pada sampel dengan tingkat degradasi berbeda. Sebaliknya, metode usulan menunjukkan varians lebih rendah (21.8) berkat pembelajaran adaptif dari data yang memungkinkan generalisasi lebih baik pada berbagai tingkat degradasi.
\end{enumerate}

\subsubsection{Perbandingan Visual Metode}

Gambar~\ref{fig:method_comparison} menyajikan perbandingan visual kualitatif antara metode yang diusulkan dan metode \textit{baseline} klasik pada 5 sampel representatif dari set uji.

\begin{figure*}[!t]
\centering
\includegraphics[width=\textwidth]{../visualizations/fig_method_comparison_4methods.pdf}
\caption{Perbandingan visual metode restorasi pada 5 sampel representatif dari set uji. Baris menunjukkan: (1) citra terdegradasi masukan, (2) hasil Otsu, (3) hasil Sauvola, (4) hasil metode usulan, (5) \textit{ground truth} bersih. Metode usulan menghasilkan kualitas visual tinggi (PSNR 30.74~dB, SSIM 0.987) dan akurasi HTR lebih tinggi (CER 34.9\%) dibanding metode klasik yang menghasilkan artefak binarisasi.}
\label{fig:method_comparison}
\end{figure*}

Analisis visual mengungkapkan pola performa yang konsisten:
\begin{itemize}
    \item Kualitas restorasi: Metode usulan menghasilkan latar belakang bersih dengan pelestarian goresan optimal, sementara Otsu menghasilkan artefak \textit{salt-and-pepper} dan Sauvola cenderung menghasilkan goresan yang terlalu tebal atau putus. Perbedaan ini menjelaskan mengapa metode klasik justru memperburuk kinerja HTR meskipun dirancang untuk pembersihan dokumen.
    \item Korelasi visual-HTR: Sampel dengan PSNR/SSIM tinggi berkorelasi dengan CER rendah, memvalidasi bahwa peningkatan kualitas visual secara langsung berkontribusi pada peningkatan akurasi HTR. Korelasi ini membenarkan pendekatan penelitian yang mengoptimalkan metrik visual dan HTR secara simultan.
    \item Konsistensi performa: Metode usulan mempertahankan performa konsisten pada sampel dengan tingkat degradasi berbeda (lihat variasi PSNR: 24.25--33.84 dB), sedangkan metode klasik menunjukkan inkonsistensi (CER bervariasi drastis). Stabilitas ini penting untuk aplikasi restorasi arsip skala besar dengan variasi degradasi tinggi.
\end{itemize}

\subsubsection{Analisis Trajektori Metrik Validasi}

Gambar~\ref{fig:validation_metrics} menyajikan trajektori terperinci dari empat metrik validasi kunci selama 50 epoch dengan interval kepercayaan 95\%, menunjukkan proses konvergensi model dan keandalan statistik hasil pelatihan.

\begin{figure*}[!t]
\centering
\includegraphics[width=\textwidth]{../data_dukung/fig_validation_metrics_trajectory.pdf}
\caption{Trajektori metrik validasi dengan interval kepercayaan 95\% selama 50 \textit{epoch} (n=710): (a) PSNR mencapai 30.91±0.42~dB, (b) SSIM mencapai 0.987±0.001, (c,d) CER dan WER mendekati \textit{baseline} citra bersih. Model terbaik (\textit{epoch} 44, bintang merah) dipilih berdasarkan kriteria gabungan PSNR-CER. Interval kepercayaan sempit mengindikasikan kinerja konsisten. Evaluasi akhir pada set uji: PSNR 30.74~dB, SSIM 0.987, CER 34.9\% (Tabel~\ref{table_synthetic_results}).}
\label{fig:validation_metrics}
\end{figure*}

Observasi trajektori validasi yang mengkonfirmasi stabilitas pelatihan:
\begin{enumerate}
    \item Konvergensi bertahap pada metrik visual: Target PSNR tercapai pada epoch 38 (30.0 dB) dan target SSIM pada epoch 22 (0.95), mengindikasikan bahwa pelatihan fase \textit{warmup} visual (epoch 1--10) dan fase peningkatan CTC bertahap (epoch 11--30) berhasil membangun kemampuan rekonstruksi dasar sebelum optimasi HTR penuh.
    
    \item Stabilitas interval kepercayaan: Interval kepercayaan yang sempit untuk PSNR (±0.42 dB pada epoch 44) dan SSIM (±0.0011) menunjukkan varians rendah pada set validasi, mengkonfirmasi bahwa model belajar pola restorasi yang konsisten tanpa fluktuasi berlebihan.
    
    \item Pemilihan model berbasis kriteria gabungan: Model terbaik dipilih pada epoch 44 (bukan epoch terakhir) berdasarkan optimasi kriteria gabungan PSNR-CER pada set validasi, mengikuti protokol \textit{early stopping} dengan \textit{patience}=25. Pemilihan ini mencegah \textit{overfitting} pada fase pelatihan akhir.
    
    \item Generalisasi dari validasi ke uji: Performa pada \textit{set} uji (PSNR 30.74~dB, SSIM 0.987, CER 34.9\%) konsisten dengan tren validasi, memvalidasi bahwa model tidak mengalami \textit{overfitting}. Perbedaan CER antara validasi dan uji (27.11\% vs 34.9\%, $\Delta$=7.79\%) konsisten dengan perbedaan \textit{recognizer} beku pada citra bersih (26.57\% vs 34.1\%, $\Delta$=7.53\%), mengonfirmasi bahwa gap disebabkan kompleksitas intrinsik teks \textit{set} uji yang lebih tinggi, bukan degradasi performa model. Rasio peningkatan CER konsisten: validasi 67.5\% (83.4\%$\rightarrow$27.11\%), uji 58.2\% (83.4\%$\rightarrow$34.9\%).
\end{enumerate}

Trajektori validasi ini mengkonfirmasi efektivitas strategi \textit{curriculum learning} tiga fase yang diusulkan (Bagian III-F), di mana konvergensi metrik visual terjadi lebih awal (epoch 22--38) dibanding metrik HTR, mengindikasikan bahwa pembangunan kemampuan rekonstruksi visual menjadi fondasi untuk optimasi keterbacaan teks pada fase pelatihan penuh.

\subsection{Evaluasi Kualitatif pada Dokumen Historis Nyata (ANRI)}

Untuk mengevaluasi generalisasi model terhadap degradasi historis nyata yang tidak terlihat selama pelatihan, analisis kualitatif dilakukan pada dokumen dari Arsip Nasional Republik Indonesia (ANRI). Karena keterbatasan transkripsi acuan untuk dokumen paleografi historis, evaluasi berfokus pada penilaian kualitas visual berdasarkan observasi sistematis terhadap karakteristik restorasi.

\subsubsection{Dataset dan Metodologi Evaluasi}

Dokumen uji: Dataset terdiri dari 15 dokumen representatif dari koleksi ANRI abad ke-17 hingga ke-18 (era VOC dan Hindia Belanda, 1650--1800) dengan karakteristik degradasi autentik:
\begin{itemize}
    \item Penuaan kertas: diskolorasi kuning/coklat dengan intensitas tinggi
    \item Variasi kontras: rentang kontras dari rendah hingga tinggi, mencakup degradasi parah hingga sedang
    \item \textit{Foxing} dan noda organik: bercak kecoklatan dari \textit{fungi} dan oksidasi selulosa
    \item \textit{Bleed-through}: transparansi tinta dari halaman \textit{verso} dengan intensitas bervariasi
    \item Artefak tinta \textit{iron gall}: korosi tinta historis dengan halo gelap dan perubahan pH lokal
\end{itemize}

Karakteristik teknis dataset:
\begin{itemize}
    \item Resolusi: 300 DPI hasil pemindaian material arsip
    \item Orientasi: potret
    \item Distribusi degradasi: 2 dokumen parah (13,3\%), 8 sedang (53,3\%), 5 ringan-sedang (33,3\%)
\end{itemize}
% bagaimana perintah tambah baris baru di latex



Prosedur inferensi:
\begin{enumerate}
    \item Dokumen halaman penuh diproses menggunakan model dengan performa validasi terbaik
    \item Strategi tumpang tindih berorientasi potret untuk mempertahankan autentisitas fitur paleografi
    \item Format keluaran: TIFF tanpa kompresi untuk standar preservasi arsip
    \item Evaluasi visual: perbandingan sistematis antara masukan terdegradasi dan keluaran terestorasi untuk 5 aspek kualitas (pembersihan latar, kejelasan teks, keberadaan artefak, pelestarian goresan, kegunaan keseluruhan)
\end{enumerate}

\subsubsection{Hasil Kualitatif dan Observasi Sistematis}

Gambar~\ref{fig:anri_qualitative} menunjukkan contoh representatif restorasi pada dokumen ANRI autentik. Analisis visual mengungkapkan pola performa yang konsisten dengan karakteristik degradasi masukan.

% [GAMBAR: ANRI Qualitative Results - 6 examples arranged in 2 rows]
\begin{figure*}[!t]
\centering
\includegraphics[width=7in]{../data_dukung/evaluasi_anri/fig_anri_qualitative_results.pdf}
\caption{Restorasi kualitatif pada dokumen ANRI autentik abad ke-17--18 (n=15): setiap pasangan menunjukkan masukan terdegradasi (kiri) dan keluaran terestorasi (kanan). \textbf{Baris 1 (a--c):} kasus sukses kontras tinggi dengan pembersihan latar substansial dan pelestarian teks. \textbf{Baris 2 (d--f):} kasus menantang kontras rendah-sedang menggunakan strategi konservatif untuk menghindari artefak.}
\label{fig:anri_qualitative}
\end{figure*}

\paragraph{Observasi Berdasarkan Kategori Degradasi (n=15 dokumen)}

\textit{Dokumen Kontras Tinggi (n=5, 33,3\%):}
\begin{itemize}
    \item Pembersihan latar: \textit{foxing}, noda penuaan, dan bercak organik dihilangkan secara efektif
    \item Pelestarian teks: ketebalan goresan dan morfologi karakter dipertahankan tanpa peningkatan berlebihan atau artefak penipisan
    \item Penghilangan \textit{bleed-through}: transparansi teks \textit{verso} diminimalkan dengan pelestarian yang baik terhadap teks \textit{recto}
    \item Artefak: artefak minimal, keluaran sesuai untuk alur kerja transkripsi
\end{itemize}

\textit{Dokumen Kontras Sedang (n=8, 53,3\%):}
\begin{itemize}
    \item Sukses parsial: derau latar berkurang substansial, tetapi artefak residual terlihat pada wilayah dengan \textit{bleed-through} berat
    \item Peningkatan teks: keterbacaan karakter meningkat, namun peningkatan lebih konservatif dibanding kasus kontras tinggi
    \item \textit{Trade-off}: model memprioritaskan negatif palsu (derau terlewat) daripada positif palsu (goresan halusinasi) untuk mempertahankan autentisitas
    \item Kegunaan: keluaran lebih baik dibanding metode binarisasi tradisional untuk degradasi sedang
\end{itemize}

\textit{Dokumen Kontras Rendah (n=2, 13,3\%):}
\begin{itemize}
    \item Peningkatan terbatas: model mengadopsi strategi konservatif untuk menghindari pembangkitan artefak palsu
    \item Tantangan tinta memudar: tinta dengan kontras rendah hanya mendapat peningkatan minimal
    \item Pelestarian autentisitas: tidak ada halusinasi atau peningkatan berlebihan yang dapat menyesatkan analisis paleografi
\end{itemize}

\subsubsection{Analisis Performa Berdasarkan Kontras Masukan}

Untuk memvalidasi hipotesis bahwa efektivitas restorasi berkorelasi dengan tingkat kontras masukan, analisis kuantitatif dilakukan terhadap distribusi degradasi dataset ANRI.

\begin{table}[!t]
\renewcommand{\arraystretch}{1.3}
\caption{Performa Restorasi pada Dokumen ANRI Berdasarkan Tingkat Kontras Masukan (n=15)}
\label{table:anri_contrast_analysis}
\centering
\small
\begin{tabular}{|l|c|c|c|}
\hline
\textbf{Kategori Kontras} & \textbf{n} & \textbf{\% Total} & \textbf{Efektivitas Teramati} \\
\hline
Tinggi ($\geq$60) & 5 & 33,3\% & Baik Sekali$^a$ \\
Sedang (30--60) & 8 & 53,3\% & Baik$^b$ \\
Rendah ($<$30) & 2 & 13,3\% & Terbatas$^c$ \\
\hline
\multicolumn{4}{l}{\footnotesize $^a$Baik Sekali: pembersihan latar >80\%, artefak minimal} \\
\multicolumn{4}{l}{\footnotesize $^b$Baik: perbaikan substansial, derau residual minor} \\
\multicolumn{4}{l}{\footnotesize $^c$Terbatas: peningkatan konservatif, tanpa halusinasi} \\
\end{tabular}
\end{table}

Temuan analisis kontras menunjukkan tiga pola kunci:
\begin{enumerate}
    \item Efek ambang kontras: Degradasi performa model mulai terlihat pada kontras $<$30, mengindikasikan keterbatasan distribusi \textit{dataset} pelatihan yang mungkin tidak sepenuhnya mencakup skenario kontras rendah. Hasil ini konsisten dengan distribusi kontras \textit{dataset} semi-sintetis realistis yang memiliki median kontras 35--45.
    
    \item Strategi konservatif pada ambiguitas: Pada kasus dengan SNR rendah, model memprioritaskan presisi daripada \textit{recall} untuk menghindari pembangkitan artefak palsu---pendekatan yang penting untuk preservasi autentisitas arsip historis dimana integritas dokumen lebih penting daripada peningkatan visual maksimal.
    
    \item Sukses pada degradasi tipikal: Dari 15 dokumen ANRI, 13 dokumen (86,7\%) dengan kontras $\geq$30 menunjukkan perbaikan berguna untuk alur kerja transkripsi, sementara 10 dokumen (66,7\%) menunjukkan peningkatan substansial dengan pembersihan latar >70\% dan artefak minimal. Perbedaan ini mencerminkan gradasi kualitas restorasi: kategori "berguna" mencakup semua dokumen dengan peningkatan minimal hingga optimal, sedangkan "signifikan" terbatas pada dokumen dengan pembersihan latar substansial.
\end{enumerate}

Keterbatasan generalisasi yang teridentifikasi:
\begin{enumerate}
    \item Pemudaran ekstrem: Dokumen dengan kontras $<$20 (1/15, 6,7\%) mendapat peningkatan minimal karena distribusi kontras \textit{dataset} pelatihan semi-sintetis berbeda. Hal ini konsisten dengan observasi pada \textit{set} validasi dimana performa menurun pada PSNR masukan $<$8 dB (wilayah di luar distribusi pelatihan).
    
    \item Variansi \textit{bleed-through}: Pola \textit{bleed-through} dunia nyata dengan tata letak halaman \textit{verso} yang kompleks lebih menantang dibanding augmentasi sintetis yang menggunakan \textit{flip} horizontal sederhana, menghasilkan sukses parsial pada 4/15 kasus (26,7\%).
    
    \item Artefak tinta \textit{iron gall}: Degradasi spesifik kimiawi historis (efek halo, pergeseran warna pH-induced) tidak sepenuhnya tereplikasi dalam pelatihan sintetis, menyebabkan sesekali artefak \textit{over-suppression} pada 2/15 dokumen (13,3\%).
\end{enumerate}

Implikasi praktis untuk digitisasi arsip: Hasil evaluasi ANRI mengonfirmasi bahwa model dapat menggeneralisasi ke dokumen historis nyata dengan karakteristik degradasi yang tumpang tindih dengan distribusi pelatihan sintetis (kontras $\geq$30, penuaan sedang hingga berat). Untuk kasus ekstrem di luar domain pelatihan, model mengadopsi strategi konservatif yang mempertahankan autentisitas dokumen, menghindari mode kegagalan katastrofik seperti halusinasi atau korupsi struktural. Pendekatan ini sesuai dengan prinsip preservasi arsip dimana negatif palsu (peluang restorasi terlewat) lebih dapat diterima daripada positif palsu (artefak terintroduksi yang menyesatkan interpretasi historis).

\subsubsection{Validasi Ahli (Asesmen Paleografi)}

Untuk memvalidasi kualitas restorasi dari perspektif paleografi profesional, penelitian ini melibatkan dua ahli paleografi independen dengan pengalaman lebih dari 10 tahun dalam analisis dokumen historis periode VOC dan Hindia Belanda abad ke-16 hingga ke-18.

Metodologi validasi:
\begin{itemize}
    \item Panel ahli: 
    \begin{itemize}
        \item[$-$] \textit{Penilai 1:} Ahli paleografi, spesialisasi skrip Belanda kuno dan Latin abad ke-16--18, 16 tahun pengalaman konservasi dokumen VOC (afiliasi dirahasiakan untuk menjaga kebutaan evaluasi)
        \item[$-$] \textit{Penilai 2:} Ahli restorasi arsip historis, spesialisasi dokumen administratif Hindia Belanda, berpengalaman dalam restorasi arsip fisik dan digitisasi koleksi (afiliasi dirahasiakan)
    \end{itemize}
    \item Protokol asesmen tertutup: Evaluasi independen pada 15 pasangan dokumen (terdegradasi-terestorasi) tanpa informasi teknis model atau metodologi restorasi. Kedua penilai tidak mengetahui bahwa restorasi dilakukan menggunakan GAN atau pendekatan berbasis HTR. Setiap dokumen dinilai menggunakan skala Likert 4 poin (1=Buruk, 2=Cukup, 3=Baik, 4=Sangat Baik) untuk empat kriteria kualitas.
    \item Kriteria evaluasi:
    \begin{enumerate}
        \item \textit{Keterbacaan Teks:} Kemudahan membaca karakter paleografi, terutama huruf \textit{minuscule} dan ligatur khas abad ke-17
        \item \textit{Autentisitas Paleografi:} Preservasi ciri khas tulisan tangan era (\textit{ductus}, tekanan pena, variasi ketebalan goresan)
        \item \textit{Minimasi Artefak} (\textit{inverse scoring}): Tingkat gangguan artefak restorasi terhadap interpretasi paleografi (skor tinggi = artefak minimal)
        \item \textit{Kegunaan Transkripsi:} Kesesuaian keluaran untuk alur kerja digitisasi arsip dan transkripsi sistematis
    \end{enumerate}
    \item Reliabilitas antarpenilai: Cohen's kappa dihitung untuk mengukur konsistensi kesepakatan antara dua penilai independen, dengan ambang batas $\kappa > 0.60$ untuk kesepakatan substansial (Landis \& Koch, 1977)
\end{itemize}

\begin{table}[!t]
\renewcommand{\arraystretch}{1.3}
\caption{Hasil Validasi Ahli Paleografi pada Dokumen ANRI (n=15)}
\label{table:anri_expert_validation}
\centering
\small
\begin{tabular}{|l|c|c|c|}
\hline
\textbf{Kriteria} & \textbf{Penilai 1} & \textbf{Penilai 2} & \textbf{Rerata} \\
\hline
Keterbacaan Teks & 3,47 & 3,33 & 3,40 \\
Autentisitas Paleografi & 3,60 & 3,53 & 3,57 \\
Minimasi Artefak (inv.) & 3,27 & 3,13 & 3,20 \\
Kegunaan Transkripsi & 3,53 & 3,47 & 3,50 \\
\hline
\textbf{Skor Keseluruhan} & \textbf{3,47} & \textbf{3,37} & \textbf{3,42} \\
\hline
\multicolumn{4}{l}{\footnotesize Cohen's kappa (antarpenilai): $\kappa$ = 0,78 (\textit{substantial agreement})} \\
\multicolumn{4}{l}{\footnotesize Skala: 1=Buruk, 2=Cukup, 3=Baik, 4=Sangat Baik} \\
\multicolumn{4}{l}{\footnotesize 95\% CI untuk rerata keseluruhan: [3,28, 3,56]} \\
\end{tabular}
\end{table}

Distribusi skor per kategori kontras:

Analisis stratifikasi menunjukkan korelasi kuat antara kontras masukan dan skor paleografi:

\begin{itemize}
    \item Kontras Tinggi ($\geq$60, n=5): Rerata skor 3,75/4,0. Kedua penilai memberikan apresiasi tinggi terhadap preservasi detail mikroskopis seperti \textit{hairline strokes} dan modulasi tekanan pena yang krusial untuk analisis paleografi komparatif.
    
    \item Kontras Sedang (30--60, n=8): Rerata skor 3,38/4,0. Evaluasi positif dengan catatan minor bahwa sesekali terjadi \textit{over-smoothing} pada region \textit{bleed-through} berat, dapat mengurangi kejelasan diferensiasi antara teks \textit{recto} dan \textit{verso}, namun tidak menghambat transkripsi utama.
    
    \item Kontras Rendah ($<$30, n=2): Rerata skor 2,63/4,0. Penilai mengapresiasi strategi konservatif model yang menghindari \textit{false reconstruction}. Dokumen kontras sangat rendah dinilai mempertahankan integritas paleografi meskipun peningkatan visual minimal.
\end{itemize}

Umpan balik kualitatif ahli:

\textit{Penilai 1 (ahli paleografi):} ``Restorasi menunjukkan pemahaman implisit terhadap karakteristik skrip Belanda abad ke-17, khususnya preservasi \textit{aspect ratio} huruf-huruf \textit{Gothic minuscule} dan ligatur khas seperti `st', `ct', dan `ck'. Pada dokumen kontras tinggi, pembersihan \textit{foxing} sangat efektif tanpa mengorbankan detail superfisial seperti \textit{pen lift marks} yang penting untuk analisis autentisitas. Catatan kritik: pada 2--3 kasus sedang, terdapat \textit{slight blur} pada region dengan tinta \textit{iron gall} terkorosi, namun ini tidak mengurangi utilitas untuk transkripsi semantik.''

\textit{Penilai 2 (ahli restorasi arsip):} ``Dari perspektif arsip digital, keluaran model sangat sesuai untuk digitisasi massal dengan koreksi manual minimal. Reduksi derau latar pada 86,7\% dokumen (13/15) memfasilitasi OCR/HTR \textit{downstream} dengan signifikan. Apresiasi khusus untuk pendekatan konservatif pada dokumen memudar ekstrem---model tidak melakukan \textit{creative reconstruction} yang dapat menyesatkan interpretasi historis. Untuk alur kerja ANRI, disarankan penggunaan untuk praproses \textit{batch} dengan verifikasi ahli pada kasus kontras $<$30.''

Kesepakatan antarpenilai: Cohen's kappa = 0,78 mengindikasikan \textit{substantial agreement} (Landis \& Koch, 1977), memvalidasi konsistensi evaluasi paleografi. \textit{Disagreement} terbesar terjadi pada kriteria ``Minimasi Artefak'' untuk dokumen kontras sedang (30--60), dimana Penilai 1 lebih toleran terhadap \textit{residual noise} dibanding Penilai 2 yang memprioritaskan \textit{cleanliness} untuk otomasi OCR. Namun, kedua penilai sepakat pada \textit{ranking} relatif dokumen (Spearman's $\rho$ = 0,91, p$<$0,001).

Implikasi untuk digitisasi arsip: Validasi ahli mengonfirmasi bahwa model memenuhi standar kualitas paleografi untuk tiga aplikasi kunci: (1) Preservasi arsip digital---autentisitas fitur paleografi terjaga (skor 3,57/4,0), memungkinkan analisis komparatif skrip dan identifikasi juru tulis (\textit{scribe identification}); (2) Aksesibilitas publik---peningkatan keterbacaan (skor 3,40/4,0) mengurangi hambatan untuk peneliti nonspesialis; (3) Alur kerja semiotomatis---kegunaan transkripsi (skor 3,50/4,0) mendukung integrasi dengan \textit{pipeline} HTR modern dengan supervisi ahli minimal.

\subsection{Analisis Korelasi Tingkat Degradasi}

Untuk memahami hubungan antara tingkat degradasi dokumen dan efektivitas restorasi, dilakukan analisis korelasi komprehensif menggunakan set validasi (n=710 sampel). Analisis ini menjawab pertanyaan penelitian kunci: \textit{``Apakah restorasi lebih efektif untuk dokumen dengan degradasi berat atau ringan?''}

\begin{figure*}[!t]
\centering
\includegraphics[width=\textwidth]{../data_dukung/fig_combined_correlation_analysis.pdf}
\caption{Analisis korelasi tingkat degradasi pada set validasi (n=710): (a) korelasi negatif kuat (Pearson r=--0,963, p<0,001) antara SNR dan peningkatan CER menunjukkan restorasi lebih efektif pada dokumen sangat terdegradasi, (b) PSNR vs CER pascarestaorasi menunjukkan model membawa semua tingkat degradasi mendekati nilai referensi citra bersih (26,6\%, garis hijau), (c) SSIM vs CER menunjukkan korelasi lemah mengindikasikan kesamaan struktural bukan prediktor tunggal untuk akurasi HTR, (d) perbandingan CER sebelum-sesudah menunjukkan mayoritas sampel mengalami perbaikan (poin di bawah diagonal y=x, abu-abu). Catatan: Data dari set validasi (n=710) dengan statistik PSNR 30,91±5,60 dB, CER 27,11±20,83\%, SSIM 0,987±0,014. Hasil akhir pada set uji (n=712) dilaporkan dalam Tabel~\ref{table_synthetic_results}: PSNR 30.74 dB, CER 34.9\%, SSIM 0.987.}
\label{fig:degradation_correlation}
\end{figure*}

\subsubsection{Temuan Korelasi}

\textit{SNR terhadap Peningkatan CER (Panel a):}
\begin{itemize}
    \item Korelasi negatif sangat kuat: Pearson's r = --0,963, R² = 0,927, p < 0,001
    \item Interpretasi: dokumen dengan SNR rendah (degradasi berat) mendapatkan peningkatan CER lebih besar dibanding dokumen dengan SNR tinggi (degradasi ringan)
    \item Implikasi praktis: model menunjukkan efektivitas lebih tinggi untuk dokumen historis yang sangat terdegradasi, dengan peningkatan CER yang berbanding terbalik terhadap tingkat SNR awal
    \item Validasi empiris: korelasi kuat (R² > 0,9) mengonfirmasi bahwa tingkat degradasi awal merupakan prediktor signifikan untuk manfaat restorasi
\end{itemize}

\textit{PSNR terhadap CER Terestorasi (Panel b):}
\begin{itemize}
    \item Korelasi lemah: Pearson's r mendekati nol
    \item Interpretasi: setelah restorasi, PSNR (kualitas citra) tidak berkorelasi kuat dengan CER (akurasi HTR), mengindikasikan bahwa peningkatan kualitas visual tidak secara otomatis menjamin peningkatan akurasi pengenalan teks
    \item Implikasi praktis: model cenderung membawa berbagai tingkat degradasi ke kinerja HTR yang relatif seragam mendekati nilai dasar bersih
    \item Visualisasi densitas: konsentrasi di sekitar CER dasar (26,6\%) mengonfirmasi normalisasi kinerja HTR pascarestaorasi
\end{itemize}

\textit{SSIM terhadap CER (Panel c):}
\begin{itemize}
    \item Korelasi lemah teramati antara SSIM dan CER, mengindikasikan SSIM saja tidak cukup untuk memprediksi akurasi HTR
    \item Wawasan: kesamaan struktural (SSIM) mengukur kualitas visual global, tetapi kinerja HTR bergantung pada detail tingkat karakter (ketajaman goresan, kontras lokal) yang tidak sepenuhnya tercakup oleh metrik struktural
    \item Implikasi: justifikasi untuk evaluasi multimetrik (PSNR + SSIM + CER) daripada mengandalkan metrik kualitas visual tunggal untuk memvalidasi efektivitas restorasi berorientasi HTR
\end{itemize}

\textit{CER Sebelum vs Sesudah (Panel d):}
\begin{itemize}
    \item Mayoritas sampel terletak di bawah garis diagonal (y=x), mengonfirmasi bahwa restorasi umumnya meningkatkan atau mempertahankan akurasi HTR dibanding citra terdegradasi tanpa pemrosesan
    \item Distribusi poin menunjukkan bahwa restorasi memberikan manfaat konsisten pada spektrum tingkat degradasi, dengan kasus penurunan kinerja minor (<2\% peningkatan CER) yang terbatas pada sampel dengan degradasi sangat ringan di mana intervensi restorasi kurang diperlukan
\end{itemize}

\subsubsection{Implikasi untuk Digitalisasi Dokumen Historis}

Berdasarkan analisis korelasi, implikasi operasional untuk aplikasi digitalisasi arsip dapat diidentifikasi:

\begin{enumerate}
    \item Strategi prioritisasi: arsip dapat memprioritaskan dokumen dengan degradasi berat (SNR rendah) untuk restorasi, karena korelasi kuat (r = --0,963) mengonfirmasi bahwa dokumen terdegradasi berat mendapatkan peningkatan akurasi HTR lebih besar dibanding dokumen dengan degradasi ringan
    \item Jaminan kualitas: untuk dokumen dengan degradasi ringan, restorasi tetap bermanfaat untuk normalisasi kualitas dan konsistensi \textit{pipeline}, meskipun dengan margin peningkatan lebih kecil mengingat akurasi dasar yang sudah tinggi
    \item Penilaian risiko: distribusi perbaikan yang konsisten memberikan kepercayaan untuk alur kerja pemrosesan berkelompok, dengan risiko penurunan kinerja terbatas pada kasus degradasi sangat ringan (margin <2\%)
    \item Pemilihan metrik: PSNR/SSIM berguna untuk pemantauan kualitas visual, tetapi CER/WER esensial untuk memvalidasi efektivitas restorasi berorientasi HTR karena korelasi lemah antara metrik visual dan akurasi pengenalan teks
\end{enumerate}

\subsubsection{Analisis \textit{Domain Shift}: Sintetis terhadap Nyata}

Performa yang konsisten pada dokumen ANRI nyata (meskipun hanya dilatih pada degradasi sintetis) menunjukkan bahwa \textit{pipeline} degradasi sintetis berhasil mensimulasikan karakteristik degradasi historis yang relevan. Namun, beberapa efek kesenjangan domain teramati:

\begin{table}[!t]
\renewcommand{\arraystretch}{1.15}
\caption{Perbandingan karakteristik data pelatihan sintetis terhadap dokumen ANRI nyata: analisis kesenjangan domain antara degradasi sintetis dan autentik.}
\label{table:synthetic_vs_real_comparison}
\centering
\begin{tabular}{|p{2.5cm}|p{2.8cm}|p{2.8cm}|}
\hline
\textbf{Aspek} & \textbf{Sintetis} & \textbf{Nyata ANRI} \\
\hline
Degradasi & Pipeline terkontrol & Penuaan alami \\
& 6 tahap & 250--450 tahun \\
\hline
Keseragaman & Relatif seragam & Sangat bervariasi \\
\hline
Keparahan & Sedang (15--25 dB) & Ekstrem (<10 dB) \\
\hline
Latar belakang & Tekstur sampel & Kimia kertas menua \\
\hline
Jenis tinta & Simulasi memudar & Korosi iron gall \\
\hline
Ground truth & Data berpasangan & Tanpa referensi \\
\hline
\end{tabular}
\end{table}

\paragraph{Efek Kesenjangan Domain yang Teramati}
\begin{enumerate}
    \item Bias restorasi konservatif: model cenderung melakukan peningkatan kurang optimal pada degradasi ekstrem (PSNR <10 dB) karena distribusi pelatihan tidak mencakup tingkat keparahan yang sama. Ini merupakan \textit{trade-off} yang dapat diterima untuk menghindari halusinasi
    
    \item Sensitivitas pola latar belakang: dokumen nyata dengan margin dekoratif kompleks atau tepi ornamental kadang tersupresi sebagian karena kebingungan dengan derau. Frekuensi: 2/15 dokumen
    
    \item Variansi kimia tinta \textit{iron gall}: pola korosi kimia berbeda dengan degradasi yang disimulasikan sehingga pemulihan hanya parsial. Data pelatihan perlu diaugmentasi dengan model degradasi spesifik korosi
\end{enumerate}

\paragraph{Analisis Generalisasi}
Meskipun dilatih secara eksklusif pada data berpasangan sintetis, model menunjukkan generalisasi yang wajar ke pola degradasi dunia nyata. Dari 15 dokumen ANRI: 10 dokumen (66,7\%) menunjukkan peningkatan signifikan dengan pembersihan latar >70\% dan artefak minimal; 3 dokumen (20,0\%) menunjukkan peningkatan moderat dengan derau residual yang dapat diterima; dan 2 dokumen (13,3\%) dengan degradasi ekstrem (kontras $<$20) mendapat peningkatan terbatas karena berada di luar domain pelatihan. Secara keseluruhan, 13 dari 15 dokumen (86,7\%) mendapat manfaat berguna untuk alur kerja transkripsi, mengindikasikan bahwa kesenjangan domain sintetis-ke-nyata dapat diatasi melalui desain \textit{pipeline} degradasi yang hati-hati. Perbedaan antara "peningkatan signifikan" (66,7\%) dan "manfaat berguna" (86,7\%) mencerminkan spektrum kualitas restorasi dari optimal hingga minimal namun tetap fungsional.

\subsubsection{Keterbatasan dan Pekerjaan Mendatang untuk Evaluasi ANRI}

\paragraph{Keterbatasan Evaluasi Saat Ini}
\begin{itemize}
    \item Tidak ada metrik HTR kuantitatif: karena dokumen ANRI tidak memiliki transkripsi \textit{ground truth} yang andal (usia 250--450 tahun, transkripsi paleografi memerlukan keahlian khusus), CER/WER tidak dapat diukur secara objektif. Evaluasi terbatas pada penilaian visual kualitatif dan penilaian pakar
    
    \item Ukuran sampel kecil: evaluasi pada 15 dokumen halaman penuh yang dipilih secara representatif dari koleksi arsip ANRI. Bias sampel mungkin ada
    
    \item Tidak ada pengukuran PSNR: dokumen nyata tidak memiliki citra referensi bersih, sehingga PSNR tidak dapat dihitung. Penilaian kualitas visual sepenuhnya subjektif
\end{itemize}

\paragraph{Rencana Pekerjaan Mendatang}
\begin{enumerate}
    \item Pembuatan subset \textit{ground truth}: kerja sama dengan paleografer ANRI untuk transkripsi manual 100--200 citra baris yang representatif, memungkinkan evaluasi HTR kuantitatif
    
    \item Adaptasi domain melalui \textit{fine-tuning}: \textit{fine-tuning few-shot} pada subset ANRI (50--100 citra) dengan label semu untuk meningkatkan generalisasi terhadap kasus degradasi ekstrem
    
    \item \textit{Pipeline} degradasi yang diperluas: augmentasi degradasi sintetis dengan model korosi kimia dan tingkat keparahan ekstrem untuk cakupan domain yang lebih baik
    
    \item Validasi skala besar: perluasan evaluasi ke 100+ dokumen dengan metrik kualitas otomatis (peningkatan kontras, rasio pengurangan derau) sebagai proksi untuk keterbacaan teks
\end{enumerate}

\paragraph{Kesimpulan Evaluasi ANRI}
Evaluasi kualitatif pada dokumen ANRI nyata menunjukkan bahwa metode yang diusulkan berhasil melakukan generalisasi dari data pelatihan sintetis ke degradasi historis dunia nyata dengan tingkat kesuksesan 66,7\%. Penilaian paleografer pakar memberikan validasi independen bahwa kualitas restorasi sesuai untuk alur kerja digitalisasi arsip (skor kesesuaian rerata 3,42/4,0). Meskipun terdapat keterbatasan metrik kuantitatif (tidak ada \textit{ground truth}), bukti kualitatif yang kuat---dikombinasikan dengan validasi pakar dan keandalan antarpenilai yang tinggi---memberikan kepercayaan bahwa model dapat diterapkan untuk aplikasi arsip praktis pada dokumen historis nyata. Pekerjaan mendatang akan fokus pada pembuatan subset \textit{ground truth} untuk memungkinkan evaluasi kuantitatif yang ketat dan adaptasi domain untuk meningkatkan kinerja pada kasus degradasi ekstrem.

\subsection{Studi Ablasi}

Studi ablasi dilakukan untuk mengukur kontribusi setiap komponen.

\subsubsection{Ablasi Arsitektur Diskriminator}

Tabel~\ref{table_ablation_discriminator} membandingkan varian diskriminator untuk memvalidasi kontribusi arsitektur \textit{dual-modal} terhadap kinerja restorasi.

% Ablation study discriminator architecture
\begin{table}[!t]
\renewcommand{\arraystretch}{1.3}
\caption{Studi Ablasi Arsitektur Diskriminator pada Set Validasi (n=710): perbandingan diskriminator CNN tunggal (baseline) terhadap arsitektur dual-modal yang diusulkan untuk memvalidasi kontribusi cabang teks dan mekanisme cross-attention bilateral terhadap kualitas restorasi (PSNR/SSIM) dan akurasi HTR (CER). Kedua varian dilatih dengan konfigurasi identik untuk isolasi dampak arsitektural.}
\label{table_ablation_discriminator}
\centering
\begin{tabular}{|l|c|c|c|}
\hline
\textbf{Diskriminator} & \textbf{PSNR (dB)} & \textbf{SSIM} & \textbf{CER (\%)} \\
\hline
Hanya CNN (baseline) & 30.63 & 0.9854 & 27.12 \\
\textbf{\textit{Dual-Modal} (usulan)} & \textbf{30.91} & \textbf{0.9869} & \textbf{27.11} \\
\hline
\multicolumn{4}{|l|}{$\Delta$ (Dual-Modal - CNN): +0.28 dB, +0.0015, -0.01\%} \\
\hline
\end{tabular}
\vspace{2mm}

\small{Kedua eksperimen menggunakan \textit{generator} identik (Enhanced U-Net), 50 \textit{epoch}, \textit{checkpoint} epoch 44 dipilih pada \textit{set} validasi ($n$=710) berdasarkan kriteria gabungan PSNR-CER. Perbedaan hanya pada arsitektur diskriminator untuk mengisolasi kontribusi komponen \textit{dual-modal}. Hasil evaluasi akhir pada \textit{set} uji ($n$=712) dilaporkan dalam Tabel~\ref{table_synthetic_results}: PSNR 30.74~dB, CER 34.9\%. \textit{Catatan}: CER\textsubscript{val-GT} = 26.57\% vs CER\textsubscript{test-GT} = 34.1\% ($\Delta$=7.53\%) mencerminkan perbedaan kompleksitas teks antar \textit{split}, bukan \textit{recognizer drift}.}
\end{table}

Analisis Empiris: Hasil eksperimen menunjukkan perbedaan yang sangat minimal ($\Delta$PSNR = +0.28 dB, $\Delta$SSIM = +0.0015, $\Delta$CER = -0.01\%), tidak signifikan secara statistik (p > 0.05). Secara praktis, kedua arsitektur menghasilkan keluaran yang identik dari perspektif pengguna, dengan metrik kualitas visual dan HTR yang sebanding.

Analisis Akar Permasalahan: Analisis mendalam mengidentifikasi lima faktor yang menjelaskan temuan ini:

(1) Hambatan \textit{Generator}: Kedua eksperimen menggunakan arsitektur \textit{generator} yang identik (\textit{enhanced U-Net}). Batas kualitas keluaran ditentukan oleh kapasitas \textit{generator}, bukan kompleksitas umpan balik diskriminator. \textit{Generator} dengan kapasitas terbatas tidak dapat menghasilkan keluaran yang jauh lebih baik meskipun diskriminator memberikan sinyal yang lebih sophisticated.

(2) Keterbatasan Modus Teks Prediksi: Konfigurasi pelatihan menggunakan modus teks prediksi, yang berarti masukan teks ke diskriminator \textit{dual-modal} berasal dari keluaran \textit{generator} (CER 27.11\%), bukan \textit{ground truth}. Cabang LSTM menerima teks dengan tingkat kesalahan 27\%, sehingga tidak dapat belajar dari pola teks yang akurat. Mekanisme atensi \textit{cross-modal} terbatas oleh kualitas teks prediksi yang mengandung derau.

(3) Dominasi Fungsi Loss: Analisis kontribusi fungsi loss menunjukkan bahwa kehilangan \textit{adversarial} hanya berkontribusi kecil terhadap total sinyal pelatihan. Perbedaan arsitektur diskriminator hanya memengaruhi komponen minoritas dari total gradien.

(4) Kendala \textit{Recognizer} Beku: \textit{Recognizer} HTR dibekukan dengan CER\textsubscript{val-GT} = 26.57\% pada citra bersih \textit{set} validasi. Citra yang dihasilkan mencapai CER\textsubscript{val-restored} = 27.11\%, hanya 0.54 poin persentase dari batas atas teoretis. Ruang perbaikan sangat terbatas terlepas dari arsitektur diskriminator.

(5) Kinerja Hampir Optimal: Model sudah mencapai dataran kinerja. Peningkatan marjinal diharapkan dari modifikasi arsitektural tanpa perubahan paradigma dalam protokol pelatihan.

Interpretasi Ilmiah: Temuan ini bukan kegagalan arsitektur \textit{dual-modal}, melainkan wawasan ilmiah berharga yang menunjukkan bahwa dalam kerangka GAN-HTR: (a) desain protokol pelatihan (modus supervisi teks, keseimbangan bobot kehilangan) memiliki dampak lebih signifikan dibanding kompleksitas arsitektural diskriminator, (b) kapasitas \textit{generator} merupakan hambatan utama untuk peningkatan kualitas keluaran, (c) modus teks prediksi membatasi efektivitas \textit{dual-modal} karena LSTM belajar dari teks yang mengandung derau, bukan \textit{ground truth}.

Validitas Studi Ablasi: Meskipun hanya membandingkan dua varian (CNN tunggal vs \textit{dual-modal}), studi ini tetap valid untuk tujuan utama: memvalidasi apakah penambahan cabang teks dan mekanisme \textit{cross-attention} memberikan peningkatan signifikan. Hasil menunjukkan bahwa dalam konfigurasi saat ini, kontribusi arsitektural \textit{dual-modal} bersifat marjinal. Studi lebih lanjut dengan varian \textit{dual-modal} tanpa \textit{cross-attention} atau dengan supervisi teks \textit{ground truth} dapat memberikan wawasan tambahan tentang komponen mana yang paling berpengaruh.

Catatan Metodologis: Evaluasi ablasi dilakukan pada set validasi (n=710) yang digunakan untuk pemilihan model (\textit{checkpoint} epoch 44), mengikuti praktik standar studi ablasi dimana komponen dievaluasi pada \textit{dataset development} untuk iterasi cepat dan mencegah kontaminasi set uji. Pendekatan ini konsisten dengan protokol pemilihan model dimana \textit{checkpoint} epoch 44 dioptimasi berdasarkan kriteria gabungan PSNR-CER pada validation set yang sama. Test set (n=712) di-\textit{reserve} untuk evaluasi generalisasi akhir yang dilaporkan dalam Tabel~\ref{table_synthetic_results}. Perbedaan metrik antara validation (PSNR 30.91 dB, CER 27.11\%) dan test (PSNR 30.74 dB, CER 34.9\%) mencerminkan variasi alami sampel berbeda, bukan \textit{overfitting}, dan mengkonfirmasi bahwa temuan ablasi tetap valid untuk generalisasi.

Implikasi untuk Penelitian Lanjutan: Untuk memaksimalkan keuntungan arsitektur \textit{dual-modal}, diperlukan modifikasi: (1) supervisi teks \textit{ground truth}, (2) peningkatan bobot \textit{adversarial} untuk memperkuat pengaruh diskriminator, (3) arsitektur \textit{generator} sadar-teks dengan \textit{encoder} teks terintegrasi, atau (4) strategi pelatihan progresif (fase 1: CNN tunggal fokus pada visual, fase 2: \textit{dual-modal} memperbaiki keterbacaan teks).

\subsubsection{Studi Ablasi Inkremental: Kontribusi Komponen Kehilangan}

Untuk memahami kontribusi setiap komponen kehilangan secara sistematis, studi ablasi inkremental dilakukan dengan menambahkan komponen kehilangan secara bertahap. Lima eksperimen dilatih selama 15 \textit{epoch} dengan \textit{dataset} identik untuk mengisolasi dampak setiap komponen. Tabel~\ref{table_ablation_loss} dan Gambar~\ref{fig:ablation_comparison} menyajikan hasil komprehensif dari studi ablasi.

% Incremental ablation table
\begin{table*}[!t]
\renewcommand{\arraystretch}{1.3}
\caption{Studi Ablasi Inkremental Komponen Kehilangan: kontribusi setiap komponen kehilangan terhadap kualitas restorasi dokumen. Lima konfigurasi eksperimen dengan penambahan inkremental komponen: (1) \textit{baseline} kehilangan piksel (L1), (2) penambahan kehilangan \textit{adversarial}, (3) penambahan kehilangan perseptual (VGG-19), (4) penambahan kehilangan CTC untuk kesadaran HTR, (5) penambahan kehilangan fitur pengenal. Setiap eksperimen dilatih selama 15 \textit{epoch} dengan \textit{hyperparameter} identik untuk mengisolasi dampak individual setiap komponen. Metrik evaluasi mencakup PSNR (kualitas piksel), SSIM (kesamaan struktural), dan CER (kinerja HTR) pada \textit{set} validasi ($n$=710 citra).}
\label{table_ablation_loss}
\centering
\small
\begin{tabular}{|c|l|c|c|c|c|}
\hline
\textbf{Eks} & \textbf{Komponen Kehilangan} & \textbf{PSNR (dB)} & \textbf{SSIM} & \textbf{CER (\%)} & \textbf{$\Delta$ PSNR} \\
\hline
1 & $\mathcal{L}_{\text{pixel}}$ (L1) & 24.59 $\pm$ 4.10 & 0.9614 $\pm$ 0.0278 & N/A$^{\dagger}$ & \textit{baseline} \\
2 & $\mathcal{L}_{\text{pixel}}$ + $\mathcal{L}_{\text{adv}}$ & 24.86$^{\star}$ $\pm$ 4.21 & 0.9626 $\pm$ 0.0279 & N/A$^{\dagger}$ & +0.27 \\
3 & + $\mathcal{L}_{\text{perc}}$ (VGG) & 24.55 $\pm$ 4.69 & 0.9630$^{\star}$ $\pm$ 0.0273 & N/A$^{\dagger}$ & -0.04 \\
4 & + $\mathcal{L}_{\text{ctc}}$ & 24.76 $\pm$ 4.71 & 0.9627 $\pm$ 0.0294 & 29.66$^{\star}$ & +0.17 \\
5 & + $\mathcal{L}_{\text{rec-feat}}$ & 24.75 $\pm$ 4.60 & 0.9629 $\pm$ 0.0312 & 30.16 & +0.16 \\
\hline
\multicolumn{6}{|l|}{\footnotesize $^{\dagger}$CER tidak tersedia karena pengenal tidak diaktifkan} \\
\multicolumn{6}{|l|}{\footnotesize $^{\star}$Nilai terbaik untuk metrik ini: PSNR (Eks 2), SSIM (Eks 3), CER (Eks 4)} \\
\hline
\end{tabular}
\end{table*}

% Ablation comparison figure
\begin{figure*}[!t]
\centering
\includegraphics[width=\textwidth]{../data_dukung/ablation_comparison_4panel.pdf}
\caption{Studi ablasi inkremental komponen loss: (a--b) Metrik absolut PSNR dan SSIM untuk setiap konfigurasi, (c--d) Peningkatan inkremental relatif terhadap baseline hanya-piksel. Error bars menunjukkan deviasi standar pada validation set. Loss adversarial memberikan peningkatan PSNR tertinggi (+0.27 dB), sementara loss perceptual mengoptimalkan kesamaan struktural (+0.0016 SSIM).}
\label{fig:ablation_comparison}
\end{figure*}

Integrasi Kehilangan CTC:
Dalam implementasi, kehilangan CTC ($\mathcal{L}_{\text{ctc}}$) di-\textit{backpropagate} ke \textit{generator} bersama komponen kehilangan lainnya, memberikan sinyal kesadaran-teks langsung untuk pelestarian keterbacaan karakter. Formula total kehilangan \textit{generator}:

\begin{equation}
\mathcal{L}_G = 50\mathcal{L}_{\text{pixel}} + 3\mathcal{L}_{\text{adv}} + 8\mathcal{L}_{\text{rec-feat}} + \mathcal{L}_{\text{perc}} + 0.15\mathcal{L}_{\text{ctc}}
\end{equation}

Gradien dari semua komponen mengalir ke \textit{generator} melalui mekanisme \textit{backpropagation} standar.

Analisis Kontribusi Komponen: Kehilangan \textit{adversarial} memberikan peningkatan PSNR terbesar (+0.27 dB), meningkatkan realisme tekstur tanpa mengorbankan stabilitas. Kehilangan perseptual mengoptimalkan kesamaan struktural (SSIM tertinggi 0.9630) dengan \textit{trade-off} penurunan PSNR minor (-0.04 dB). Penambahan CTC mengaktifkan kapabilitas HTR (CER 29.66\%) sambil mempertahankan kualitas visual (PSNR 24.76 dB). Temuan penting: Kehilangan fitur pengenal (Eks 5) menurunkan akurasi HTR (+0.50\% CER) tanpa peningkatan visual, mengindikasikan ketidaksesuaian arsitektural yang dianalisis pada subseksi berikutnya.

Temuan Penting - Perbandingan dengan Pelatihan Produksi: Studi ablasi 15 \textit{epoch} mencapai PSNR maksimal 24.86 dB (Eks 2), jauh di bawah model produksi (30.91 dB). Kesenjangan ini (6.05 dB) mengonfirmasi bahwa pelatihan diperluas ($\geq$40 \textit{epoch}) diperlukan untuk konvergensi penuh. Studi ablasi pada 15 \textit{epoch} efektif untuk mengukur kontribusi relatif setiap komponen, tetapi tidak mencapai batas kinerja absolut.

Rekomendasi Konfigurasi Kehilangan Berdasarkan Kasus Penggunaan:
\begin{enumerate}
    \item Kualitas visual maksimal: Eks 2 (Piksel + \textit{Adversarial}, PSNR 24.86 dB) untuk aplikasi tampilan tanpa persyaratan HTR.
    
    \item Restorasi berorientasi HTR (DIREKOMENDASIKAN): Eks 4 (+ CTC, CER 29.66\%) mencapai keseimbangan kualitas-akurasi tanpa \textit{overhead} kehilangan fitur pengenal.
    
    \item TIDAK direkomendasikan: Eks 5 (\textit{full model}) karena degradasi CER (+0.50\%), \textit{overhead} komputasional signifikan (+15-20\% memori), dan kontribusi \textit{negligible} (analisis detail pada subseksi berikutnya).
\end{enumerate}

Temuan ini memvalidasi arsitektur fungsi loss yang diusulkan dan memberikan panduan berbasis bukti untuk penyetelan \textit{hyperparameter} pada penerapan praktis. Wawasan kunci: Pemilihan fitur tingkat \textit{feature-level} sangat berpengaruh — fitur tugas tambahan harus selaras secara semantik dengan tujuan tugas utama.

\subsubsection{Analisis Mendalam: Ketidakefektifan Kehilangan Fitur Pengenal}
\label{sec:recfeat_analysis}

Meskipun kehilangan fitur pengenal secara konseptual menjanjikan untuk meningkatkan penyelarasan antara citra yang direstorasi dan representasi internal pengenal HTR, hasil studi ablasi menunjukkan bahwa penambahan komponen ini justru menurunkan kinerja HTR (CER meningkat dari 29.66\% menjadi 30.16\%, degradasi +0.50\%). Investigasi mendalam mengungkapkan empat faktor fundamental yang saling terkait.

\textit{Ketidaksesuaian tingkat fitur.} Kehilangan fitur pengenal mengekstrak representasi dari \textit{layer} antara dalam arsitektur pengenal HTR, yang terletak sebelum pemrosesan \textit{Transformer}. Posisi ekstraksi ini menciptakan ketidaksesuaian fundamental antara tujuan optimasi:

\begin{itemize}
    \item Fitur tingkat menengah: Merepresentasikan pola arsitektural konvolusional (tepi, komponen karakter) yang dioptimalkan untuk persiapan masukan \textit{Transformer}, bukan untuk persepsi visual manusia atau kualitas semantik tingkat tinggi.
    
    \item Kebutuhan restorasi dokumen: Memerlukan optimasi pada kualitas semantik tingkat tinggi (kejelasan teks, keterbacaan) dan estetika visual (latar belakang halus, tekstur alami, penekanan artefak).
    
    \item Konflik tingkat representasi: Kehilangan kuadrat rata-rata pada fitur tingkat menengah tidak selaras dengan metrik kualitas perseptual tingkat tinggi maupun rekonstruksi tingkat piksel.
\end{itemize}

\textit{Ketidakseimbangan magnitudo kehilangan.} Analisis \textit{log} pelatihan pada \textit{epoch} 15 mengungkapkan disparitas magnitudo yang ekstrem antar komponen kehilangan:

\begin{table}[!t]
\centering
\footnotesize
\caption{Perbandingan Magnitudo Komponen Kehilangan: analisis disparitas magnitudo untuk memvalidasi konfigurasi pembobotan. Menunjukkan nilai mentah, bobot yang diterapkan, dan kontribusi tertimbang untuk setiap komponen terhadap total kehilangan \textit{generator}.}
\label{table_loss_magnitude}
\begin{tabular}{lrrc}
\hline
\textbf{Komponen} & \textbf{Nilai Mentah} & \textbf{Bobot} & \textbf{Tertimbang} \\
\hline
CTC (dominan) & 365.00 & 0.15 & 54.75 \\
Perseptual (utama) & 41.94 & 1.00 & 41.94 \\
\textit{Adversarial} & 0.80 & 3.00 & 2.40 \\
Piksel & 0.16 & 50.00 & 8.00 \\
Fitur Pengenal & 0.06$^{\ddagger}$ & 8.00 & 0.48$^{\ddagger}$ \\
\hline
\multicolumn{3}{l}{Total Kehilangan \textit{Generator}} & 107.57 \\
\multicolumn{3}{l}{Proporsi Fitur Pengenal} & 0.45\%$^{\ddagger}$ \\
\multicolumn{4}{l}{\footnotesize $^{\ddagger}$Kontribusi terkecil (hanya 0.45\% dari total kehilangan \textit{generator})} \\
\hline
\end{tabular}
\end{table}

Observasi utama:
\begin{itemize}
    \item Nilai mentah fitur pengenal ($\approx$0.06) adalah 2-3 urutan magnitudo lebih kecil dari CTC ($\approx$365) dan perseptual ($\approx$42).
    \item Bahkan dengan pembobotan tinggi, kontribusi efektif fitur pengenal dapat diabaikan dalam optimasi.
    \item Aliran gradien: Komponen CTC dan perseptual mendominasi dengan kekuatan $\approx$100-200$\times$ lebih besar, secara efektif menutupi gradien fitur pengenal dalam optimasi.
\end{itemize}

Implikasi: Fitur pengenal memiliki pengaruh yang dapat diabaikan pada pembaruan gradien \textit{generator}, sehingga tidak dapat memberikan panduan yang bermakna untuk trajektori optimasi.

\textit{Interferensi gradien.} Kehilangan perseptual dan kehilangan fitur pengenal mengoptimalkan terhadap ruang fitur yang berbeda dengan tujuan yang berpotensi bertentangan:

\begin{itemize}
    \item Kehilangan perseptual: Mendorong \textit{generator} ke arah estetika halus yang selaras dengan persepsi manusia (gradien mengarah pada tekstur alami, tepi halus, penekanan artefak).
    
    \item Kehilangan fitur pengenal: Mendorong \textit{generator} mencocokkan pola karakter tajam yang dioptimalkan untuk pengenalan (gradien mengarah pada peningkatan ketajaman tepi untuk pemrosesan \textit{Transformer}).
\end{itemize}

Hasil empiris: eksperimen 5 (dengan fitur pengenal) mencapai CER 30.16\% dibandingkan eksperimen 4 (tanpa fitur pengenal) dengan CER 29.66\%, mengindikasikan bahwa gradien yang berlawanan mengganggu konvergensi optimal untuk kinerja HTR. Meskipun eksperimen 5 menambah komponen kehilangan tambahan, peningkatan CER menunjukkan fitur pengenal mengoptimalkan tujuan yang tidak selaras dengan kinerja HTR \textit{downstream}.

\textit{Beban komputasional tanpa manfaat.} Penggunaan kehilangan fitur pengenal memerlukan modifikasi arsitektur pengenal untuk menghasilkan keluaran ganda, meningkatkan jejak memori sekitar 15--20\% dan waktu pelatihan 5--10\%.

Analisis biaya-manfaat menunjukkan:
\begin{itemize}
    \item Biaya: Peningkatan memori 15--20\%, peningkatan waktu pelatihan 5--10\%
    \item Manfaat: Degradasi CER +0.50\% (dari 29.66\% ke 30.16\%), penurunan PSNR marjinal (-0.01 dB), peningkatan SSIM dapat diabaikan (+0.0002)
    \item Kesimpulan: Rasio biaya-manfaat negatif, beban komputasional tidak memberikan peningkatan kinerja yang bermakna
\end{itemize}

Kesimpulan: Analisis mendalam mengungkapkan bahwa kehilangan fitur pengenal secara fundamental tidak selaras dengan tujuan restorasi dokumen karena ketidaksesuaian tingkat fitur, magnitudo dapat diabaikan (0.45\% dari total kehilangan), potensi interferensi gradien, dan \textit{overhead} komputasional tanpa manfaat.

Rekomendasi: Konfigurasi eksperimen 4 (Piksel + \textit{Adversarial} + Perseptual + CTC) tanpa kehilangan fitur pengenal mencapai kinerja optimal. Kehilangan CTC sudah cukup untuk memberikan sinyal kesadaran-teks kepada \textit{generator} tanpa memerlukan penyelarasan fitur tambahan.

Kontribusi ilmiah: Temuan ini menunjukkan bahwa pemilihan tingkat fitur sangat berpengaruh dalam optimasi multi-objektif untuk restorasi dokumen. Fitur tugas tambahan harus selaras secara semantik dengan tujuan tugas utama, bukan hanya secara konseptual terkait.

\subsubsection{Analisis Konfigurasi Optimal: Mengapa Eksperimen 04 Merupakan Pilihan Terbaik}
\label{sec:exp04_optimal}

Berdasarkan hasil studi ablasi komprehensif, penelitian ini mengidentifikasi eksperimen 04 (konfigurasi Piksel + \textit{Adversarial} + Perseptual + CTC \textit{tanpa} fitur pengenal) sebagai konfigurasi optimal untuk penerapan produksi. Analisis mengungkapkan empat faktor kunci: keseimbangan multi-objektif, sinergi komponen kehilangan, efisiensi komputasional, dan cakupan informasi multi-resolusi.

\textit{Optimalitas Pareto multi-objektif.} eksperimen 04 mencapai keseimbangan optimal dalam ruang tri-objektif (kualitas visual, kesamaan struktural, kinerja HTR). Referensi Tabel~\ref{table_ablation_loss} menunjukkan eksperimen 04 memiliki: CER terbaik (29.66\%), PSNR kompetitif (24.76 dB, hanya -0.10 dB dari maksimum), dan SSIM kompetitif (0.9627, hanya -0.0003 dari maksimum).

Observasi utama:
\begin{itemize}
    \item Kinerja HTR terbaik: CER 29.66\% adalah terendah di antara semua konfigurasi dengan kemampuan HTR
    
    \item Kualitas visual kompetitif: PSNR 24.76 dB hanya -0.10 dB dari maksimum (Eks 2: 24.86 dB), dan SSIM 0.9627 hanya -0.0003 dari maksimum (Eks 3: 0.9630)
    
    \item Optimalitas Pareto: Tidak ada konfigurasi lain yang secara simultan meningkatkan \textit{semua} metrik. eksperimen 04 berada pada \textit{frontier} Pareto untuk restorasi berorientasi HTR dengan CER terbaik sambil mempertahankan kualitas visual mendekati optimal
\end{itemize}

\textit{Sinergi komponen kehilangan.} Setiap komponen kehilangan memberikan kontribusi unik: \textit{Adversarial} (PSNR, realisme tekstur), Perseptual (SSIM puncak, pelestarian struktur), CTC (kapabilitas HTR, keterbacaan karakter). Keempat komponen selaras tanpa konflik — semua mendukung tujuan ``merestorasi dokumen yang baik secara visual dan dapat dibaca HTR'' dengan efek sinergis pada hierarki multi-resolusi (piksel $\rightarrow$ struktur $\rightarrow$ semantik).

\textit{Efisiensi komputasional.} Eksperimen 04 superior secara komputasional dibandingkan Eksperimen 05:

\begin{table}[!t]
\centering
\footnotesize
\caption{Perbandingan Komputasional: Eksperimen 04 vs Eksperimen 05}
\label{table_exp04_computational}
\begin{tabular}{lcc}
\hline
\textbf{Metrik} & \textbf{Eks 04} & \textbf{Eks 05} \\
\hline
Memori/\textit{Batch} & \textit{Baseline} & +15--20\% \\
Waktu Pelatihan & \textit{Baseline} & +5--10\% lebih lambat \\
Kinerja CER & 29.66\% & 30.16\% \\
\hline
\end{tabular}
\end{table}

Analisis biaya-manfaat: eksperimen 05 menambah peningkatan memori 15--20\% dan waktu pelatihan 5--10\% lebih lambat, namun CER malah 0.50\% lebih buruk (dari 29.66\% ke 30.16\%). Perubahan kualitas visual sangat kecil: penurunan PSNR -0.01 dB dan peningkatan SSIM +0.0002 (keduanya dapat diabaikan). Rasio biaya-manfaat negatif.

\textit{Cakupan informasi multi-resolusi.} eksperimen 04 menggunakan sumber informasi optimal dengan cakupan: Piksel (rekonstruksi tingkat rendah) $\rightarrow$ Perseptual (struktur tingkat menengah) $\rightarrow$ CTC (semantik tingkat tinggi) + \textit{Adversarial} (realisme global). Fitur pengenal (fitur CNN pra-\textit{Transformer}) redundan karena berada antara Piksel dan Perseptual, sudah tercakup oleh kombinasi keduanya.

Kesimpulan: eksperimen 04 adalah konfigurasi optimal dengan justifikasi: (1) HTR terbaik (CER 29.66\% pada 15 \textit{epoch} ablasi), (2) Optimalitas Pareto, (3) Sinergi komponen, (4) Efisiensi komputasional. Model produksi menggunakan bobot kehilangan eksperimen 04 (tanpa kehilangan fitur pengenal) dengan pelatihan penuh 50 \textit{epoch} mencapai PSNR 30.74 dB (uji) atau 30.91 dB (validasi) dan CER 34.9\% (uji) atau 27.1\% (validasi). Temuan penting: Kesederhanaan dengan penyelarasan yang tepat mengungguli kompleksitas (4 kehilangan $>$ 5 kehilangan).

\subsubsection{Justifikasi Strategi Integrasi Pengenal HTR}
\label{sec:ablation}

Pemilihan strategi pengenal beku dalam penelitian ini didasarkan pada alasan desain yang mempertimbangkan efisiensi komputasional, stabilitas pelatihan, dan panduan dari literatur terkait. Tabel~\ref{table_ablation_recognizer_strategy} menyajikan perbandingan analitik dari strategi integrasi alternatif.

% Recognizer integration strategy comparison
\begin{table}[!t]
\renewcommand{\arraystretch}{1.2}
\caption{Analisis Komparatif Strategi Integrasi Pengenal HTR}
\label{table_ablation_recognizer_strategy}
\centering
\small
\begin{tabular}{|p{2.5cm}|p{2.2cm}|p{2.5cm}|}
\hline
\textbf{Strategi} & \textbf{Keuntungan} & \textbf{Limitasi} \\
\hline
Pelatihan bersama & Optimasi \textit{end-to-end}, adaptasi spesifik-tugas & Biaya komputasi sangat tinggi, ketidakstabilan pelatihan dari konflik gradien, risiko pelupaan katastrofik~\cite{kirkpatrick2017overcoming} \\
\hline
Beku$^\star$ & Konvergensi stabil (48 GPU-jam untuk 50 \textit{epoch}), efisiensi komputasional, melestarikan kinerja pengenal (CER 33.72\%) & Adaptasi spesifik-tugas pengenal terbatas \\
\hline
Dua tahap & Optimasi seimbang, stabilitas sekuensial & Waktu pelatihan berlipat ganda, penjadwalan kompleks, memerlukan infrastruktur pelatihan pengenal \\
\hline
\textit{Baseline} (Tanpa R) & Arsitektur sederhana, komputasi minimal & Tidak ada optimasi sadar-HTR, CER sangat tinggi pada masukan terdegradasi (83.4\%) \\
\hline
\multicolumn{3}{|l|}{\footnotesize $^\star$Strategi yang dipilih dalam penelitian ini} \\
\hline
\end{tabular}
\end{table}

Justifikasi Pemilihan Strategi Pengenal Beku:

1. Efisiensi komputasional: Pengenal pra-latih (27.86M parameter) tidak memerlukan komputasi gradien, sehingga hanya \textit{generator} (89.4M parameter) dan diskriminator (18.9M parameter) yang dioptimasi. Pendekatan ini mengurangi kompleksitas optimasi dan mencegah konflik gradien antar komponen. Waktu pelatihan aktual yang tercatat adalah 48 GPU-jam untuk 50 \textit{epoch} dengan ukuran \textit{batch} 2 pada GPU RTX A4000 16GB, lebih efisien dibandingkan strategi pelatihan bersama yang memerlukan optimasi simultan seluruh komponen.

2. Stabilitas pelatihan: Pengenal beku mencegah gradien yang bertentangan antara tujuan rekonstruksi ($\mathcal{L}_{\text{pixel}}$, $\mathcal{L}_{\text{adv}}$) dan tujuan pengenalan ($\mathcal{L}_{\text{ctc}}$, $\mathcal{L}_{\text{rec-feat}}$). Literatur menunjukkan pelatihan bersama pada GAN multi-tugas rentan terhadap kolaps modus dan konvergensi berosilasi~\cite{souibgui2022,kang2021pay}.

3. Pelestaran kinerja: Pengenal hibrida CNN-\textit{Transformer} (6-layer transformer encoder dengan CTC decoder) yang dilatih pada dataset ANRI mencapai CER 33.72\% saat dievaluasi pada subset ANRI terpisah (n=948) untuk validasi awal kemampuan generalisasi pada domain paleografi Belanda historis abad 16-18. Membekukan pengenal mempertahankan kapabilitas generalisasi ini dan mencegah pelupaan katastrofik~\cite{kirkpatrick2017overcoming} yang dapat terjadi jika pengenal dilatih ulang pada pergeseran distribusi dari citra yang ditingkatkan. Kerangka kerja yang diusulkan mencapai CER 27.11\% pada \textit{set} validasi, menunjukkan bahwa restorasi berbasis GAN meningkatkan keterbacaan melampaui kemampuan intrinsik pengenal.

4. Diferensiasi dari literatur: Berbeda dari Souibgui \textit{et al.}~\cite{souibgui2022} yang menggunakan pelatihan bersama pengenal-generator (berpotensi ketidakstabilan gradien), penelitian ini membekukan pengenal praterlatih untuk stabilitas optimisasi. Strategi ini mencegah konflik gradien yang terdokumentasi pada GAN multi-tugas~\cite{souibgui2022,kang2021pay}, memungkinkan generator fokus pada restorasi visual tanpa mendegradasi pengenal.

Bukti Empiris - Indikasi Awal \textit{Baseline}: Pengujian eksploratori menggunakan \textit{baseline} U-Net (Strategi: Tanpa Pengenal, hanya kehilangan L1+SSIM) pada subset data validasi (n=160 sampel, 5 \textit{epoch} pelatihan cepat\footnote{Terbatas pada 5 \textit{epoch} karena kendala komputasional. Pelatihan penuh 50 \textit{epoch} diharapkan meningkatkan metrik tetapi tidak mungkin menutup kesenjangan CER vs pendekatan sadar-HTR berdasarkan literatur~\cite{souibgui2022,kang2021pay}.}) menunjukkan:
\begin{itemize}
    \item Kualitas visual (kurang terlatih): PSNR 13.75 $\pm$ 2.41 dB, SSIM 0.817 $\pm$ 0.075 (pola konvergensi \textit{baseline})
    \item Pelatihan penuh yang diharapkan: Literatur~\cite{souibgui2022} menyarankan kehilangan hanya-piksel konvergen ke SSIM $\approx$0.96-0.97 setelah 50 \textit{epoch}, tetapi CER tetap 3-5\% lebih tinggi vs pendekatan sadar-HTR
    \item Catatan: Eksperimen ini dirancang sebagai \textit{proof-of-concept} protokol pelatihan, bukan perbandingan terkontrol penuh. Namun, konsisten dengan temuan literatur~\cite{souibgui2022,kang2021pay} bahwa komponen kehilangan sadar-HTR ($\mathcal{L}_{\text{ctc}}$, $\mathcal{L}_{\text{rec-feat}}$) memberikan kontribusi pada pelestarian teks
\end{itemize}

Validasi H2\textsubscript{B}: Hipotesis bahwa pendekatan pengenal beku menghasilkan stabilitas pelatihan yang memadai TERVALIDASI melalui: (1) Tidak ada kejadian divergensi selama pelatihan 50 \textit{epoch} (kurva kehilangan menurun mulus pada Gambar~\ref{fig:convergence_analysis}), (2) Trajektori metrik validasi konsisten tanpa lonjakan tidak menentu, (3) Kesesuaian dengan praktik literatur~\cite{souibgui2022}. Catatan: Validasi ini mengukur stabilitas konvergensi, bukan superioritas komparatif terhadap pelatihan bersama yang berada di luar ruang lingkup penelitian. Pendekatan ini memungkinkan restorasi mencapai keterbacaan mendekati batas atas teoretis (gap CER hanya 0.8 poin persentase dari GT bersih pada set uji).

\subsection{Analisis Kuantitatif Kasus Kegagalan}

Evaluasi komprehensif dilakukan pada set uji yang diisolasi secara ketat ($n = 712$ citra, 15\% dari total set data) mengikuti protokol akademis standar dengan pemisahan ketat antara set pelatihan dan set uji.

Kinerja Keseluruhan pada Set Uji:
\begin{itemize}
    \item CER: $34.9\% \pm 21.8\%$ (interval kepercayaan 95\%: [33.3\%, 36.5\%])
    \item WER: $82.4\% \pm 26.1\%$ (interval kepercayaan 95\%: [80.5\%, 84.3\%])
    \item PSNR: $30.74 \pm 5.09$ dB (interval kepercayaan 95\%: [30.36, 31.11] dB)
    \item SSIM: $0.987 \pm 0.014$ (interval kepercayaan 95\%: [0.986, 0.988])
\end{itemize}

Target kualitas visual tercapai dengan PSNR $>$ 30 dB dan SSIM $>$ 0.95, menunjukkan peningkatan signifikan dalam keterbacaan dokumen terdegradasi.

Kategorisasi Pola Kegagalan Kualitatif\protect\footnotemark:
Analisis sistematis terhadap 712 sampel uji mengidentifikasi empat pola kegagalan utama berdasarkan inspeksi prediksi pengenal teks dan karakteristik kesalahan:
\begin{enumerate}
    \item Teks memudar ekstrem (3 kasus, 0.4\%): Metode yang diusulkan gagal merekonstruksi goresan dengan intensitas sangat rendah, menghasilkan karakter yang hilang atau terdistorsi. Kasus representatif menunjukkan prediksi tunggal (``1'' untuk sekuens 40+ karakter, ``ei'' untuk teks panjang) dengan CER mendekati 100\%.
    
    \item Teks bertumpang tindih dengan stempel\textsuperscript{\textdagger}: Pemisahan antara teks utama dan elemen bertindih (stempel, anotasi) tidak sempurna, menyebabkan penggabungan karakter dari kedua lapisan. Pola ini teridentifikasi pada evaluasi kualitatif dokumen ANRI (subsection evaluasi dokumen historis nyata), bukan pada set uji sintetis.
    
    \item Ligatur paleografi kompleks (448 kasus, 62.9\%): Goresan terhubung yang khas dalam tulisan tangan abad ke-16 hingga ke-18 disederhanakan atau salah diinterpretasi. Contoh terverifikasi: ``verleden'' diprediksi sebagai ``verleeden'', ``jar'' sebagai ``saar'' (CER 17.6\%, sampel uji \#9). Pola ini merupakan sumber kegagalan dominan pada set data paleografi.
    
    \item Artefak numerik dan simbol (13 kasus, 1.8\%): Angka dan simbol khusus (ƒ, §, =, titik-titik) sering kali tidak terestorasi dengan baik. Contoh: simbol paragraf ``§ 73:'' diprediksi sebagai ``537201'' (CER 100\%), menunjukkan keterbatasan metode yang diusulkan pada karakter non-alfabet.

\end{enumerate}
\footnotetext{Kategorisasi dilakukan melalui analisis prediksi pengenal teks pada 712 sampel uji dengan kriteria: (1) Teks memudar: prediksi kosong atau sangat pendek untuk teks panjang, (2) Ligatur: kesalahan pada kata dengan goresan terhubung kompleks dan CER$>$15\%, (3) Numerik/simbol: CER$>$30\% pada teks mengandung karakter non-alfabet. Data lengkap tersedia dalam berkas hasil evaluasi.}
\footnotetext{\textsuperscript{\textdagger}Set data sintetis tidak mengandung stempel bertumpang tindih. Pola ini teridentifikasi pada evaluasi kualitatif dokumen historis ANRI ($n=15$) di mana teks utama bertumpang tindih dengan stempel arsip atau anotasi marginal.}

Analisis distribusi (Tabel~\ref{table_failure_distribution}) mengungkapkan: meskipun 75\% sampel menunjukkan CER$>$15\%, mayoritas (62.9\%) adalah ligatur paleografi dengan CER moderat (35.2\%), bukan kegagalan total. Kasus kegagalan parah (CER$\approx$100\%) bersifat jarang ($<$2\%) dan terkait degradasi ekstrem atau karakter non-standar, mengindikasikan bahwa sebagian besar kesalahan bersumber dari kompleksitas intrinsik tulisan tangan paleografi, bukan kegagalan restorasi.

Korelasi dengan Citra Bersih:
Delta CER antara citra yang direstorasi (CER rata-rata 35.0\%) dan citra bersih tanpa degradasi (34.2\%, lihat Tabel~\ref{table_synthetic_results}) adalah $+0.8$ poin persentase (interval kepercayaan 95\%: [+0.6, +1.0]), mengindikasikan bahwa proses restorasi mempertahankan sebagian besar keterbacaan dokumen asli sambil menghilangkan degradasi visual. Korelasi Pearson yang sangat kuat ($r = 0.96$, $p<0.001$, $n=712$) antara CER pada citra bersih dan yang direstorasi menunjukkan konsistensi kinerja pengenal teks di seluruh set data, memvalidasi bahwa metode yang diusulkan melestarikan karakteristik kesulitan intrinsik setiap sampel.

% [Table: Distribusi Pola Kegagalan]
\begin{table}[!t]
\renewcommand{\arraystretch}{1.15}
\caption{Distribusi Pola Kegagalan pada Set Uji ($n=712$)}
\label{table_failure_distribution}
\centering
\footnotesize
\begin{tabularx}{\columnwidth}{|X|c|c|c|}
\hline
\textbf{Pola Kegagalan} & \textbf{Jml} & \textbf{\%} & \textbf{CER} \\
\hline
Ligatur paleografi & 448 & 62.9 & 35.2 \\
Lainnya (CER$>$50\%) & 67 & 9.4 & 78.3 \\
Numerik \& simbol & 13 & 1.8 & 92.1 \\
Teks memudar ekstrem & 3 & 0.4 & 100.0 \\
Prediksi sangat pendek & 3 & 0.4 & 97.5 \\
\hline
\textbf{Total CER Tinggi} & \textbf{534} & \textbf{75.0} & \textbf{42.8} \\
\textbf{($>$15\%)} & & & \\
Normal/Rendah ($<$15\%) & 178 & 25.0 & 6.2 \\
\hline
\multicolumn{4}{|X|}{\scriptsize Kategorisasi berdasarkan karakteristik kesalahan HTR. CER dalam \%.} \\
\hline
\end{tabularx}
\end{table}

Tabel~\ref{table_failure_distribution} dan Gambar~\ref{fig:failure_distribution} menunjukkan distribusi kuantitatif pola kegagalan, mengonfirmasi bahwa ligatur paleografi merupakan tantangan dominan pada set data tulisan tangan historis (62.9\% dari kasus CER tinggi), sementara kegagalan ekstrem (teks memudar, artefak simbol) bersifat jarang tetapi parah dengan CER rata-rata mendekati 100\%.

% [Figure: Distribusi Pola Kegagalan - Bar Charts]
\begin{figure*}[!t]
\centering
\includegraphics[width=0.95\textwidth]{../data_dukung/fig_failure_distribution.pdf}
\caption{Distribusi pola kegagalan (kiri) dan tingkat keparahan (kanan) pada set uji ($n=712$). Ligatur paleografi mendominasi (62.9\%) tetapi dengan CER moderat (35.2\%), sementara kasus jarang (teks memudar, artefak simbol) menunjukkan kegagalan parah (CER$\approx$100\%), mengindikasikan bahwa mayoritas kesalahan bersumber dari kompleksitas tulisan tangan paleografi, bukan kegagalan restorasi total.}
\label{fig:failure_distribution}
\end{figure*}

Analisis kualitatif menunjukkan bahwa kasus kegagalan yang teridentifikasi mencerminkan tantangan pelestarian detail goresan halus pada tulisan tangan abad ke-16 hingga ke-18, seperti ligatur ``verleden''$\rightarrow$``verleeden'' dan ``jar''$\rightarrow$``saar'' (sampel \#9) yang menunjukkan kesulitan dalam pengenalan karakter paleografi kompleks.

\subsection{Validasi Hipotesis Penelitian}

Berdasarkan hasil eksperimental pada Tabel~\ref{table_synthetic_results}, pengujian hipotesis dilakukan menggunakan \textit{independent samples t-test} untuk H1 (perbandingan utama: hasil yang direstorasi vs masukan terdegradasi) dan studi ablasi sistematis untuk H2\textsubscript{A}-H2\textsubscript{C}, dengan koreksi Bonferroni ($\alpha = 0.0167$) untuk ketiga hipotesis spesifik.

% [Table 8: Hasil Validasi Hipotesis]
\begin{table*}[!t]
\renewcommand{\arraystretch}{1.2}
\caption{Hasil Pengujian Hipotesis pada Himpunan Uji ($n=712$, $\alpha=0.05$ untuk H1, $\alpha=0.0167$ untuk H2)\protect\footnotemark}
\label{table_hypothesis_test}
\centering
\scriptsize
\begin{tabular}{|p{1.5cm}|p{1.8cm}|p{1.3cm}|p{1.3cm}|p{1cm}|p{0.9cm}|p{1.2cm}|}
\hline
\textbf{Hipotesis} & \textbf{Metrik} & \textbf{Pembanding} & \textbf{Usulan} & \textbf{Nilai-p} & \textbf{Cohen's d} & \textbf{Hasil} \\
\hline
H1 (Utama) & CER & 83.4\% & 34.9\%$^\dagger$ & $<$0.001 & 0.85 & Diterima$^\dagger$ \\
& Pengurangan Absolut & - & 48.5 poin & - & - & - \\
& Pengurangan Relatif & - & 58.2\% & - & - & - \\
\hline
H2\textsubscript{A} & CER (Diskriminator ganda vs tunggal) & 27.12\% & 27.11\% & $>$0.05 & 0.004 & Marjinal$^\ddagger$ \\
& PSNR (Diskriminator ganda vs tunggal) & 30.63 dB & 30.91 dB & $>$0.05 & 0.065 & Marjinal$^\ddagger$ \\
\hline
H2\textsubscript{B} & Stabilitas Konvergensi & \textit{Baseline} & Stabil & - & - & Terkonfirmasi$^\S$ \\
& Varians Kehilangan (Epoch 40-50) & - & Rendah & $<$0.01 & - & Terkonfirmasi \\
\hline
H2\textsubscript{C} & PSNR (Kualitas Visual) & - & 30.74 dB & - & - & Tinggi \\
& SSIM (Kesamaan Struktural) & - & 0.987 & - & - & Tinggi \\
& Kesenjangan CER vs GT Bersih & 34.1\% & 34.9\% & - & 0.8 poin & Minimal \\
\hline
\multicolumn{7}{|l|}{\footnotesize $^\dagger$H1 diterima: CER 83.4\% $\rightarrow$ 34.9\% (himpunan uji, $n$=712, $p<$0.001, $d=0.85$). Himpunan validasi (n=710): CER 27.1\%. \textit{Test set} lebih konservatif.} \\
\multicolumn{7}{|l|}{\footnotesize $^\ddagger$H2A: Peningkatan marjinal tidak signifikan ($p>$0.05, berdasarkan studi ablasi diskriminator Tabel~\ref{table_ablation_discriminator})} \\
\multicolumn{7}{|l|}{\footnotesize $^\S$H2B: Mengukur stabilitas konvergensi \textit{frozen recognizer}, bukan perbandingan dengan \textit{joint training} (di luar ruang lingkup)} \\
\hline
\end{tabular}
\end{table*}
\footnotetext{Data terverifikasi dari evaluasi himpunan uji ($n=712$): H1 menggunakan CER pembanding dan usulan, H2\textsubscript{C} menggunakan PSNR, SSIM, dan CER rujukan bersih.}

Analisis Hasil Pengujian:

Hipotesis Utama (H1): DITERIMA dengan bukti empiris kuat.
\begin{itemize}
\item Pengurangan CER: 83.4\% $\rightarrow$ 34.9\% (pengurangan absolut 48.5 poin, relatif 58.2\%)
\item Signifikansi statistik: $p<0.001$, Cohen's $d=0.85$ (efek besar), IK 95\%: [33.12\%, 36.72\%]
\item Himpunan validasi: CER 27.1\% (n=710), menunjukkan performa konsisten lintas partisi data
\item Kesenjangan dengan GT bersih (34.1\%): 0.8 poin persentase, mendekati batas teoretis \textit{recognizer}
\end{itemize}

Catatan: H1 mengukur efektivitas restorasi dengan membandingkan keluaran model terhadap masukan terdegradasi (\textit{baseline}), bukan perbandingan dengan metode restorasi lain yang dilatih pada \textit{dataset} berbeda.

Hipotesis Spesifik: Status validasi berdasarkan data ablasi.
\begin{itemize}
\item \textbf{H2\textsubscript{A} (Diskriminator \textit{Dual-Modal}):} TIDAK TERKONFIRMASI SECARA STATISTIK - Berdasarkan studi ablasi diskriminator (Tabel~VI): arsitektur diskriminator \textit{dual-modal} menunjukkan peningkatan marjinal dibanding diskriminator tunggal ($\Delta$PSNR +0.28 dB, $\Delta$CER $-0.01$\%), namun tidak signifikan secara statistik ($p>0.05$). Meskipun menunjukkan tren positif konsisten, efek sangat kecil (Cohen's $d=0.004$ untuk CER).

\item \textbf{H2\textsubscript{B} (Stabilitas Konvergensi \textit{Frozen Recognizer}):} TERKONFIRMASI - Metode \textit{frozen recognizer} menghasilkan konvergensi stabil tanpa divergensi (varians kehilangan epoch 40-50 rendah, $p<0.01$), kurva fungsi loss menunjukkan penurunan konsisten tanpa fluktuasi signifikan. Catatan: Hipotesis ini mengukur stabilitas pelatihan, bukan perbandingan superioritas terhadap \textit{joint training} yang berada di luar ruang lingkup penelitian.

\item \textbf{H2\textsubscript{C} (Pelestarian Kualitas Visual):} TERKONFIRMASI - Pada himpunan uji: PSNR 30.74 dB dan SSIM 0.987 menunjukkan kualitas visual tinggi, CER 34.9\% mendekati rujukan bersih (34.1\%, kesenjangan hanya 0.8 poin), mengonfirmasi tidak ada \textit{trade-off} antara kualitas visual dan keterbacaan HTR. Himpunan validasi menunjukkan PSNR 30.91 dB dan CER 27.1\%, konsisten dengan performa \textit{test set}.
\end{itemize}

Kesimpulan Validasi Hipotesis: H1 DITERIMA dengan pengurangan CER 58.2\% ($p<0.001$, $d=0.85$). Untuk hipotesis komponen: H2\textsubscript{A} tidak terkonfirmasi (\textit{dual-modal} marjinal, $p>0.05$), H2\textsubscript{B} terkonfirmasi (\textit{frozen recognizer} stabil), H2\textsubscript{C} terkonfirmasi (kualitas visual tinggi tanpa \textit{trade-off} keterbacaan). Temuan mengonfirmasi keberhasilan berasal dari \textit{frozen recognizer} dan optimasi kehilangan, bukan kompleksitas diskriminator.

% ============================================================
% VI. DISKUSI
% ============================================================
\section{Diskusi}

Bagian ini membahas temuan kunci, implikasi teoretis dan praktis, keterbatasan, serta arah penelitian masa depan.

\subsection{Kontribusi Arsitektur \textit{Dual-Modal} dan Wawasan Protokol Pelatihan}

Studi ablasi diskriminator (Tabel~\ref{table_ablation_discriminator}) mengungkap temuan penting tentang dinamika pelatihan GAN berorientasi HTR. Meskipun arsitektur \textit{dual-modal} secara teoretis lebih kuat, hasil empiris menunjukkan peningkatan yang sangat minimal dibanding diskriminator CNN tunggal ($\Delta$PSNR +0.28 dB, $\Delta$CER -0.01\%, tidak signifikan secara statistik dengan $p > 0.05$). 

Analisis mendalam mengidentifikasi bahwa kontribusi marginal arsitektur \textit{dual-modal} dalam penelitian ini terutama disebabkan oleh \textit{limitation implementasi}, bukan limitation inherent arsitektur itu sendiri. Tiga faktor kunci dari pilihan desain penelitian yang membatasi efektivitas \textit{dual-modal}:

\textit{Modus teks prediksi:} Implementasi penelitian ini menggunakan teks prediksi sebagai masukan diskriminator, sehingga cabang LSTM menerima teks dari keluaran \textit{generator} (CER 27.11\%) bukan teks \textit{ground truth}. Masukan teks yang mengandung derau (tingkat kesalahan 27\%) membatasi kemampuan LSTM untuk mempelajari pola teks yang akurat, dan mekanisme atensi \textit{cross-modal} tidak dapat sepenuhnya memanfaatkan semantik teks yang bersih. Dengan modifikasi menggunakan supervisi teks \textit{ground truth}, potensi arsitektur \textit{dual-modal} dapat lebih terealisasi.

\textit{Ketidakseimbangan bobot kehilangan:} Kontribusi kehilangan \textit{adversarial} hanya sekitar 5\% dari total sinyal pelatihan (bobot 3.0 dibanding kehilangan piksel 50.0 dan fitur pengenalan 8.0). Perubahan arsitektur diskriminator hanya memengaruhi komponen minoritas, sementara 95\% gradien berasal dari komponen kehilangan yang identik untuk kedua eksperimen. Untuk memperbesar pengaruh diskriminator, bobot \textit{adversarial} perlu ditingkatkan (misalnya, dari 3.0 menjadi 10.0).

\textit{Hambatan kapasitas generator:} Kedua eksperimen menggunakan \textit{generator} yang identik. Batas kualitas keluaran ditentukan oleh arsitektur \textit{generator}, bukan sophistication umpan balik diskriminator. \textit{Generator} dengan kapasitas terbatas tidak dapat sepenuhnya memanfaatkan diskriminator yang lebih kompleks.

Arsitektur \textit{dual-modal} tetap berharga secara teoretis karena dapat: (a) memodelkan koherensi teks sekuensial secara eksplisit melalui LSTM, (b) mendeteksi masalah konektivitas karakter yang halus, (c) menyediakan sinyal evaluasi multi-modal. Untuk merealisasikan potensi penuh dalam penelitian masa depan, diperlukan modifikasi protokol: supervisi teks \textit{ground truth}, bobot \textit{adversarial} yang lebih tinggi (misalnya dari 3.0 menjadi 10.0), atau arsitektur \textit{generator} yang sadar-teks.

Temuan ini memberikan wawasan bahwa dalam optimisasi GAN multi-objektif, inovasi arsitektural harus disertai dengan protokol pelatihan yang selaras untuk mencapai manfaat yang diinginkan. Kontribusi marginal \textit{dual-modal} dalam penelitian ini mencerminkan ketidakselarasan implementasi, bukan kegagalan konseptual arsitektur. Kompleksitas diskriminator saja tidak menjamin peningkatan jika konfigurasi pelatihan tidak dirancang untuk memanfaatkannya secara optimal.

Analisis kualitatif pada kasus sukses mengungkapkan bahwa diskriminator \textit{dual-modal} membantu menjaga konektivitas karakter dan kontinuitas goresan pada tulisan tangan \textit{kursif}. Namun, untuk degradasi yang menantang, kedua varian diskriminator menunjukkan performa yang sebanding, mengonfirmasi bahwa kapasitas \textit{generator} merupakan faktor pembatas utama.

\subsection{Strategi Pengenal Beku dan Stabilitas Pelatihan}

Pemilihan pengenal praterlatih yang dibekukan sebagai pengganti pelatihan bersama (sebagaimana diusulkan dalam penelitian sebelumnya~\cite{souibgui2021enhance}) terbukti memberikan manfaat signifikan. Pelatihan bersama memperkenalkan tantangan optimisasi fundamental: pengenal dan \textit{generator} dapat berevolusi ke arah yang bertentangan, dan degradasi pengenal selama pelatihan menghasilkan gradien yang tidak stabil.

Pendekatan pengenal beku memberikan umpan balik sadar-teks yang konsisten dan andal selama pelatihan. Dengan melatih pengenal sebelumnya pada \textit{dataset} tulisan tangan bersih yang beragam, representasi fitur teks yang kuat telah dikodekan dan dapat memandu restorasi tanpa memerlukan label teks selama pelatihan GAN. Stabilitas ini dikonfirmasi melalui trajektori konvergensi yang halus tanpa osilasi signifikan (Gambar~\ref{fig:convergence_analysis}).

\subsection{Wawasan Optimisasi Bobot Kehilangan}

Pencarian \textit{grid search} pada \textit{hyperparameter} mengungkapkan bahwa bobot kehilangan \textit{adversarial} dan kehilangan perseptual merupakan \textit{hyperparameter} paling berpengaruh untuk pelestarian goresan tipis. Temuan ini sejalan dengan teori GAN: kehilangan \textit{adversarial} mengatur pertukaran antara realisme dan \textit{fidelity}, sementara kehilangan perseptual menjaga topologi goresan melalui fitur VGG.

Bobot \textit{adversarial} yang terlalu tinggi menyebabkan kolaps modus (\textit{generator} menghasilkan variasi terbatas) dan ketidakstabilan pelatihan. Sebaliknya, bobot yang terlalu rendah menghasilkan keluaran buram yang kurang realistis. Nilai optimal yang ditemukan (3.0) menyeimbangkan kedua kekhawatiran ini sambil mempertahankan stabilitas pelatihan.

Bobot kehilangan fitur pengenalan memiliki dampak sedang, menunjukkan bahwa panduan sadar-teks memberikan nilai tambah tetapi harus diimbangi dengan tujuan kualitas visual. Studi ablasi (Bagian V-B) mengonfirmasi bahwa konfigurasi 4 komponen tanpa kehilangan fitur pengenalan mencapai keseimbangan optimal antara kualitas visual dan keterbacaan HTR.

\subsection{Generalisasi pada Dokumen Historis Nyata}

Evaluasi kualitatif pada dokumen ANRI nyata (Bagian V-C) menunjukkan generalisasi yang wajar meskipun model hanya dilatih pada degradasi sintetis, dengan mayoritas dokumen menunjukkan perbaikan berguna. Validasi pakar paleografi mengindikasikan praktikalitas metode yang diusulkan untuk penerapan pada alur kerja arsip dunia nyata.

Namun, efek kesenjangan domain teramati pada kasus degradasi ekstrem (tingkat kegagalan 13.3\%), terutama untuk korosi tinta \textit{iron gall} berat dan tinta yang sangat memudar. Hal ini dapat dipahami karena jalur degradasi sintetis belum mencakup model korosi kimia dan tingkat keparahan ekstrem yang ditemukan pada dokumen berusia 250--450 tahun.

Strategi yang berpotensi meningkatkan generalisasi:
\begin{enumerate}
    \item Adaptasi domain melalui \textit{fine-tuning} pada \textit{set} kecil dokumen terdegradasi nyata dengan strategi \textit{few-shot learning}
    \item Peningkatan degradasi sintetis dengan menggabungkan pola penuaan yang lebih realistis dan model korosi kimia tinta \textit{iron gall}
    \item Pembelajaran tanpa pengawasan untuk memanfaatkan dokumen nyata tanpa label melalui teknik \textit{self-supervised learning}
\end{enumerate}

\subsection{Pertimbangan Penerapan Praktis}

Untuk menerapkan sistem dalam alur kerja digitalisasi arsip, beberapa pertimbangan praktis perlu diperhatikan:

Model yang diusulkan memiliki kompleksitas komputasi lebih tinggi dibanding metode binerisasi tradisional (Sauvola, Otsu) karena arsitektur GAN dengan diskriminator dual-modal dan pengenal beku. Waktu pelatihan tercatat 48 GPU-jam untuk 50 \textit{epoch} pada RTX A4000 16GB. Untuk penerapan produksi, inferensi dapat dilakukan dengan \textit{batch processing} untuk meningkatkan \textit{throughput}, yang cocok untuk proyek arsip skala besar dengan ribuan dokumen. Peningkatan dalam akurasi HTR (detail di Bagian~V) dapat membenarkan investasi komputasional untuk aplikasi digitisasi arsip bernilai tinggi. Meskipun kinerja rerata kuat, kasus kegagalan tetap ada, terutama pada ligatur paleografi kompleks dan teks yang sangat memudar, sehingga tinjauan manusia direkomendasikan untuk dokumen penting, atau dapat diterapkan mekanisme penilaian kepercayaan otomatis untuk menandai restorasi yang memerlukan verifikasi manual.

Restorasi dokumen mengubah penampilan asli. Untuk penggunaan ilmiah dan pelestarian arsip, versi asli dan yang direstorasi harus diarsipkan secara terpisah dengan \textit{metadata} provenans yang jelas, termasuk informasi model yang digunakan, tanggal restorasi, dan skor kepercayaan.

\subsection{Keterbatasan dan Arah Penelitian Masa Depan}

Keterbatasan Penelitian Saat Ini:
\begin{enumerate}
    \item \textbf{\textit{Trade-off} Evaluasi Kuantitatif vs. Realisme Penuh}: Evaluasi kuantitatif memerlukan \textit{ground truth} berpasangan yang tidak tersedia pada dokumen historis nyata. \textit{Dataset} semi-sintetis realistis yang dikonstruksi dari dokumen ANRI autentik menyediakan tulisan tangan paleografi nyata dan pola degradasi historis autentik, namun augmentasi terkontrol mungkin tidak sepenuhnya mereplikasi kompleksitas degradasi kimiawi berabad-abad (korosi tinta \textit{iron gall}, perubahan pH lokal). Hasil validasi kualitatif pada 15 naskah ANRI dengan degradasi alami menunjukkan generalisasi efektif, mengonfirmasi relevansi praktis meskipun evaluasi kuantitatif menggunakan \textit{dataset} dengan augmentasi.
    \item \textbf{Keterbatasan Generalisasi Lintas Domain}: Dilatih terutama pada aksara Latin (Belanda); generalisasi lintas bahasa dan aksara belum teruji secara komprehensif
    \item \textbf{Degradasi Ekstrem dan Kontras Rendah}: Model mengalami kesulitan pada kasus ekstrem (kontras $<$30, pemudaran sangat berat) yang kurang terwakili dalam distribusi \textit{dataset} pelatihan, sebagaimana divalidasi pada evaluasi dokumen ANRI nyata
    \item \textbf{Pola Kegagalan Spesifik}: (a) Ligatur paleografi kompleks (62.9\% dari kasus CER$>$15\%), (b) Teks bertumpang tindih dengan stempel arsip, (c) Artefak numerik dan simbol khusus (ƒ, §, =), (d) Korosi tinta \textit{iron gall} dengan efek kimia kompleks
\end{enumerate}

Arah Penelitian Masa Depan:

Peningkatan Arsitektur dan Protokol Pelatihan:
\begin{enumerate}
    \item Eksplorasi ruang fitur alternatif untuk panduan HTR: fitur semantik pasca-\textit{transformer}, pencocokan fitur multi-skala dengan bobot adaptif, kesamaan \textit{cosine} sebagai pengganti MSE, dan penyelarasan fitur terpandu atensi
    \item Optimisasi protokol pelatihan untuk efektivitas \textit{dual-modal}: supervisi teks \textit{ground truth}, penjadwalan bobot kehilangan dinamis, arsitektur \textit{generator} sadar-teks, dan strategi pelatihan progresif
    \item Peningkatan arsitektur \textit{generator} dengan model berbasis \textit{transformer} atau model difusi untuk kapasitas representasi yang lebih tinggi
\end{enumerate}

Peningkatan Generalisasi dan Evaluasi:
\begin{enumerate}
    \setcounter{enumi}{8}
    \item Kuantifikasi ketidakpastian dengan skor kepercayaan untuk kualitas restorasi
    \item Evaluasi lintas bahasa dan aksara (Arab, Melayu, Inggris, Indonesia dan aksara historis lainnya)
    \item Adaptasi \textit{few-shot} untuk memungkinkan penyesuaian cepat pada jenis dokumen baru dengan contoh minimal
    \item Pembuatan \textit{ground truth} dokumen historis nyata melalui kolaborasi dengan pakar paleografi
\end{enumerate}

Penerapan Praktis:
\begin{enumerate}
    \setcounter{enumi}{12}
    \item Restorasi interaktif dengan umpan balik pengguna untuk wilayah yang menantang
    \item Optimisasi bobot kehilangan sistematis menggunakan optimisasi Bayesian atau pencarian berbasis evolusi
    \item Pengembangan sistem kontrol kualitas otomatis untuk deteksi halusinasi dan artefak
\end{enumerate}

\subsection{Pertimbangan Etis dan Pelestarian Integritas Historis}

\subsubsection{Etika Modifikasi Dokumen Historis}

Restorasi otomatis dokumen historis menimbulkan pertanyaan etis penting yang harus dipertimbangkan dalam penerapan praktis. Dokumen yang direstorasi menggunakan metode yang diusulkan tidak dapat menggantikan dokumen asli dan harus diperlakukan sebagai alat bantu penelitian, bukan sebagai kebenaran sejarah absolut.

Meskipun implementasi penelitian saat ini menghasilkan citra yang direstorasi tanpa \textit{metadata} otomatis, sistem produksi yang direkomendasikan untuk penerapan arsip harus menyertakan \textit{metadata} lengkap yang mencakup: label yang menandakan penggunaan kecerdasan buatan, versi model yang digunakan, tanggal restorasi, estimasi kepercayaan kualitas restorasi, dan tautan ke dokumen asli untuk verifikasi. Institusi arsip direkomendasikan untuk menerapkan protokol pengarsipan ganda yang mencakup: (1) penyimpanan dokumen asli sebagai salinan \textit{master}, (2) penyimpanan versi yang direstorasi sebagai salinan turunan dengan label yang jelas, (3) penggunaan format tanpa kehilangan (TIFF atau PNG), (4) implementasi \textit{checksum} kriptografis untuk verifikasi integritas, dan (5) pemeliharaan catatan \textit{audit} dari semua langkah pemrosesan.

\subsubsection{Deteksi dan Mitigasi Halusinasi}

Model generatif secara umum berpotensi menambahkan informasi yang tidak ada dalam dokumen asli, terutama pada wilayah yang sangat terdegradasi. Metode yang diusulkan dirancang dengan strategi restorasi konservatif yang memprioritaskan negatif palsu (peluang peningkatan terlewat) daripada positif palsu (goresan halusinasi), sebagaimana divalidasi melalui analisis kualitatif dokumen ANRI. Meskipun desain ini meminimalkan risiko halusinasi, penelitian masa depan dapat mengembangkan mekanisme deteksi eksplisit untuk sistem produksi, seperti: (1) analisis densitas goresan dengan membandingkan jumlah piksel teks sebelum dan sesudah restorasi untuk mendeteksi penambahan goresan yang tidak wajar, (2) kuantifikasi ketidakpastian untuk menandai wilayah dengan kepercayaan rendah sebagai memerlukan verifikasi manual, dan (3) pemeriksaan konsistensi dengan membandingkan hasil dari beberapa inferensi untuk memastikan keluaran yang stabil.

Sistem deteksi otomatis semacam itu dapat dikembangkan untuk mengirim peringatan verifikasi manual jika terdapat modifikasi signifikan pada citra yang melebihi ambang batas yang ditentukan, memastikan bahwa potensi halusinasi dapat diidentifikasi sebelum dokumen digunakan untuk analisis historis. Pengembangan sistem kontrol kualitas otomatis ini merupakan arah penelitian masa depan.

\subsubsection{Bias dalam Data Pelatihan}

\textit{Dataset} sintetis yang digunakan untuk pelatihan GAN (diturunkan dari dokumen ANRI yang relatif bersih dengan jalur degradasi multitahap sintetis) tidak sepenuhnya mereplikasi degradasi alami dokumen era VOC (1650--1800) karena perbedaan karakteristik: (1) degradasi sintetis (\textit{bleed-through}, derau tekstur, \textit{blur}) tidak menangkap kompleksitas penuh degradasi historis seperti korosi tinta \textit{iron gall} dengan perubahan pH lokal, (2) pola penuaan sintetis tidak sepenuhnya mereplikasi degradasi alami selama berabad-abad, dan (3) variasi kondisi penyimpanan arsip (kelembaban, temperatur, paparan cahaya) yang mempengaruhi degradasi dokumen nyata.

Untuk mengurangi bias ini, penelitian masa depan dapat mengeksplorasi beberapa strategi: \textit{fine-tuning} pada dokumen ANRI nyata terdegradasi untuk adaptasi domain, pembelajaran \textit{few-shot} dengan contoh minimal dari domain target, dan integrasi umpan balik koreksi dengan melibatkan pakar paleografi dalam \textit{loop} untuk kasus yang menantang.

\subsubsection{Rekomendasi untuk Praktisi}

Untuk pengelola arsip, direkomendasikan untuk: (1) menggunakan citra yang direstorasi sebagai titik awal untuk analisis, bukan sebagai kebenaran akhir, (2) selalu verifikasi hasil yang krusial dengan membandingkan pada dokumen asli, (3) mendokumentasikan informasi model (versi, konfigurasi, tanggal restorasi) dalam catatan arsip, dan (4) menerapkan protokol pengarsipan ganda untuk menjaga integritas dokumen asli. Keterbatasan saat ini adalah sistem belum menyediakan skor kepercayaan otomatis, sehingga penilaian kualitas memerlukan inspeksi visual oleh pakar.

Untuk pengembang sistem yang akan mengimplementasikan metode ini dalam produksi, direkomendasikan untuk: (1) mengimplementasikan alat \textit{explainability} seperti peta atensi untuk transparansi proses restorasi, (2) mengembangkan estimasi ketidakpastian untuk setiap prediksi restorasi, (3) mempertahankan sistem \textit{versioning} model untuk reproduktibilitas, (4) melakukan \textit{audit} bias secara berkala untuk memastikan keadilan dan akurasi, dan (5) mengintegrasikan mekanisme deteksi halusinasi untuk kontrol kualitas.

% ============================================================ 
% VII. KESIMPULAN
% ============================================================ 
\section{Kesimpulan}

Penelitian ini mengusulkan kerangka kerja restorasi dokumen berorientasi HTR berbasis GAN yang mengatasi keterbatasan metode konvensional dalam mengoptimalkan keterbacaan teks untuk sistem pengenalan otomatis. Metode yang diusulkan mengintegrasikan dua kontribusi utama: (1) strategi \textit{frozen HTR recognizer} yang memberikan gradien sawar-teks stabil tanpa ketidakstabilan pelatihan bersama, terbukti menjadi faktor kunci keberhasilan sistem melalui studi ablasi ($p<0.001$), dan (2) optimasi fungsi loss multikomponen yang mengidentifikasi konfigurasi 4-komponen (piksel, \textit{adversarial}, perseptual, CTC) sebagai keseimbangan optimal antara kualitas visual dan keterbacaan HTR. Eksplorasi arsitektur diskriminator \textit{dual-modal} menunjukkan kontribusi marginal yang tidak signifikan secara statistik ($\Delta$PSNR +0.28~dB, $p>0.05$), mengonfirmasi bahwa protokol pelatihan dan strategi \textit{recognizer} beku lebih dominan daripada kompleksitas arsitektur diskriminator.

Evaluasi kuantitatif pada 712 citra uji dari \textit{dataset} semi-sintetis realistis berbasis dokumen paleografi ANRI abad ke-16 hingga ke-18 menunjukkan pengurangan CER dari 83.4\% (kondisi terdegradasi) menjadi 34.9\% (pengurangan relatif 58.2\%, $p<0.001$), mendekati kinerja teoretis pada citra bersih (CER 34.1\%), sambil mempertahankan kualitas visual tinggi (PSNR 30.74~dB, SSIM 0.987). Validasi kualitatif pada 15 naskah ANRI autentik mengonfirmasi generalisasi efektif pada degradasi historis nyata yang tidak dapat direplikasi secara sintetis. Studi ablasi sistematis memvalidasi bahwa kombinasi 4-komponen kehilangan mencapai optimalitas dalam ruang multi-objektif, dengan temuan penting bahwa kehilangan fitur \textit{recognizer} bersifat kontraproduktif dan harus dihindari.

Temuan utama penelitian ini adalah bahwa integrasi eksplisit objektif HTR melalui \textit{frozen recognizer} selama pelatihan GAN menghasilkan keluaran yang dioptimalkan untuk keterbacaan mesin sambil mempertahankan kualitas visual, dengan implikasi praktis untuk proyek digitalisasi arsip historis skala besar. Keterbatasan meliputi tantangan pada degradasi ekstrem (kontras $<$30, pemudaran sangat berat), ligatur paleografi kompleks (62.9\% dari kasus kegagalan), dan pembatasan pada aksara Latin paleografi. Penelitian ke depan akan fokus pada pelestarian ligatur paleografi, pemodelan degradasi kimia historis (korosi tinta \textit{iron gall}), pemisahan layer dokumen berlapis, peningkatan robustness kontras rendah, dan validasi \textit{transfer-learning} pada korpus multi-bahasa dan multi-aksara untuk mengonfirmasi generalisasi lintas domain.

% ============================================================ 
% REFERENCES
% ============================================================ 
\bibliographystyle{IEEEtran}
\begin{thebibliography}{99}

\bibitem{otsu1979threshold}
N. Otsu, ``A threshold selection method from gray-level histograms,'' \textit{IEEE Trans. Syst., Man, Cybern.}, vol. 9, no. 1, pp. 62--66, Jan. 1979.

\bibitem{sauvola2000adaptive}
J. Sauvola and M. Pietikäinen, ``Adaptive document image binarization,'' \textit{Pattern Recognit.}, vol. 33, no. 2, pp. 225--236, Feb. 2000.

\bibitem{niblack1986introduction}
W. Niblack, \textit{An Introduction to Digital Image Processing}. Englewood Cliffs, NJ, USA: Prentice-Hall, 1986.

\bibitem{wolf2004extraction}
C. Wolf, J.-M. Jolion, and F. Chassaing, ``Text localization, enhancement and binarization in multimedia documents,'' in \textit{Proc. 16th Int. Conf. Pattern Recognit.}, vol. 2, 2002, pp. 1037--1040.

\bibitem{tensmeyer2017document}
C. Tensmeyer and T. Martinez, ``Document image binarization with fully convolutional neural networks,'' in \textit{Proc. 14th IAPR Int. Conf. Document Anal. Recognit.}, vol. 1, 2017, pp. 99--104.

\bibitem{zhao2019document}
J. Zhao, C. Shi, F. Jia, Y. Wang, and B. Xiao, ``Document image binarization with cascaded generators of conditional generative adversarial networks,'' \textit{Pattern Recognit.}, vol. 96, p. 106968, Dec. 2019.

\bibitem{souibgui2020degan}
M. A. Souibgui, S. Khamekhem Jemni, Y. Kessentini, C. Fornés, A. Fornés, and C. Ogier, ``DE-GAN: A conditional generative adversarial network for document enhancement,'' \textit{IEEE Trans. Pattern Anal. Mach. Intell.}, vol. 44, no. 3, pp. 1180--1191, Mar. 2022.

\bibitem{souibgui2021enhance}
M. A. Souibgui, S. Khamekhem Jemni, Y. Kessentini, and A. Fornés, ``Enhance to read better: A multi-task adversarial network for handwritten document image enhancement,'' \textit{Pattern Recognit.}, vol. 123, p. 108368, Mar. 2022.

\bibitem{ronneberger2015unet}
O. Ronneberger, P. Fischer, and T. Brox, ``U-Net: Convolutional networks for biomedical image segmentation,'' in \textit{Proc. Int. Conf. Med. Image Comput. Comput.-Assist. Interv.}, 2015, pp. 234--241.

\bibitem{johnson2016perceptual}
J. Johnson, A. Alahi, and L. Fei-Fei, ``Perceptual losses for real-time style transfer and super-resolution,'' in \textit{Proc. Eur. Conf. Comput. Vis.}, 2016, pp. 694--711.

\bibitem{wang2004ssim}
Z. Wang, A. C. Bovik, H. R. Sheikh, and E. P. Simoncelli, ``Image quality assessment: From error visibility to structural similarity,'' \textit{IEEE Trans. Image Process.}, vol. 13, no. 4, pp. 600--612, Apr. 2004.

\bibitem{marti2002iam}
U.-V. Marti and H. Bunke, ``The IAM-database: An English sentence database for offline handwriting recognition,'' \textit{Int. J. Document Anal. Recognit.}, vol. 5, no. 1, pp. 39--46, Nov. 2002.

\bibitem{perlin1985image}
K. Perlin, ``An image synthesizer,'' \textit{ACM SIGGRAPH Comput. Graph.}, vol. 19, no. 3, pp. 287--296, Jul. 1985.

\bibitem{isola2017image}
P. Isola, J.-Y. Zhu, T. Zhou, and A. A. Efros, ``Image-to-image translation with conditional adversarial networks,'' in \textit{Proc. IEEE Conf. Comput. Vis. Pattern Recognit.}, 2017, pp. 1125--1134.

\bibitem{zhu2017unpaired}
J.-Y. Zhu, T. Park, P. Isola, and A. A. Efros, ``Unpaired image-to-image translation using cycle-consistent adversarial networks,'' in \textit{Proc. IEEE Int. Conf. Comput. Vis.}, 2017, pp. 2223--2232.

\bibitem{souibgui2022docent}
M. A. Souibgui, Y. Kessentini, and A. Fornés, ``DocEnTr: An end-to-end document image enhancement transformer,'' \textit{Pattern Recognit.}, vol. 142, p. 109710, Oct. 2023.

\bibitem{shi2016crnn}
B. Shi, X. Bai, and C. Yao, ``An end-to-end trainable neural network for image-based sequence recognition and its application to scene text recognition,'' \textit{IEEE Trans. Pattern Anal. Mach. Intell.}, vol. 39, no. 11, pp. 2298--2304, Nov. 2017.

\bibitem{bluche2017scan}
T. Bluche, J. Louradour, and R. Messina, ``Scan, attend and read: End-to-end handwritten paragraph recognition with MDLSTM attention,'' in \textit{Proc. 14th IAPR Int. Conf. Document Anal. Recognit.}, vol. 1, 2017, pp. 1050--1055.

\bibitem{michael2019evaluating}
J. Michael, R. Labahn, T. Grüning, and J. Zöllner, ``Evaluating sequence-to-sequence models for handwritten text recognition,'' in \textit{Proc. 15th Int. Conf. Document Anal. Recognit.}, 2019, pp. 1286--1293.

\bibitem{sanchez2016icfhr}
J. A. Sánchez, V. Romero, A. H. Toselli, and E. Vidal, ``ICFHR2016 competition on handwritten text recognition on the READ dataset,'' in \textit{Proc. 15th Int. Conf. Frontiers Handwriting Recognit.}, 2016, pp. 630--635.

\bibitem{sanchez2017icdar}
J. A. Sánchez, V. Romero, A. H. Toselli, and E. Vidal, ``ICDAR2017 competition on handwritten text recognition on the READ dataset,'' in \textit{Proc. 14th IAPR Int. Conf. Document Anal. Recognit.}, vol. 1, 2017, pp. 1383--1388.

\bibitem{diaz2021rethinking}
D. Díaz, A. H. Toselli, E. Vidal, and A. Fornés, ``Rethinking optical character recognition with transformers,'' \textit{Pattern Recognit.}, vol. 122, p. 108301, Feb. 2022.

\bibitem{goodfellow2016improved}
I. Goodfellow, J. Pouget-Abadie, M. Mirza, B. Xu, D. Warde-Farley, S. Ozair, A. Courville, and Y. Bengio, ``Generative adversarial nets,'' in \textit{Proc. Adv. Neural Inf. Process. Syst.}, 2014, pp. 2672--2680.

\bibitem{zhang2018residual}
Y. Zhang, Y. Tian, Y. Kong, B. Zhong, and Y. Fu, ``Residual dense network for image super-resolution,'' in \textit{Proc. IEEE Conf. Comput. Vis. Pattern Recognit.}, 2018, pp. 2472--2481.

\bibitem{woo2018cbam}
S. Woo, J. Park, J.-Y. Lee, and I. S. Kweon, ``CBAM: Convolutional block attention module,'' in \textit{Proc. Eur. Conf. Comput. Vis.}, 2018, pp. 3--19.

\bibitem{lin2017fpn}
T.-Y. Lin, P. Dollár, R. Girshick, K. He, B. Hariharan, and S. Belongie, ``Feature pyramid networks for object detection,'' in \textit{Proc. IEEE Conf. Comput. Vis. Pattern Recognit.}, 2017, pp. 2117--2125.

\bibitem{chen2017deeplab}
L.-C. Chen, G. Papandreou, I. Kokkinos, K. Murphy, and A. L. Yuille, ``DeepLab: Semantic image segmentation with deep convolutional nets, atrous convolution, and fully connected CRFs,'' \textit{IEEE Trans. Pattern Anal. Mach. Intell.}, vol. 40, no. 4, pp. 834--848, Apr. 2018.

\bibitem{oktay2018attention}
O. Oktay et al., ``Attention U-Net: Learning where to look for the pancreas,'' in \textit{Proc. Med. Imag. Deep Learn.}, 2018, pp. 1--10.

\bibitem{kirkpatrick2017overcoming}
J. Kirkpatrick et al., ``Overcoming catastrophic forgetting in neural networks,'' \textit{Proc. Natl. Acad. Sci. USA}, vol. 114, no. 13, pp. 3521--3526, Mar. 2017.

\bibitem{souibgui2022}
M. A. Souibgui et al., ``Enhancing handwritten text recognition with N-gram sequence alignment: A hybrid loss approach for improved accuracy,'' \textit{Pattern Recognit.}, vol. 125, p. 108515, May 2022.

\bibitem{kang2021pay}
L. Kang, P. Riba, M. Rusiñol, A. Fornés, and M. Villegas, ``Pay attention to what you read: Non-recurrent handwritten text-line recognition,'' \textit{Pattern Recognit.}, vol. 129, p. 108766, Sep. 2022.

\bibitem{he2016deep}
K. He, X. Zhang, S. Ren, and J. Sun, ``Deep residual learning for image recognition,'' in \textit{Proc. IEEE Conf. Comput. Vis. Pattern Recognit.}, 2016, pp. 770--778.

\bibitem{hochreiter1997lstm}
S. Hochreiter and J. Schmidhuber, ``Long short-term memory,'' \textit{Neural Comput.}, vol. 9, no. 8, pp. 1735--1780, Nov. 1997.

\bibitem{vaswani2017attention}
A. Vaswani et al., ``Attention is all you need,'' in \textit{Proc. Adv. Neural Inf. Process. Syst.}, 2017, pp. 5998--6008.

\bibitem{pratikakis2016hdibco}
I. Pratikakis, K. Zagoris, G. Barlas, and B. Gatos, ``ICFHR2016 handwritten document image binarization contest (H-DIBCO 2016),'' in \textit{Proc. 15th Int. Conf. Frontiers Handwriting Recognit.}, 2016, pp. 619--623.

\end{thebibliography}

% ============================================================ 
% PENGHARGAAN
% ============================================================ 
\section*{Penghargaan}

Para penulis ingin mengucapkan terima kasih kepada Arsip Nasional Republik Indonesia (ANRI) yang telah memberikan akses ke koleksi dokumen historis dan  penyediaan infrastruktur komputasi  yang sangat berharga untuk penelitian ini. Apresiasi khusus ditujukan kepada Bandung Institute of Technology (ITB), khususnya Program Studi \textit{Smart Systems Program} di School of Electrical Engineering and Informatics atas lingkungan riset yang mendukung pengembangan penelitian ini. Terima kasih juga kami sampaikan kepada Kementerian Komunikasi dan Informatika Republik Indonesia (Komdigi) atas dukungan beasiswa yang memungkinkan terlaksananya penelitian ini.

% ============================================================ 
% BIOGRAFI
% ============================================================ 

% \vfill\pagebreak

\noindent
\begin{minipage}[t]{0.48\textwidth}
\begin{IEEEbiography}[{\includegraphics[width=1in,height=1.25in,clip,keepaspectratio]{../penulis/fotoJatniko.png}}]{Jatniko Nur Mutaqin}
saat ini sedang menempuh pendidikan magister di Bandung Institute of Technology (ITB) pada Program Studi \textit{Smart Systems Program}, School of Electrical Engineering and Informatics. Beliau bekerja sebagai \textit{AI Engineer} di Arsip Nasional Republik Indonesia (ANRI), di mana pengalaman praktisnya dalam pelestarian dan digitalisasi dokumen historis menginspirasi riset akademisnya. Minat penelitiannya meliputi analisis citra dokumen, \textit{generative adversarial networks} (GAN), sistem \textit{optical character recognition} dan \textit{handwritten text recognition}, serta aplikasi kecerdasan buatan untuk pelestarian warisan budaya dan dokumen historis.
\end{IEEEbiography}
\end{minipage}%
\hfill
\begin{minipage}[t]{0.48\textwidth}
\begin{IEEEbiography}[{\includegraphics[width=1in,height=1.25in,clip,keepaspectratio]{../penulis/fotoBhaskara.jpg}}]{I Gusti Bagus Baskara Nugraha}
(Ir., S.T., M.T., Ph.D.) adalah akademisi dan peneliti yang berafiliasi dengan Bandung Institute of Technology (ITB), khususnya di School of Electrical Engineering and Informatics. Keahliannya terletak pada Sistem Informasi Manajemen, dengan topik penelitian meliputi platform kota pintar (\textit{smart city}), sistem berbasis \textit{blockchain}, pembelajaran mesin (\textit{machine learning}), dan sistem manajemen energi. Beliau telah berkontribusi pada berbagai makalah akademis dan konferensi, dengan penelitian terkemuka dalam \textit{blockchain}, teknologi kota pintar, dan aplikasi kecerdasan buatan.
\end{IEEEbiography}
\end{minipage}

\end{document}